% ============================================================
%  Harald Høffding, "Den store Humor. En psykologisk Studie"
%  "Humor as a Feeling of Life (The Great Humor). A Psychological Study"
%  (title after the German edition, "Humor als Lebensgefühl (Der große Humor).
%  Eine psychologische Studie", Leipzig 1918, so that the English title does
%  not read as a book about comedy; "Livsfølelse", feeling of life, is
%  Høffding's own word, and § 13 sets the great humor apart from a mere mood)
%  København og Kristiania: Gyldendalske Boghandel, Nordisk Forlag,
%  1916.  178 pp.
%
%  TRANSLATION (English), from transcription.tex (Danish) in this directory,
%  the text of record (collated 2026-10-04; see its header).
%
%  DRAFT COMPLETE (2026-10-04): the whole book, pp. 7--178, translated in
%  reading order from the collated transcription; \opage 7--178 (no 176) in
%  order, once each; 108 notes, as in the Danish.  To be read by an
%  independent reviewer.
%
%  BOOK-SPECIFIC CONVENTIONS (beyond TRANSLATION-PLAYBOOK.md §2):
%   * Høffding's terms: Totalfølelse = total feeling, Enkeltfølelse = single
%     feeling; Totaltilstand / Enkelttilstand = total state / single state;
%     Sammensmeltning = fusion; Organisation = organization; Grundværdi =
%     fundamental value; den store Humor / den lille Humor = the great humor /
%     the little humor; Spøg / Alvor = jest / earnest; Vemod = wistfulness;
%     Sindshøjhed = loftiness of mind; Overlegenhed = superiority; Livsanskuelse =
%     view of life; realt Jeg = real self (Høffding's term, Psykologi V B, 5).
%   * Italics in the Danish (\textit: titles, foreign words) -> \textit;
%     letterspacing (\emph) -> \emph.
%   * The sections are numbered 1.--54. continuously, run in at the start of a
%     paragraph, as printed.  The print numbers two sections "49." (pp. 160,
%     161); the second is given as 50 with a comment.
%   * Quotations in German, French, Latin, English, Swedish and Greek stay in
%     the original.  Quotations Høffding gives in Danish translation are
%     translated from his Danish (Darwin's letter among them), except where the
%     English original is certain (Swift, Carlyle; marked at the site).  Danish
%     verse is translated in the text, the Danish kept in a comment.  Titles
%     cited in the notes stay as cited (Høffding's own works in Danish);
%     names and citation numbers as printed (Wolf, Fontaine, Etienne, Rhode,
%     Yrjo, King Lear V, 3, 203).
%   * Footnotes are numbered through each chapter (the print numbers them per
%     page, 1), 2), ...).  "Smlgn." / "Smlg." -> "Cf."
%   * Misprints: translate the sense, with a comment at the site (the Danish
%     marks them PRINTED AS IS).
%   * Page markers \opage{N} placed by hand, inline, at the English word where
%     printed page N begins.
%   * The Register (pp. 177--178) is translated as the Index: the print's order,
%     its page numbers (the original pagination, in the margin of this
%     translation), the Danish headword in italics after the English.
% ============================================================
\documentclass[12pt]{book}
\usepackage[utf8]{inputenc}
\usepackage[T1]{fontenc}
\usepackage[english]{babel}
\usepackage[protrusion=true,expansion=true]{microtype}
\usepackage[margin=1.3in]{geometry}
\usepackage{setspace}
\usepackage{amsmath}
\usepackage{enumitem}
\usepackage{libertinus}
\usepackage{libertinust1math}
\usepackage{textalpha}
\usepackage{fancyhdr}
\usepackage[colorlinks=true,linkcolor=black,urlcolor=black,
  pdftitle={Humor as a Feeling of Life},
  pdfauthor={Harald Høffding}]{hyperref}
\setlength{\headheight}{15pt}
\pagestyle{fancy}
\fancyhf{}
\fancyhead[LE,RO]{\thepage}
\fancyhead[RE]{\textit{Humor as a Feeling of Life}}
\fancyhead[LO]{\textit{\leftmark}}
\setlength{\parindent}{1.5em}
\setstretch{1.3}
\newcommand{\tocline}[3]{\noindent\makebox[2.2em][r]{#1}\hspace{0.6em}#2\dotfill\ #3\par}
\newcommand{\entry}{\par\noindent\hangindent=1.5em }
\usepackage{marginnote}
\newcommand{\opage}[1]{\marginnote{\footnotesize\textit{[#1]}}}

\title{Humor as a Feeling of Life\\[0.4em]
  \large (The Great Humor)\\[0.2em]
  A Psychological Study}
\author{Harald Høffding}
\date{Translated from the Danish\\[0.6em]
  \small \textit{Den store Humor. En psykologisk Studie.} Copenhagen and
  Kristiania: Gyldendalske Boghandel, Nordisk Forlag, 1916}

\begin{document}
\frontmatter
\maketitle

% ===== The author's other works (p. [2] of the print), titles as printed =====
\chapter*{By the Same Author}
{\small\setlength{\parindent}{0pt}\setlength{\parskip}{0.2em}
\everypar{\hangindent=2em}
% The list is given as printed (the titles are Danish; Høffding's own titles stay in Danish).
\textit{Den antike Opfattelse af Menneskets Villie}. 1870.\par
\textit{Philosophien i Tyskland efter Hegel}. 1872.\par
\textit{Den engelske Philosophi i vor Tid}. 1874.\par
\textit{Om Grundlaget for den humane Ethik}. 1876.\par
\textit{Psykologi i Omrids paa Grundlag af Erfaring}. (1882). Sixth, newly revised edition. 1911.\par
\textit{Formel Logik. Til Brug ved Forelæsninger}. (1884). Sixth edition. 1913.\par
\textit{Etik. En Fremstilling af de etiske Principer og deres Anvendelse paa de vigtigste Livsforhold}. (1887). Fourth, newly revised and partly altered edition. 1913.\par
\textit{Psykologiske Undersøgelser}. 1889. (Videnskabernes Selskabs Skrifter).\par
\textit{Etiske Undersøgelser. Om Muligheden af en filosofisk Etik. --- Om Velfærdsprincipet. --- Forholdsloven i Etiken}. 1891.\par
\textit{Søren Kierkegaard som Filosof}. 1892.\par
\textit{Kontinuiteten i Kants filosofiske Udviklingsgang}. 1893. (Videnskabernes Selskabs Skrifter).\par
\textit{Den nyere Filosofis Historie. En Fremstilling af Filosofiens Historie fra Renaissancens Slutning til vore Dage}. Two volumes. (1894--95). Second, revised edition. 1904.\par
\textit{Jean Jacques Rousseau og hans Filosofi}. (1896). Second edition. 1912.\par
\textit{Kort Oversigt over den nyere Filosofis Historie}. (1898). Fifth, newly revised edition. 1913.\par
\textit{Mindre Arbejder}. First series. 1899.\par
\textit{Det psykologiske Grundlag for logiske Domme}. 1899. (Videnskabernes Selskabs Skrifter).\par
\textit{Ved Aarhundredeskiftet}. Address delivered at the University of Copenhagen. 1901.\par
\textit{Religionsfilosofi} (1901). Second printing. 1906.\par
\textit{Filosofiske Problemer}. 1902. (University program).\par
\textit{Om nogle religionsfilosofiske Arbejder fra den nyeste Tid}. 1903. (University program).\par
\textit{Moderne Filosofer}. 1904.\par
\textit{Mindre Arbejder}. Second series. 1905.\par
\textit{Danske Filosofer}. 1909.\par
\textit{Den menneskelige Tanke, dens Former og dens Opgaver}. 1910.\par
\textit{Religion og Videnskab}. 1910.\par
\textit{Udvalgte Stykker af dansk filosofisk Litteratur}. With introductions by Harald Høffding. 1910.\par
\textit{Personlighetsprincipen i filosofin}. Stockholm. 1911.\par
\textit{Mindre Arbejder}. Third series. 1913.\par
\textit{Henri Bergson's Filosofi. Karakteristik og Kritik}. 1914.\par}

% ===== CONTENTS (p. [5]) =====
\chapter*{Contents}
\hfill{\scriptsize Page}\par
\tocline{I.}{Humor as a Total Feeling}{7}
\tocline{II.}{Laughter and Humor}{48}
\tocline{III.}{Irony and Humor}{59}
\tocline{IV.}{Main Forms of Humor}{75}
\tocline{V.}{The Tragic and Humor}{88}
\tocline{VI.}{Understanding and Humor}{112}
\tocline{VII.}{Action and Humor}{124}
\tocline{VIII.}{Historical Conditions for Humor}{135}
\tocline{IX.}{Humor and Philosophy}{145}
\medskip
\noindent{\small [Page numbers are those of the Danish edition, given in the margin of this translation.]}

\mainmatter

% ===== I. HUMOR AS A TOTAL FEELING (pp. 7--47) =====
\chapter*{I.\quad Humor as a Total Feeling}
\addcontentsline{toc}{chapter}{I. Humor as a Total Feeling}
\markboth{Humor as a Total Feeling}{Humor as a Total Feeling}
\opage{7}1.~No mental state is altogether simple. As far as observation
can reach, leaving aside everything due to the influence of reflection, it
always finds a greater or lesser multiplicity of need or drive, sensation,
pleasure and pain, representation --- in short, of what in a word can be
called psychical elements; and it also finds a looser or firmer connection
among these elements. Every mental state is a totality within which a certain
continuity asserts itself, so that a sharp distinction among the different
elements always remains more or less artificial.

Continuity asserts itself also among the different mental states in their
relation to one another. The states often pass imperceptibly into one another,
and where the transition happens suddenly, with a jolt or a blow, it is just
this suddenness that makes the following state have its character determined
by its contrast to the preceding one, so that it cannot be understood without
it. There is always something artificial in keeping the different mental states
sharply apart.

When, nonetheless, by way of introduction to the inquiry to be undertaken
here, I distinguish between single states and total states, I find the
difference between them in this: the former are determined by a single
experience, while the latter are determined by a series of experiences. However
full a content a single state may have, this content --- sensations, thoughts,
drives, feelings --- is still conditioned by its common relation to a single
experience. In the coming-to-be of total states, on the other hand, several,
perhaps many experiences have had a part. There are of course many
\opage{8}transitional forms between a single experience and a series of
experiences. Altogether single is only the experience of a multiplicity given
simultaneously, for example the sight of a group of people or things at rest.
When it is a successive multiplicity that is experienced, as when one hears a
melody or follows the outline of a figure, or when a historical event unfolds
over a longer time, the experience consists of a series of single experiences.
It then depends on whether the single episodes, in their relation to the one
who experiences them, appear in an inner connection; only then will a real
total state be able to arise. Were there absolutely no connection among the
experiences, there would in the end, so far as we can see, be no mental states
whatsoever.

The arising of single states already presents difficulties enough for
psychological research; but total states contain by the nature of the case
still greater problems. Simple single states can perhaps be called forth
experimentally, and in any case one need not extend one's consideration beyond
certain relatively simple and surveyable relations. One can use analogies from
physiology, indeed perhaps from mechanics. Total states, on the other hand, are
harder to analyze. They demand a greater knowledge of the whole developmental
history of the mental life of the person concerned, perhaps also of the time
and the people to which he belongs, indeed even of the spiritual development of
the past. A historical (comparative, sociological) method must be used beside
the observing and analyzing method. No wonder that psychologists have by
preference kept to single states, and have often thought that total states could
without more ado be understood from them, just as mechanical wholes can be
understood when one only knows the parts.

Before we go further, let us give some examples of single states.

In the midst of the mood that a fine spring morning has called forth, I
suddenly receive sad news. Everything now gathers for me around the message;
its content lays claim to all reflection; its intrusion into relations that are
dear to me is drawn before me in broad strokes; the bright mood gives way to
dejection \opage{9}and dull sinking; everything that otherwise gave life worth
vanishes from consciousness. Attempts to think away what has happened, or to
find means of averting its consequences, are checked by the insight into the
inevitable. The organic feeling of life, my whole state, takes on the stamp of
inhibition, of depression. No thought can come up if it does not concern what
has happened. The day's work is done under compulsion, without inner assent.
After a single moment's forgetfulness the dark reality stands before me again
in its whole extent. And there may even stir a striving to immerse oneself
fully in the dark reality and to keep all light out; indeed, one can feel pangs
of conscience at having been able to think of other things at all.

Another example. A new idea dawns on a thinker. He sees the possibility of an
entirely new direction in research. With the force of intuition, which is all
the greater the more a period of doubt, discouragement and unsuccessful attempts
has gone before, views now open onto new regions, and scattered paths meet
before his eye in a common goal. A sudden concentration has taken place; a
synthetic intuition\footnote{See my paper on \textit{Begrebet Intuition} (Vid.\
Selsk.\ Oversigt 1914) p.~81. (Also \textit{Henri Bergson's Filosofi} p.~24 f.).}
has formed. Not only the life of thought, the whole mental life is set in strong
motion. No interest other than the intellectual can assert itself, no other
striving; and yet the highest enthusiasm is present. An experience of this kind
was had by René Descartes on the 10th of November 1619, while he lay in winter
quarters in Bavaria, when after long seeking it dawned on him which way he was
to follow in his research, the main features of his future philosophical and
mathematical work being revealed to him.\footnote{\textit{Den nyere Filosofis
Historie.}\textsuperscript{2} I p.~205.} When the enthusiasm, the sudden single
state, had spread itself out, there came the work of long years; indeed, only
eighteen years later did his fundamental work see the light. --- The history of
art offers similar examples. Well known, thus, is Mozart's description of how
motifs he had often hummed to himself involuntarily gathered into a whole of
tones that stood before him all at once (gleich Alles zusammen), so that he
\opage{10}could survey it as if it were a painting or a piece of sculpture.
Rousseau's conception of the problem of culture shows similar traits.

A third example. After long deliberation I have taken a decision in an
important matter. My whole being, that is, all thought, feeling and striving,
gathers on one point, in one wish, which is nevertheless more than a wish, since
I already see myself in the act of carrying out the action, and since its
effects, which I acknowledge as mine, are drawn before me in more or less clear
outline. I seem to notice that something from within me, perhaps from my
innermost, streams out into the imagined future in which the action is to take
place. My whole self, such as it now is, finds vent in the decision, comes
through it to breathe out.\footnote{Cf.\ on the psychology of decision my
\textit{Psykologi i Omrids paa Grundlag af Erfaring} VII B, 2 (sixth edition
p.~409--411).}

A last example I take from Ludvig Feilberg's rich collection of single states.
Feilberg's strength and weakness is that he is all too taken up with single
states; but he describes them with mastery. Total states he is inclined to look
down on as less genuine. Nor do they let themselves be so easily caught in
flight, and they cannot be described as vividly as, for example, the following
little experience. ``A small, freshly fallen green chestnut lay on the paving
under Holberg's tree beside some fallen twigs, older fallen chestnuts and
half-withered chestnut blossoms, which the wind had swept together. It was
surprising to see how beautiful it was in these surroundings; it had not been so
at all while it hung on the tree. On the contrary, the little prickly thing, which
was a reminder that the time of blossoms was over, had then rather aroused
displeasure. Here it was quite peculiarly cozy. One could not help noticing
it.''\footnote{\textit{Ludvig Feilberg: Samlede Skrifter.} I p.~172.}

2.~These examples of single states have deliberately not been chosen among the
simplest. They have the distinguishing mark that in each of them a manifold
content gathers around a relatively simple outer or inner experience. Even had
we chosen quite simple single states, it would still have turned out that
\opage{11}conditions owed to the person's earlier experiences are always
presupposed. Thousands have seen green chestnuts lying in Fiolstræde without
experiencing what Feilberg experienced, and the prize question that made the
problem of culture dawn on Rousseau has been seen by many without working as it
worked on him. Single states too can thus have a historical character. Whether
individuals in experiments on physiological time react sensorially or motorially
will, at least at the beginning, depend on their practice and habits. In small
things as in great it turns out that our present cannot be understood without
our past. But this holds in a quite special degree of total states, which are
determined by a long series of different experiences. Each of these experiences
has called forth its single state, and it is through these single states that a
total state can come to be. The presupposition is that they do not vanish
without a trace, but leave behind dispositions that can take part in the arising
of later states. All recognition and all recollection, as well as all
reflection, rests on this. Not least does this accumulation work in the domain
of the life of feeling. Through the alternation of joy and sorrow, hope and
fear, enthusiasm and dejection, new, lasting fundamental moods can work their way
forward. In conflict with inherited dispositions there is shaped through the
series of experiences what, in a more involuntary, nature-bound form, we call
temperament, and, by more or less conscious striving and acting, call character.
In my Psychology (V B, 5) I introduced the expression \emph{real self} to denote
what is thus partly the presupposition, partly the effect, of the influence of
experiences on mental states.

To forestall a misunderstanding that is all too frequent even among excellent
psychologists, it must be expressly stressed that this real self need not be an
object of consciousness for the individual himself. It need not give rise to any
``representation of the self.'' That the word ``I'' is used relatively early
means nothing; it is at first a mere name, a name of which the child later
discovers to its great surprise that it is a name all people claim with the same
right to have. There can be a ruling interest and striving --- and with it a real
self --- without any representation of it being formed, --- without its becoming
an object. The real self can undergo even very far-reaching changes
\opage{12}without the individual becoming conscious of it. The changes show
themselves immediately and involuntarily in his conduct and perhaps thereby
become apparent to others, although the individual himself does not discover
them.

Ludvig Feilberg even went so far as to want to abolish altogether words like
``self'' or ``personality,'' except insofar as they could be used in a
reproachful sense, namely of conceited self-consciousness, refined
self-absorption, that circles in a demanding way about its own little world
instead of forgetting itself in ``Ligeløb.'' % Feilberg's term (self-forgetful, even-paced absorption in the task); kept in Danish
He was so taken up with the
flowers he found in his psychological botanizing that he did not see how, in and
through the series of single states, a mental order of nature can show itself,
which forms precisely the basis of what psychology means by words like I, self or
personality. --- In his later writings, however, he had become clear about the
misunderstandings to which his polemic against ``personality'' could give rise.

It is true enough that express consciousness of what stirs in us can be
hazardous insofar as it leads beyond the immediate and involuntary, which is the
fruitful natural ground of spiritual life. But just as instinct can lead the
animal into the trap, so the involuntary unfolding of life can very easily turn
to harm, and it can be necessary for self-preservation that we bethink ourselves
which way can lead us to our goal. It will then often be necessary to push the
involuntary urge aside. It is especially at such turning points that we become
conscious of our real self, as we turn our attention to our inner being, to our
central striving, and it can then mean a great advance that consciousness makes
itself its object.

With regard to the real self, total states are more decisive than single
states. These latter can vanish without far-reaching effects, without standing,
either in their arising or through their influence, in close connection with what
in us is the properly central. Psychology does assume that all states and
experiences deposit effects and contribute to the development of the inner order
of nature, and so directly or indirectly, positively or \opage{13}negatively,
work to produce total states.\footnote{The so-called psychoanalytic school in
psychiatry (Freud and his disciples) has lately laid great weight on the fact
that experiences of far-reaching significance continually go on exerting their
influence on the mind and can thereby become the cause of mental illnesses. Under
hypnosis or in dreams the patient may come to betray such experiences, while
otherwise he perhaps involuntarily or voluntarily seeks to repress them. In
particular a certain group of mental illnesses (the ``hysterical'') are explained
as effects of mental wounds (as traumatic neuroses), and it becomes the doctor's
task to discover these wounds and to heal them by drawing them forth into clear
consciousness and making them an object of understanding.} But with this
assumption what is set up is properly only a great ideal for psychological
research, an ideal that must constantly be kept in view, but whose realization
will always be incomplete. Moreover, it will not always be a single total state
that comes to rule. Several total states can follow one another in the same
person's life, perhaps even to the degree that the very concept of a real self
can become problematic in the particular case. In very marked personalities,
however, there will be a total state that gives everything else in the mental
life its peculiar character. It need not be present uninterruptedly; but it
constantly exerts its after-effects and in any case plays an indirect part. And
in the moments decisive for the personality it will be the one that has the
word. It is to it that the personality swings back when collection,
concentration is needed. It expresses the inner man, whether or not this is
accessible to others. In some natures, which have been called the ``once-born''
(``the singly born''), the whole of life will be borne by the same total state,
even if moods and ripples arise that are not conditioned by it through and
through, and even if, when action is called for under particular circumstances,
emotions can arise that seem to have nothing to do with it. The ends a person
sets himself are on the whole determined by the total state ruling in him, the
prevailing disposition. But when work and struggle are needed for these ends,
thought and emotion can be released, owing to obstacles that present themselves
or means that must be procured, \opage{14}which in and of themselves do not
belong to the total state. Love of mankind can thus lead to fighting resistance
with courage and indignation; love of country can release warlike moods;
intellectual or artistic enthusiasm can lead to bitterness against spiritual
dullness and brutality. The famous criminologist Anselm v.~Feuerbach laid
particular weight on this apparent contrast between disposition and emotion as
necessary for understanding criminal characters. (See Psychology VI E, 5).

In the ``twice-born'' (``the doubly born'') the ruling total state is acquired
through a break or a crisis by which tendencies present beforehand are pushed
back. Continuity is not so plainly evident here as in the once-born, although a
deeper investigation can often show that the new and the old disposition have a
common stamp despite the wholly different direction in which the development of
the mental life takes place within them. While the development of the once-born
consists in an unfolding, the mental life of the twice-born is conditioned by the
overcoming of inner disharmonies.\footnote{The two designations go back to
Francis Newman (the Cardinal's humanist brother). The two types are often
described, under various names, in recent psychology of religion. Cf.\ my
\textit{Religionsfilosofi}\textsuperscript{2} p.~258--261 and \textit{Mindre
Arbejder.} Second series. p.~124 f.; 134 f.}

Examples of total states have already been indicated in what precedes, and I
shall add only a few more. Reverence, wistfulness, resignation, faithfulness,
the joy of knowledge are total states when they presuppose a series of
experiences of life, experiences of a more or less composite and intensive kind.
They are then different from a momentary flaring up and from instinctive
striving, from decisions called forth by passing situations, as well as from
moods and impulses that vanish as quickly as they come. But in their conflict
with practical circumstances they can take such more elementary psychical
functions into their service. In general it is ruling total states that underlie
every attempt to characterize personalities, every statement of the qualities
regarded as peculiar to a person.

Already in Plato there is an apt description of how \opage{15}the total state
comes to be. In the fourth book of the ``Republic'' (p.~429 D E) he compares the
different experiences with the different layers of color that must be applied to
a material before the particular color one wishes the material to have can be
applied. ``You know, surely, that when the dyers want to dye wool purple, they
first select from the many colors one kind, namely the white, and then work it
over with manifold preparations, so that it may take the purple in the best way,
and only then do they dye. What is dyed in this way is fast, and no washing can
rob it of its color!'' The quality or total state whose condition Plato wants to
state is courage, not in the sense of involuntary boldness, but in the sense of
the energy to hold fast under all circumstances to the conviction of what is
rightly to be feared and what not. It is the total state of firmness of character
or of faithfulness that he wants to describe. He wants to point out the way
education must go to be able to lead to such a state. Real life naturally does
not go about it so systematically as he wants to proceed in preparing the
qualities of character he wants to cultivate. In return, life's course can lead
to a richer multiplicity of qualities of character in different people. Fate
``grounds'' us with a series of strokes of color before the total state comes to
be that conditions the peculiar tone of color of our character.

3.~By different paths experiences can come into contact with one another, so
that they can enter into connections, and under the influence of these new total
states can be formed within the mental life.

It is customary to distinguish three kinds of elements, sides or qualities that
reflection can distinguish in every mental state, namely cognition (sensation and
representation), feeling of pleasure and pain, striving or willing. If we use
this threefold division, which there is no reason to abandon,\footnote{See my
paper on \textit{Begrebet Villie} (Mindre Arbejder. Third series).} there will
be three groups of laws or tendencies that can be assumed to play a part in the
arising of total states. Whether they also \opage{16}assert themselves in the
arising of the very simplest single states is a very difficult question. For we
may not assume without more ado that the divisions or distinctions that
reflection can and must make when it makes the mental life an object of inquiry
also hold for the primitive forms of the mental life. Psychological analysis
offers no complete analogy to chemical analysis. A compound substance can be
dissolved into parts that can exist each by itself, independently of the
compound. But what we can distinguish in the psychical states actually before us
are not parts that could have existence apart from the whole in which we think
we discover them. There exists no pure knowing, feeling or willing, always only
complexes of these psychical elements. And these complexes prove to be the more
intimately joined the further back we go in the history of the mental life. So
far as we know this history, for the individual person and for the race, it is a
main feature that the later stages present marked differences that do not appear
at earlier stages, and that we have no right to presuppose there. That threefold
division therefore holds only for the later stages of the mental life, in the
adult as opposed to the child, in man as opposed to the animal, in civilized
people as opposed to barbarians and savages. A schematic and deductive treatment
of psychology, which would combine the states that appear with the character of a
whole out of the elements that can be distinguished within the more highly
developed mental life, therefore leads to illusory results.\footnote{See more
on this my \textit{Psykologi} ch.~4 and \textit{Den menneskelige Tanke} (1909)
p.~9--19. --- On the basis of comparative psychology \textit{Edward Thorndike}
has energetically urged this truth. ``The old conception of human conscious
life,'' he says \textit{(Animal Intelligence.} New York 1911. p.~154), % the print has „In-intelligence“ (PRINTED AS IS in the Danish)
``is that it is built up of elementary sensations, --- that very small fragments
of consciousness (bits of consciousness) arise first, and that these are then
gradually built up into complex tissues.'' In contrast to this he maintains that
``development does not proceed from the small and simple to the large and
complex, but from immediate connection to mediate connection, within which a
host of isolated elements appear, --- from pure experience or undifferentiated
states to the arising of differences on the one hand, generalizations and
abstractions on the other.''}

\opage{17}As regards the elements of cognition, the association of
representations is the most important phenomenon. Most psychologists assume that
at the ground of all connection of representations there lies a need or
tendency to fill out representations that appear relatively isolated and to
incorporate them into a whole. What has formerly been experienced in connection,
but is now experienced only in isolation, is filled out to the earlier whole, and
this can be further filled out and widened when new experiences show the
possibility of still more comprehensive wholes. A tendency can moreover be traced
to continue a series of representations once begun according to the same rule
that has so far led from member to member, even if there are no immediate
experiences of more members than those so far given. Such transitions from given
representations to related representations will occur especially when the mind is
lively and when more energy has been released by the given representations than
can be used up in holding them fast; then the mind moves circling about them.

When, again, different series of representations are held together, new
connections or formations of wholes can arise in that a member of the one series
proves capable of being put in (substituted) in place of a member of the other
series. A single series too can be altered so that a new whole is formed, namely
when members within the series can be dropped out in such a way that the law of
the series continues to hold for the remaining members. In these ways arise
inferences, which are logical formations of wholes.\footnote{See more on this
\textit{Den menneskelige Tanke.} p.~168--173.}

In the course of such processes of representation the elements of feeling
behave in different ways.\footnote{Cf.\ with regard to what follows the section
on feeling (ch.~VI) in my Psychology.} Often they are drawn, through
connections of representations, over from their original causes or objects to
others. A shifting of value thereby becomes possible, since what at first had
value merely as means can acquire value as end, or what stood as end can in
future come to stand as means to a more remote end. But it is also possible that
the elements of feeling hinder or check connections of representations that,
seen purely objectively or intellectually, \opage{18}would lie near at hand. A
representation once tied to strong feeling is isolated by it, and if association
is possible at all, only such representations will have a chance of asserting
themselves as immediately or mediately favor the holding fast of the given
representation. In addition, an immediate influence can undoubtedly be exerted by
elements of feeling tied to one circle of representations on elements of feeling
tied to another circle of representations. Where two circles of representations
are simultaneous or in immediate succession, not only the representations
themselves but also the elements of feeling tied to them can come into immediate
contact. And here several possibilities present themselves. The contrast may be
so strong that tension arises and the unity of consciousness is perhaps even
threatened with dissolution. But it may also be that the one mood prepares the
way for the other precisely by exhausting the force available in a certain
direction, so that the opposite direction is favored as long as there is any
force left at all. In that case the one mood favors the other by the effect of
contrast. It may further be that the feeling strikes so deep a root in the mind
that it does not cease even when the occasion that called it forth ceases, but
spreads over the whole state and drives out the feelings tied to other
experiences, if such have been able to assert themselves at all.

But besides these possibilities there is yet another, which is of special
interest for the following inquiry. Elements of feeling tied to different
experiences can enter into connections with one another, so that complexes of
feeling arise. Experience shows, as will be developed more fully later in a
special context, that there are two main forms of such totalities of feeling or
total feelings. The connection of the different elements of feeling can have the
character of a \emph{fusion}, in that by virtue of the connection a new feeling
arises with a quality different from that of the preceding feelings. Only by
exploring the developmental history of the new feeling can one discover which
elements have taken part in this formation; to immediate observation it stands as
simple and undivided. Another kind of total feeling has the character of
\emph{organization}, in that the independence of the \opage{19}different
elements of feeling is not abolished, but they order themselves among one another
according to their relation to and significance for a ruling interest, a
fundamental value, which can find expression in a guiding thought.

There are analogous forms of wholes within the other sides of the mental life.
An analogy to fusion within the psychology of cognition is offered, for example,
by the way in which to a series of sense impressions following one another very
quickly there corresponds only a single sensation, whereas with a slower
succession a series of separate sensations would arise. In immediate recognition
there takes place a fusion of the momentary sensation with the after-effect of
earlier sensations of the same kind. Representations that have frequently
appeared in immediate connection can come to form a single representation; so in
the passage from discursive inquiry to intuitive knowledge. An analogy to the
organization of a group of feelings is formed by sense perception, in which a
multiplicity of contents of sensation are ordered in time and space, and the same
holds of images of memory and fantasy. In the domain of the life of the will we
have an analogy to fusion in instinct and drive with their immediate
concentration, and an analogy to organization in the way in which the thought of
an end to be realized determines the order of rank among the means to be used or
the actions to be performed; they all belong to the decision, are enclosed by it,
and make up an articulated whole within it.

All forms of wholes in the domains of cognition and feeling will in the end be
conditioned and determined by the striving for persistence and development that
asserts itself in every living being, that is, by a willing in the widest sense
of this word. What is to live must appear as a more or less closed whole. Pleasure
and pain are in the end only symptoms of whether this striving is furthered or
checked; sense sensations serve to release reactions toward the surrounding
world; and representations and connections of representations serve to hold fast
and shape ends and means where involuntary need and ability do not suffice.
Spinoza remains right in his assertion that we do not strive for, will or desire
\opage{20}anything because we regard it as good, but we regard it as good
because we strive for, will and desire it.

It is therefore no accident that total states are formed. How far the formation
of wholes advances in individual people, and in the individual person within a
single period of life, will depend on the special inner and outer circumstances.
The mental life hardly ever becomes fully totalized. Partly there will be single
states that stand relatively isolated, partly different total states will be
able to clash with one another. Therefore there is psychical work to be done as
long as life lasts.

4.~Total states may well have been able to form in human consciousness at all
times. But they will appear more easily when experiences become richer and more
manifold. Therefore they will be more prominent in recent times than in older
times, and they will be more marked, present more nuances and more individual
differences than in earlier times, when life was led within a narrower horizon,
when the very picture of the world was limited and surveyable, and when
historical experience too did not open wider vistas and numerous possibilities to
view. One must, however, take good care not to make the contrast in this respect
greater than is justified. What asserts itself with full clarity when
circumstances bring contrasts and differences with them may have been present in
a more undifferentiated form at earlier stages, just as one can see in the embryo
indications of the organs later fully developed. If the tendency to the formation
of total states has its ground in the very nature and conditions of the mental
life, it must be assumed to be present at all stages of the development of the
mental life. It would be an important task for comparative psychological
research to investigate this more closely.

Here it must not be overlooked that the appearance of total states is not
necessarily accompanied by the ability to describe and analyze them. Complex
states are often characterized by an awkward psychology by a single feature
that is thought to include the whole within itself. And since the intellectual
side of the mental life lies most clearly in view, primitive psychology will
speak most of thought, reason and knowledge, which for it then mean not merely a
single side \opage{21}of the mental life but comprise everything, including what
to a more penetrating analysis stands as particular sides of the mental life. The
personality and activity of Socrates, for example, become unintelligible if one
does not hold fast that by the knowledge in which he saw the foundation of right
action he meant ``something strong, guiding and ruling,'' ``strong enough to help
man through life'' (Plato's Protagoras p.~357 B C). A distinction between
knowing and willing, between head and heart, asserted itself only later; in any
case the ability to grasp and characterize relations of cleavage and conflict was
lacking. Of total states, again, those that have arisen by fusion will be
especially hard to describe, whereas the form of organization lets itself be
analyzed more easily.

Total states will as a rule not have the firm and balanced character of single
states. The contrasts that are overcome and joined in them can break forth anew.
They correspond, after all, to several experiences, not to a single one, and the
concentration therefore does not easily become complete. There is often a
relation of tension among the different elements, which can lead to dissolution,
or at least to transformation into new forms, perhaps through the separating out
of elements that could not be reconciled with the others. A person can perhaps
preserve his inner coherence only by forgetting, so far as possible, a single
particular experience, pushing it back into the unconscious. Lear violently
pushed away the memory of his daughters' ingratitude so that it should not drive
him mad. O, that way madness lies; let me shun that; no more of that! (King Lear
III, 4, 21). It is not only real experiences that can cause cleavage or
dissolution. Possibilities too play a great part, the more manifold the
representations are, now enticing, now frightening, now both. ``Only he,'' says
Kierkegaard,\footnote{\textit{Samlede Skrifter}. IV p.~422. --- In
\textit{Begrebet Angst} (where the statement quoted is found) and in his
posthumous papers Kierkegaard laid the greatest weight on the enticing power of
frightening possibilities. They cause dizziness; they work by the mere
consciousness that they exist.}
``who was formed by possibilities is formed according to his infinity.'' As long
as life lasts, the possibilities that seriously present themselves to us also
belong among our experiences, and they \opage{22}can be harder to master than the
firm and limited experiences we call the real ones. It is out of such
possibilities that great poets create their dramas. But in all of us inner dramas
can be played out, with knots and resolutions, with collisions and decisions,
dramas perhaps not even suspected by those around us, but which can nevertheless
leave their deep traces in the further development of our character.

The figures and fates the poets present to us also belong among our experiences.
The very fact that I must admit that life holds something such as what is
described, whether it be great or small, good or evil, can become of decisive
significance for me. It is not necessary that I should myself have experienced
something similar, provided it stands as something that can have its place in
reality. It has been thought that the consciousness that we ourselves are not
part of what we read or see, but sit in full safety as readers or spectators, is
an essential element in aesthetic feeling. But at any rate not everyone has the
calm consciousness of himself as reader or spectator alongside the apprehension
of the poet's presentation. In my personal experience such a consciousness would
make full absorption and enthrallment impossible. And since we learn to know
ourselves not only through the way in which the real events of life affect us,
but also through the way we are moved by the possibilities that our own fantasy
or the poets' presentations lay before us, it is with full right that we speak
here of experiences. It can nonetheless of course have its significance to hold
fast the difference between experiences owed to possibilities and those owed to
realities. The former undeniably contain more temptations to self-deception, to
regarding oneself as a hero without being one, or to thinking too little of
oneself in fear, perhaps throwing oneself away altogether.

In accordance with this view I cannot distinguish as sharply as Alfr.\
Lehmann\footnote{\textit{Die Hauptgesetze des menschlichen Gefühlslebens.}
Leipzig 1914. p.~312--316. --- It is regrettable that this new, enlarged and
altered edition of ``Følelseslivets Hovedlove'' is not available in Danish.}
between autopathic, sympathetic and \opage{23}aesthetic feelings, so that the
last would be characterized by the calm consciousness of standing wholly outside
what fantasy or art lets appear before us. There are undoubtedly states in which
the differences between these three kinds of feeling have vanished, since our
real or central self can consist precisely in, be made up of, the living fellow
feeling with real or imagined personal beings, and only later returns to its
isolated existence, so that this will nevertheless have undergone a change
through the states of self-forgetfulness. Our real self can swing between
self-forgetfulness and self-consciousness, and a swing in the direction of
self-forgetfulness is often the most fruitful. Lehmann does in the end recognize
transitions between the three kinds of feeling, but this recognition is not
sufficient. The main question is whether he is right that the representation of
myself as different from everything else and everyone else, that is, the
representation of the self, has the significance and the omnipresence he
ascribes to it. As already remarked above, there can be a real self without a
representation of the self. It is another matter that in psychology, when it is
asked on what we ground the right to speak of a self at all, we refer among other
things to the real self in the sense in which we have taken this word above. ---

The question whether it is a good thing that relations in the domain of the
mental life become more unstable with advancing development shall not be
discussed here. Great dangers and great goods often go together. The greater the
fullness to be held together, the more energy is needed for concentration to
become possible. It can become a question of to be or not to be that the
possibility of concentration poses.

In this respect the two kinds of formation of wholes indicated in what precedes,
to which I shall return in what follows, are not alike. In fusion the elements
enter into a closer connection than in organization. In fusion a new quality
arises, a new immediacy; the contributing elements have lost their independence
and with it the ability to rebel. The dramatic character the inner life can take
on by virtue of the mutual contrast of the elements presupposes that no fusion
has taken place.

\opage{24}It is in the investigation of total states that psychology has its
highest and inescapable tasks. Psychologists like to busy themselves by
preference with single states, and still more with the elements that can be
separated out of them. The question then becomes whether total states, complex
psychical phenomena, can without more ado be derived from the elementary
viewpoints or laws that can be found by analysis or experiment. When this cannot
be done, one often loses interest in total states and acts as if they did not
exist, or leaves it to poets and preachers to occupy themselves with them. It
seems clear, however, that here as everywhere we must proceed by two paths if we
want to gain knowledge, namely not only try to push forward from the simple to
the complex, but also, conversely, observe the complex phenomena
in their peculiarity and in their historical connection, fix them by
description, and possibly through analysis reach back to viewpoints won by the
investigation of the more elementary phenomena of the mental life. Perhaps it
will then turn out that the elementary is not so simple after all. It is a
prejudice, which in many domains pokes its head out again and again, that ``in
the beginning'' were the elements, and that they are, moreover, ``the real
thing.'' But in all the beginnings we can observe, it is states, not elements,
that are before us, and states with a more or less definite stamp of wholeness.
The processes through which total states come to be can best be studied at more
advanced stages of development, where the conditions are better known and
consciousness is more awake. And then it is not only total states that can be
illuminated by analogy from single states, but also the reverse. Just as little
as a total state can be conceived as a sum of single states, just as little can a
single state be conceived as a sum of elements. The historical character of the
mental life, its dependence on particular conditions of the time, appears by the
nature of the case most in total states. The historical method must therefore, as
already remarked, supplement the other psychological methods, and psychology
must use what the history of culture and of literature can teach us.

It is an investigation of this kind that is the aim of this treatise. But since
the total state to be investigated \opage{25}gets its chief peculiarity from the
elements of feeling it contains, we must first dwell a little on the life of
feeling and see how what has been developed so far applies there.

5.~By a \emph{single feeling} I mean a mental single state in which the
elements of feeling are predominant and decisive, and which is not only
determined by a single experience but is also relatively simple in its nature.
Such a single feeling is, for example, the agreeableness of a color or a scent,
the joy in a thought or a picture, the involuntary approval the sight of a brave
action can call forth, the satisfaction of having oneself acted with skill. Joy
and sorrow, hope and fear, anger and love can appear as single feelings.

Such single feelings can enter into connections from which \emph{total
feelings} can proceed. A total feeling presupposes experiences of different
kinds that stand in mutual connection. An experience that is objectively single
can, however, have such different effects through the different qualities it
presents that subjectively it comes to stand as several different simultaneous
experiences. Total feelings can be called forth in single states (see the
examples given in § 1), but only when they have first developed as total states
through a series of experiences. There is an analogy here to partial
recognition, in which we know a composite object by a single part or quality,
when we have earlier had sensations corresponding to the whole group of parts or
qualities. (Cf.\ Psychology V B, 2).

Single feelings already have a different character in every individual, every
people, every class and every time. But still greater differences appear in total
feelings. ---

I have already stressed the difference between single feeling and total feeling
in my Psychology (VI B, C and E), though without using these designations. I find
it first indicated in the literature in Plato at the place cited above, and, as
will be discussed later, it asserts itself in a plain and characteristic way in
the most eminent minds of the Renaissance. In recent times it was first described
clearly by F.~C. Sibbern (Psykologie 1856 p.~140. 380), who lays weight both on
the sporadic \opage{26}appearance of feelings in states of transition and crises
and on the possibility that there can arise what he calls a ``mixed feeling,'' in
which conflicting feelings have entered into so intimate a connection that they
work together to form one single total feeling. This phenomenon I would rather
call fusion, especially as ``mixed feelings'' in everyday speech is used roughly
in the sense of discomfort. But there is another kind of total feeling, which has
the character of organization in the sense in which we have used this word above.
Different feelings can namely work together under the guidance of a single
feeling that determines the value or the end in whose service they stand. This
kind of total feeling was well described by Th.~Ribot in his book ``Essai sur les
Passions'' (1907), where the word ``emotion'' is used of single feeling, the word
``passion'' of total feeling. And Alexander Shand has depicted in a masterly way
what I call organization, under the name of systematization, first in his paper
``Character of the Emotions'' (Mind 1896), later in his work ``The Foundations of
Character'' (1914). % the Danish has no closing quote or full stop here

We shall now consider more closely the two forms in which total feelings can
appear.

6.~Fusion is the most intimate connection different feelings can enter into,
and there are many intermediate forms or stages before this degree is reached.
The mind can surge back and forth as different emotions succeed one another.
Socrates' disciples now laughed, now wept during the last conversation with their
master. The situation was, after all, that his death was close at hand, but that
they nevertheless still carried on their usual work of thought together with him.
There was, says Plato (Phaedo p.~59 A), ``a strange mood,'' a mixture of sorrow
and joy. That tears and laughter alternated shows that within the mixture now the
one, now the other element had the upper hand, and that it did not come to any
unified, new feeling. There were plainly two different causes here, the thought
of the master's departure and the philosophical conversation, each of which set
its own feeling in motion. When the disciples later thought of their master, the
different emotions \opage{27}may have been able to even out with each other, so
that the representation of the loss and the representation of the incomparable
personality could work together to call forth a total feeling of wistful
admiration; a fusion may have been able to be completed.

A fusion can often take place only in single, especially concentrated moments,
while in other moments the single feelings alternate or clash with each other.
Wistfulness, for example, is then not the ruling feeling; it gains mastery only
under favorable conditions, while otherwise sorrow at the loss and joy at the
worth of what is lost succeed each other. A person's character --- the real self
--- does not come forth undivided in every single state at all. It is at the high
points of our life that we are really ourselves, i.e., are ruled by the total
feeling that has formed under the influence of our experiences. At other times
partial, more or less peripheral selves rule. And the relation between the
central and the peripheral can be very different in different natures. Often only
a very close acquaintance can find the threads that connect the peripheral and
partial with the central and total.

It is not always the case that a definite difference asserts itself for us from
the outset between causes that affect us in different directions. Their effects
on feeling can join from the very first, so that the total feeling arises
immediately, without being prepared through oscillations between different, more
or less opposed feelings. Thus a picture can at first sight awaken a total feeling
whose elements can only be traced afterwards. Examples of this are afforded by
Velázquez's pictures of the dwarfs he met at the Spanish court. These little
fellows stand out as comic through their unmistakable self-esteem in contrast to
their smallness, but the comic is complicated by a melancholy expression in the
eye, which introduces an element of sympathy into the beholder's mood, which can
thereby become of the kind we shall be occupied with in what follows --- humor,
just as humor may perhaps have lain at the ground in the artist when he went to
his work. Julius Lange remarks of Velázquez's portraits of dwarfs that, however
strongly characterized the dwarf figures are, the master's noble sense of beauty
nevertheless shows itself in his \opage{28}eye for the point where the misfortune
of the body has done the soul wrong.\footnote{\textit{Udvalgte Skrifter} II
p.~142.} It cannot of course be decided with full certainty whether the two
elements were joined into a total feeling in the artist's soul. As regards the
beholder too there may perhaps be disagreement as to whether the total feeling
arises immediately. It might be, after all, that the beholder first looked at the
figure and then dwelt on the gaze. For my part it was so that I at once caught
sight of the gaze as the center of the whole figure. Only afterwards did I spell
my way through the different parts of the picture. I must add that I have not
seen the originals.

The clearest and sharpest apprehension of contrasts we do of course get by
observing them successively. All else being equal, successive difference works
more strongly than simultaneous difference. (Psychology V A, 5). But William
James\footnote{\textit{Principles of Psychology} I p.~495.} is hardly right when
he goes a long step further and asserts that two terms or elements cannot be
apprehended as different in simultaneous apprehension if we have not
successively apprehended each of them by itself and thereby got the peculiar
transitional feeling of difference (the transitional sensation of difference).
No analysis of a phenomenon of consciousness, however complex, takes place, he
maintains, unless some of its elements have undergone change. --- To this I would
first remark that there are several intermediate forms between immediate
intuition and the analysis James speaks of. Even if a change takes place in one
part of a perceived picture (for example by a single feature being looked at more
closely or being lit more strongly from outside), the picture is not therefore
dissolved. The effect is perhaps only that the picture is, as it were, sharpened
in a single direction or acquires a new center of gravity. The intuition is
articulated, jointed, but is not dissolved in its simultaneous embrace of the
picture's content.\footnote{Cf.\ on articulating intuition as an intermediate
form between intuition and judgment \textit{Det psykologiske Grundlag for logiske
Domme.} (Vid.\ Selsk.\ Skr.\ 6th series). p.~18--21. (More briefly in \textit{Den
menneskelige Tanke.} p.~46; 53--56).} --- Next, at any rate as regards the effect
on feeling, it is not the apprehension of difference itself that is decisive, but
precisely the \opage{29}circumstance that the different elements work together
despite their difference. To the complex influence there may very well correspond
immediately a complex effect on feeling. And it may then be precisely wonder at
this effect on feeling that makes us seek, by successive observation of the
picture's single parts, that is, by analysis, to bring out the cooperating
elements each by itself. The apprehension of difference proper then becomes
secondary, but the effect on feeling is not for that reason secondary in relation
to the effects on feeling of the single elements. --- It must of course be
admitted that it is hard to distinguish between simultaneity and very rapid
succession. And the fact that we can gather our attention on a subject only for
very short moments, that is, that attention vibrates, also points to successive
apprehension setting in very quickly. ---

As examples of fusion Sibbern names two groups of phenomena.

For both groups it holds that after a relation of tension between two feelings a
total feeling arises, so that the two simpler feelings are no longer felt each by
itself. In place of the feelings that clashed with each other or succeeded each
other, a feeling of a new quality has appeared, which would nevertheless fall
away if one of the joined contrasts fell away.

The one group of examples Sibbern describes as the feelings of subjective
satisfaction. They are owed to the fulfillment of a need, the filling of a want,
the relief of discomfort, the overcoming of resistance. The feeling in the
suspenseful and the mysterious, in danger and in struggle, belongs here too.

Against the first-named examples of this group Alfr.\ Lehmann objects that since
the need falls away on fulfillment, no fusion takes place here. Sibbern could
have answered that every feeling has a tendency to persist after the original
occasion has fallen away. It spreads by a kind of expansion over the following
states and can thereby come to fuse with the elements of feeling of these states.
(Psychology VI B, 2; E, 4). Having felt a need is not necessarily a purely past
phenomenon. Linguists too teach, after all, that a perfectum absolutum besides
the element of time also \opage{30}contains something else, namely the
representation of the result, of ``what survives, the permansive.''\footnote{%
\textit{Otto Jespersen: Tid og Tempus.} (Vid.\ Selsk.\ Oversigt 1914). p.~320 f.}
And Sibbern speaks expressly of the peculiar joy of \emph{having escaped} from a
danger or distress, or of \emph{having} got through it. There is thus something one
continues to have when a danger has been got through. So too he speaks (in the
second part of Gabrielis' Letters) of ``absorbed unblessedness'' as different
from unblessedness merely got through. ---

The other group of Sibbern's examples contains feelings such as wistfulness,
reverence, the feeling of the sublime. That they can have the character of fusion
cannot be doubted. To this group also belong, as will appear later, many forms of
humor.

7.~The other form of connections of feeling appears when a series of
experiences release feelings that do indeed stand in close mutual connection,
but do not wholly lose their peculiarity each by itself; rather they form a whole
whose inner order is determined by a single value, a guiding end as organizing
force. Such an order can arise when a fundamental value asserts itself that for
its full realization demands a series of means and paths, as well as a more or
less conscious striving. In the clash of the fundamental value with the
circumstances of reality a series of different feelings will then be released,
which lie, so to speak, on the way to the real goal. There can be people whose
whole life is determined by such a fundamental value and its relation to the
reality within which it is to be carried through. For other people it is only
certain domains of life, for example art or science or social life, where a
fundamental value exists for them at all. For others again each period of life
has its fundamental value, youth its own, later age its own; or it is perhaps only
in a single period of life that one can speak of a fundamental value at all. Every
serious relation to a cause, to a person, to an endeavor (art, science, care for
others, social and political striving for freedom) will naturally give rise to an
organization of the life of feeling, so that the \opage{31}single feelings get
their definite place in the mental household according to their significance for
the maintenance of that relation. The single feelings that the ordinary course of
life brings --- joy and sorrow, hope and fear, goodwill and anger --- will be
furthered or checked by the work in the service of the fundamental value, and
this work itself will give occasion to the arising of a series of single
feelings.

The connection into which the single feelings enter when an organization takes
place around a ruling interest will bring about changes, nuances in the character
of the single feelings. Partly the direction in which the whole striving goes
will gain influence on the kind and degree of the single feelings, partly their
mutual interaction, when they meet in the course of the endeavors, will give them
a peculiar character they would not get under other circumstances. Anger, for
example, will have a different character when it has an egoistic source than when
it is owed to the injury of an altruistic feeling, and again it appears
differently when it is released together with or just after fear than when it is
released together with or just after joy --- such fear or joy as the work in the
service of the fundamental value produces.

Alexander Shand, who in his above-mentioned work ``Foundations of Character'' has
described in an excellent way the organization (or, as he calls it,
systematization) of single feelings (emotions) into total feelings (sentiments),
nevertheless seems to deny the influence that its entering into an organization
must have on a single feeling's character. The disinterested character of a
feeling depends, he says (p.~43--50), not on the feeling itself but on the cause
that calls it forth, or on the system to which it belongs; anger and fear can be
just as disinterested as pity and sympathy when the cause that calls them forth
does not fall under the individual's own need of self-preservation and the
systematization that the rule of that need brings about. This statement is not
entirely clear, but it seems to mean that the feeling itself does not acquire a
different character by serving different ends. The character of disinterested or
selfish would then concern only the end, not the single feelings that the work
for the end calls forth. But does not fear take on a different \opage{32}nuance
when one is afraid for one's own person and when one is afraid for the welfare of
others or for a cause one is taken up with? Both kind and degree will surely be
different. The very difference of the end will make different elements enter into
the feeling of fear in the different cases; here too fusion or organization can
take place, and a new timbre can thereby arise in the single feelings. There
cannot be an altogether indifferent relation between the feeling itself and the
connection in which it occurs. Manifold relations of contrast and of evening out
will assert themselves here.

It is surely connected with Shand's view of this point that he does not properly
recognize the first main form of total feeling discussed above. What he calls a
feeling's ``system''
is the sum of the thoughts the feeling is tied to --- the mood (feeling) and the
tendencies that assert themselves with it --- the organic processes tied to it
--- and the conduct or way of acting that follows from it. But these elements are
not supposed to undergo any change by entering the ``system''; the theory of full
fusion into a qualitatively new feeling would, according to Shand (p.~55), be
unable to explain how a total feeling can appear in very different ways under
different circumstances --- the presence or absence of its objects, absorption in
expectation or in memory. A fused feeling (compound emotion) would not appear in
the same way; only ``systematized'' feelings (complex emotions, --- emotions
combining into a system) can have the plasticity that experience proves for
certain feelings (p.~55; 172). Shand is right that practically there can be a
great difference between the two kinds of total feeling; but that does not
exclude that we have here two peculiar forms of total feeling, and that is what
matters here. For the rest, each may have its advantages. Fusion makes a more
concentrated mood and conduct possible, while the form of organization has room
for greater variations in the single cases. The greater plasticity of
organization can, however, very well be combined with the elements of the
``system'' mutually determining one another's nuance or timbre. Every element
makes certain demands, but demands are also made on \opage{33}it. It stands in
both a negative and a positive relation to the whole into which it enters. The
anger that a father's care for his son's upbringing has called forth finds its
limit in the end. In choleric fathers a struggle may perhaps be needed to keep
the anger within this limit. In criminals there can sometimes stir (cf.\ § 2),
while they carry out their act, a cruelty that is not conditioned by the proper
end of their conduct. Likewise in soldiers in war. It is often fear that makes
cruel, since it makes it impossible to stop once one has begun to lash out to
make oneself safe. The elements that are organized can have a constant tendency
to break out and stir unruled and unlimited by any goal except the one: to rage
themselves out! But the timbre changes in the very moment they break out. The
negation or limitation then falls away. ---

Not all single feelings can enter as members into total feeling, in the one form
or the other. If they could, human mental life would have a less chaotic and
disharmonious character than it has. There are joy and sorrow, indignation and
admiration, hope and fear, faith and doubt that do not enter as members into any
firm mental connection.

At the ground of Ingemann's splendid poem ``Der strømmer en Flod mod det evige
Hav'' (``A river streams toward the eternal sea'') lies the optimistic thought
that all the elements of the mental life can be harmonized. It is said of that
river:

\begin{quote}
Each sorrow it turns into wistful delight,\\
each sigh into melting \emph{song} --- --- ---\\
Even the stone that weighed like a mountain\\
dissolves into golden grains.
\end{quote}
% Danish: Hver Sorg den forvandler til Vemodslyst, / hvert Suk til smeltende
% \emph{Sang} --- --- --- / Selv Stenen, der knuged med Bjergevægt, / opløses til
% gyldene Gran.

Yes, this can happen. But it does not always happen. There are sighs that are not
turned into song, and there are stones in which no gold is found. And since every
total feeling has its dangers, in that the consequence can be closure,
unreceptiveness to new experiences, it is perhaps good that not all mental chaos
is harmonized.

Two great minds have spoken of how contrasts can maintain themselves in the mind
and divide it without either fusion or organization taking place. They both use
the same \opage{34}image to illuminate their experiences. Thus Carlyle says
(Sartor Resartus II, 9): ``In the forecourt there is dancing to the guitar, but
within there sounds now and then a soft lament.'' And Kierkegaard writes in his
journals (Papirer IV p.~35): ``I sit and listen to the tones within me, the
music's glad beckoning and the organ's deep earnest.''

The psychical process that in many leads involuntarily to harmonious connection
is impossible in natures like these two men, and in them it is often only by a
violent effort of will that a unity is brought about, the chaotic streams being
forced into one bed, where one can nevertheless constantly notice the difference
between the waters of the different springs.

8.~The tendency to the arising of total feelings is connected (as already
indicated in § 2) both with the innermost nature of the mental life and with the
demands the struggle for existence makes. At all stages and in all forms of
mental life a need and an ability assert themselves that aim at gathering
together, at synthesis. There works here, as Kant expressed it, a hidden art in
us, without which we would not have consciousness, but of which we only rarely
become conscious. It is the task of general psychology to show this in detail; it
was at any rate the task I set myself in my presentation of psychology published
in the year 1882. The mental life consists in and under a constant work of
gathering together and ordering a chaos. In the domain of the life of feeling too
such work is done involuntarily, and the processes that lead to the formation of
total feeling are interesting testimonies to it. There are perhaps, as we have
seen, natures in whom different interests or tendencies of mood can lie side by
side, almost independent, or so that they get the mastery by turns. But apart
from the influence outer impressions can exert, a natural alternation will assert
itself, a strong swing in the one direction giving rise to a need for a swing in
the other direction, and herein it already appears that the single elements are
subject to the demands made by the maintenance of the whole mental household. But
there will moreover for the most part be one interest, one value, which in the
long run regulates \opage{35}the appearance of the other values and determines
them more than it is determined by them. Which element it is that determines the
timbre of the fusion or the fundamental value of the organization will not always
be recognizable with certainty by the individual himself through
self-observation, but it will appear from his conduct, the way in which he takes
life and intervenes (or does not intervene) in life. It is, as I have stressed in
my Psychology, often only in a deep-going deliberation, undertaken with full
energy, in a definite situation that demands action, that it appears what our
real self, our fundamental value, really is. But it may very well be that the
real self that formerly ruled in us is changed in a decisive way in a
deliberation that sets us before the realities of life, so that what we discover
in such moments is precisely not what till then lay at our ground. When one
immerses oneself in this consideration, one gets a small glimpse of what endless
interaction goes on in the innermost of our being. Could a simple intuition help
us once and for all to clarity here, it would be easier to live. But life would
then in return lose its fullness and its tension.

When there thus lies in the very nature of the mental life, as experience shows
it to us, a tendency, at once a need and an ability, toward collection and
concentration, this agrees with the demands made by the circumstances of life.
The measure of a living being's stage of development consists in good part in the
degree to which it can appear as a whole over against its surroundings. The plant
can do this in a higher degree than the stone, and the animal in a higher degree
than the plant. Within the animal world the significance of the nervous system in
the struggle for existence is connected with its making whole and closed action
possible. The animals we call higher have a more concentrated nervous system than
those we call lower. And to this physiological concentration there corresponds
precisely the working of the mental life as synthesis. In a situation of action it
can be a matter of gathering all experiences and memories, all mood and striving,
and concentrating them in a single direction. When a total feeling has been
formed, a mobilization of the forces of the mind in a definite direction has
already taken place beforehand, and this can become of the greatest significance
for the life of action --- while on the other hand it \opage{36}can also lead to
tragic conflicts if the mobilization has taken place in a direction other than
the one in which action is to be taken. ---

In connection with this consideration it is of great significance, as Alexander
Shand shows, that every total feeling has a tendency to become the basis of an
ethics. Concentration on a value brings with it the striving to procure
everything required to hold fast and develop this value. And the judgments of good
and evil will depend on whether this striving is furthered or checked. First and
foremost it is a matter of the total feeling itself being preserved in its
freedom and strength, for then the qualities (``virtues'') that that striving
demands will come to be involuntarily. If freshness and strength are lacking,
ideals must be set up and duties laid down. Virtues, ideals and duties vary with
the total feeling. Thus love or friendship as a total feeling demands
faithfulness and honesty, parental love patience, self-control and wisdom,
intellectual interest perseverance and impartiality, etc. Each total feeling has
its conscience. The personal ethics thus developed stands in constant interaction
with the ethics determined by social demands. (Foundations of Character
p.~111--120; 172 f.).

This view of the ethical problem goes in the same direction as the one I long ago
put forward in my Ethics published in 1887 (ch.~III), and which I still hold. I
have shown in particular that there is a group of ethical concepts that all
ethical systems use, although they place them in different orders of rank
according to the fundamental value they build on. Such concepts are the concepts
goods, virtues, duties and rights (Ethics VIII, 1). Every ethical view, no matter
what total feeling it presupposes, finds application for these concepts in a
definite mutual order. These concepts can therefore be said to make up the
content of a formal ethics, which must be capable of being presented connectedly
on the basis of a comparative investigation of the different ethical systems. I
have spoken of this more closely in ``Den menneskelige Tanke'' (p.~350--359).

Since one can in fact build on different fundamental values, there is here a
limit to the rationality of ethics, a limit given by the ultimate presupposition
of every ethical system. The social or \opage{37}sociological ethics that builds
on the persistence of society as its fundamental value, whose content is the
formulation of social demands, and which often regards itself as the proper,
right ethics, is only one of the many psychologically and historically possible
ethical systems.

While I have come to this view through an epistemological investigation of the
rational justification of ethics (cf.\ Ethics III, 2, 13), Shand has come to it by
a purely psychological path, as he saw that such ``systematizations'' of the life
of feeling as he investigated had involuntarily to become the basis of the
ethical judgments passed by the people in whom ``systematization'' had taken
place in one direction or another.\footnote{In a paper on \textit{La Morale et la
Science} \textit{Henri Poincaré} (Dernières Pensées 1913) has expressed a view of
the scientific standing of ethics similar to the one set forth above.} % the Danish prints „Derniéres“

9.~As already remarked earlier, total feelings may well have formed at all
times. But they do not stand out equally strongly at all times, and circumstances
are not always equally favorable to them. As long as life unfolds within a
limited frame that stands as unchangeable, the life of feeling will have a
relatively simple character. There will be a common stamp, determined by
community within family, tribe and state. Whatever seeks beyond this frame, all
endeavors for which it is too narrow, all contrasts that threaten to burst it, are
checked and pushed back. The arising of significant total feelings presupposes
not only a great fullness of experiences and views of great horizons, but also an
independence in the individual that makes it possible for special, perhaps
unique crystallizations to take place of the feelings called forth under the
influence of manifold experiences.

If we keep to the cultural periods most important for us, the ancient and the
modern, we cannot expect to find many characteristic examples of total feeling in
the first. Antiquity was ruled by the fundamental presupposition of
unchangeableness as the expression of the innermost nature of existence and also
as the mark of the highest value. The changeable \opage{38}stood for Greek thought
in the end as illusion, and a worthless illusion at that. Even in the richest
presentation of human striving that antiquity has given us, the one in Plato's
``Symposium,'' even there this striving, with all the unrest and all the
enthusiasm it brings, stands as belonging to a lower sphere in comparison with
unchangeable being, which is the highest object of thought.\footnote{See on this
my paper on Plato's Symposium in \textit{Mindre Arbejder}. Third series.
p.~88--93.} Now feelings of pleasure and pain presuppose precisely change, and
more or less strong and sudden change, and they reach their most intensive forms
where great contrasts meet or follow one another. Pleasure and pain are not least
tied to transitions and crises in which opposed endeavors fight for mastery in
individuals and in the relations of individuals to one another. The fight is in
the end over what is to be the fundamental value and thereby assign all other
values their place. Man is drawn to opposite sides, and that there is a unity in
his being shows itself in such a conflict only by the strife and the division
being felt as pain. Pain and sorrow, remorse and doubt presuppose an inner unity;
it is precisely the dividedness that is felt as pain. Suffering is the last bond
that holds the mental life together. ``Nothing joyful,'' Sibbern makes his
Gabrielis say, ``nothing joyful did I have that could give my existence
continuity; so pain and despair became what held my existence together and kept me
from giving my soul up to being torn asunder.'' That is why in pain and sorrow one
learns to know oneself so well. Could one of the opposed terms be cast away and
forgotten, the suffering would fall away in the same moment. It is faithfulness
to oneself and to the values one has learned to know that makes one hold out in
the purgatory of suffering, or else a great hope brings it about that one does
not, despite everything, end in giving oneself up.

The chorus in the Greek tragedy does follow the hero's suffering and sorrow, but
it always returns to the wish for a life not placed in so hard a strife, and to
the resolve not to forget the limitation of mortals. And this line of thought
expressed the basic principle of the ancient view of life. Plato was well aware
\opage{39}that life was both comedy and tragedy, and that they belonged naturally
together. But this thought did not become of decisive significance for him.
Perhaps Socrates holds a special position here. This will be investigated more
closely later.

With all the reservation that must be made in considerations of this kind, one
may say that it belongs to modern times to be fully understanding toward the
manifoldness and the wealth of contrast of life. The first psychological
descriptions that bear witness to such an understanding appeared at about the
same time in Giordano Bruno's wonderful work ``De gl' heroici furori'' and in
Shakespeare's tragedies.

What Bruno calls heroic passion is conditioned by the composite character of the
higher mental life. Only in the stupid, who cannot see beyond what is immediately
before them and have no comprehensive experience, do not know the manifoldness of
existence --- only in them can feeling and need have a single and simple
character. In those who have rich and many-sided experience, the life of feeling
is composed of pleasure and pain, sweetness and bitterness, and this in such a way
that these opposed elements mutually determine each other. Ignorance is the mother
of enjoyment in the animal paradise, and he who increases his knowledge increases
his pain. But the heroic soul feels pleasure in the very pain, because this is a
sign that it yearns toward an infinite and exalted object, which must and can be
the goal of ever renewed striving. The heroic soul would rather sink in the
struggle, or at least feel its lack and imperfection over against the great goal,
than reach perfection by setting itself a lower goal.

Comparison with Plato's ``Symposium'' shows the difference between ancient and
modern thought. Plato wants to lead beyond constant yearning, beyond the world of
change in which alone there can be talk of striving. The emotions are to be
replaced by intellectual clarity, and the mind is to come to rest in the
unchangeable. But for Bruno the play of contrasts and the tension it brings are
precisely the expression of true reality, because it is the richness of life and
the infinity of the world that bring the constant changes and the great contrasts.
Just as little as the outer universe had any limit for him, just as little could
the world of \opage{40}ideas and values have a limit. The great fundamental
thought for which he became a martyr penetrated his innermost soul and determined
his view of life. --- The total feeling Bruno has described sprang from the need
for knowledge and from the experiences of life conditioned by everything that
either furthers or checks the satisfaction of this need. It was the fortunes of
advancing, independent thought that conditioned what he got out of life. His most
intensive spiritual feeling of life he had where he saw the barriers the old
picture of the world had set for spirit and thought give way to greater horizons.
As will be shown later (see ch.~IX), this kind of total feeling is the most
frequent among great philosophers.

As regards Shakespeare, there will often be occasion later to bring him forward
in connection with the total feeling this treatise is especially to be occupied
with. Here it need only be pointed out that his characters, as well as the fates he
depicts, are not built on quite single motives, but, where he reaches highest, are
each like a whole orchestra in which many motives, contrasts, resolutions and
harmonies sound together
into a mighty totality. He has thereby become a treasure trove for psychology in
a higher degree than any other poet.

10.~The difference between single feeling and total feeling shows itself in a
rather interesting way in the history of the word humor, and it is necessary to
dwell a little on this in order to understand the special meaning of the word in
which it is to be taken in the following inquiry.\footnote{Besides my own studies
I have here had much use of a paper (Addison's Humour, its matter and its form) by
the young literary scholar \textit{Mary Suddard}, who died early, in her
\textit{Studies and Essays} (1912).}

The meaning of the word is at first physical: moisture, wetness, fluid. Then it
is used in the old physiology of the fluids in the body that determine the nature
of the temperament. The well-known four old temperaments each corresponded to the
predominance of its fluid. From physiology the word passed over into psychology
to denote a certain mental state, and the psychological usage then shows
\opage{41}in its turn a series of transitions that bear witness to changes in the
life of feeling in modern times, or at least to a clearer apprehension of the
phenomena of the life of feeling. It is of course not the name that matters. All
too often psychologists still believe they can explain to us what a certain name
for a feeling or another mental phenomenon (attention, will, egoism, religion,
etc.) ``really'' means. The decisive thing is to get a feeling described and
characterized. Then one can afterwards see about finding the name for it that is
most suitable, naturally with every possible regard for traditional usage.

As regards the word humor, it can be seen plainly how the physiological usage
determined the psychological.\footnote{The transition from the physiological to
the psychological meaning seems to have taken place in France at the beginning of
the 16th century. According to \textit{Kr.\ Nyrop (Grammaire Historique de la
Langue Française.} IV. 1913. p.~228) the psychological meaning is cited by Henri
Etienne as a novelty.} Such expressions of character in which one thought one
could trace in particular the effect of a single one of the body's ``fluids'' were
called ``humorous'' in England. One thought here in particular of qualities like
stubborn one-sidedness or capriciousness or eccentricity. This usage appears in
two younger contemporaries of Shakespeare. Robert Burton (Anatomy of Melancholy.
1621. I, 3, 1, 2) describes people who are ``humorous'' thus: now they laugh from
the bottom of their hearts, now they weep without outward cause, now they sigh and
are sad and pensive, now they wrap themselves in unreasonable fancies. Ben Jonson
gives in the prologue to his comedy ``Every man out of his humour'' the following
explanation: ``When a single quality rules in a person to such a degree that it
makes all his thoughts and powers go one and the same way, such a person must be
said to bear the stamp of humor (to be a humour).'' The word ``humourist'' is used
by Ben Jonson of one who is ``humorous'' (Act I, Sc.~3). Jonson puts forward his
explanation of the word as a protest against the abuse by which every fool or fop
called himself ``humorous.'' He wanted to limit the meaning to the predominance of
a marked, temperamental quality.

When Duke Frederick in ``As you like it'' is said to be
``hu\opage{42}morous,'' the quality thought of must, from the context, be
stubbornness. When Hamlet, at the news of the players' arrival, promises among
other things that ``the humorous man'' shall be allowed to play out his part in
peace, what is probably thought of is what we would call a character actor,
particularly a performer of striking, extravagant characters. Dowden thinks (in
his edition of Hamlet, the note to II, 2, 340) of characters like Jaques and
Mercutio.

Whether by the word humorous one thought mainly of fickleness\footnote{%
\textit{Holberg} presupposes in ``the Fickle Woman'' a sudden shifting of moods,
without the preceding one having any influence on the character of the following.
This must presumably be understood to mean that Lucretia is put into a mood by the
slightest occasion, and that the mood thus aroused comes to color the whole state.
In an interesting way \textit{Johanne Louise Heiberg (Et Liv genoplevet i
Erindringen}\textsuperscript{3} II p.~27--31) has maintained the necessity of a
mimic motivation of the transition from one ``mind'' to another and has pointed to
the small indications of such transitions that the text offers on closer
examination. --- By Burton's description Lucretia would have to be called
``humorous'': ``sometimes profoundly laughing, and then again weeping without a
cause'' (Anatomy of Melancholy I, 3, 1, 2). --- A patient suffering from
successive double consciousness, observed by \textit{Morton Prince}, who after her
recovery herself described her illness \textit{(B.~C.~A.: My Life as a
Dissociated Personality.} Boston 1909), could suddenly and without outward
occasion take on an entirely different character, which showed itself in
expressions, bearing, speech, etc., and in such a way that the one state had no
memory of the other.} or of one-sidedness, it means the rule of a single feeling.
It has not yet acquired the meaning of a composite feeling arising through the
clash of opposed elements, which is prepared in the course of the eighteenth
century and definitely established in the nineteenth.\footnote{See more on this
below §§ 47--50.}

It is Shakespeare's art that has in good part prepared the way for this new
meaning and for the understanding of the peculiar mental phenomenon that lies at
its ground. Outwardly this understanding appears in the brilliant way he worked
the clownery of the older stage into the whole of a great tragic action. The
opposed currents of life unite in Shakespeare's art into a mighty river. Jest and
earnest, laughter and \opage{43}tears, loftiness and lowness --- everything makes
its contribution to the total mood the works leave behind, and which was probably
present in the poet himself as the fruit of his experience of life. He could not
have solved his artistic task if he had not in his life solved the corresponding
practical tasks. The poet need not of course have experienced himself the
contrasts he depicts. But he must have seen them as possibilities. In the
characters and fates he observed, in the moods and impulses he experienced in
himself, in the events that befell him, there must have been drawn for him
intimations, lines of cleavage, possibilities, over against which it could become
serious work to maintain the strength and harmony of the mind. What was reached
through this work he then shaped into a world of figures and fates that stands for
spectator or reader as one of the greatest pictures of life ever given.

Only from a deep understanding of human nature in its manifold forms and of human
life in its changing fortunes can one dwell on the paltry, the low, the absurd in
such a way that it is seen, despite everything, as an element in the whole, and
only on such a basis can jest have free play, although the great earnest forms the
ultimate basis.

In humor --- in the special meaning in which we are concerned with it here ---
there is expressed at once a large view of life and an honest experience of
life's pettinesses and misfortunes. The total feeling can here too appear either
as the consequence of fusion or as the consequence of organization. In the one
case the contrasts are bound, present only potentially; in the other they come
forward more freely, but are joined by their common relation to a fundamental
value.

Shakespeare's humor is, as Dowden has said, an expression of his faithfulness to
reality. He has not lost his faith in the great because life offers so much
lowness and so much suffering; nor, in his certainty of the worth of the great, has
he lost his sense for the small and the painful.

The great Spanish poet whose ``Don Quixote'' came out at the same time as
Shakespeare's ``Hamlet'' is regarded by some as the master of all humorous poetry,
because his jest is always earnest, and \opage{44}his earnest always comes forward
in the form of jest.\footnote{Thus \textit{C.~H. Weisse: Kleine Schriften zur
Aesthetik} 1867.} In his presentation of the knight-errant and his fight with
windmills, with a Sancho Panza as squire and chorus, Cervantes always has
understanding and fellow feeling for the good heart and the sincere will, although
it is only after the collapse of illusions that these qualities properly come
forward in his hero.

The elements indicated here in what I shall call \emph{the great humor} are to be
illuminated more closely in the following sections, in their different nuances and
their different mutual relations. The earnest that lies behind the jest can be owed
to sympathy and understanding, or to wistfulness and longing, or to magnanimity and
superiority. We shall come to dwell somewhat more closely on these types. ---

There is also a \emph{little humor}, the most popular form, which is one with more
or less good-natured jest. Good nature can have many degrees, through which humor
can pass over into irony, satire or scorn. It is often this little humor that is
thought of when the word humor or humorist is named. It is a single feeling, not a
total feeling. It is a single experience that is taken up for treatment in a light,
jesting form, as a rule with an undercurrent of understanding and sympathy, and in
such a way that the jester may very well draw his own person and his own conduct
in among his subjects. --- Usage also permits speaking of humor as the ability to
call forth laughter, quite apart from the disposition that lies behind it. That is
the tiny little humor.

The little humor can become a mere form that can be used from the most different
motives. It does not have the deep foundation of the great humor, which is a view
of life, or rather a disposition toward life, a total state to which all
experiences and endeavors have made their contribution. The little humor can be
used by scorn and ill will as a mere form; the earnest hidden in the jest is then
feigned. The ironist too, as will appear later, can use the little humor. Often
there can lie behind it a need to play, and through play to free oneself from what
weighs on the mind. In the single case it can be hard to decide whether \opage{45}it
is the little or the great humor one has before one, and in the first case, in
whose service it stands. In what follows the little humor will come up only by way
of exception, since it is essentially humor as a view of life or type of life, a
stage of life, that is the subject of this treatise.\footnote{To test which
meaning of the word humor would nowadays lie nearest to men versed in literature, I
asked three colleagues, each of whom from his side was occupied with literary
research, to tell me quite on the spur of the moment what representation the word
most nearly awakened in them. The first answer went mainly in the direction of the
little humor; the second took humor as a sense of the comic on the basis of
sympathy and understanding; the third characterized humor as a view of life
conditioned by life in time being seen with eternity as its background. All three
forms, and more besides, I shall come to discuss in what follows. The first answer
gave the popular meaning of the word, the second the later English usage, the third
took humor in the sense of Romanticism and of Kierkegaard.}

11.~At the outset the misunderstanding must be averted that the intention is to
prove the sole justification of a particular view of life. That would be an
impossible thing. To that extent what is often said, in season and out of season
(especially out of season), is right: that there is no disputing about views of
life. By a view of life is meant not a theory or a conception that one merely
\emph{has}, but a standpoint on which one stands, or a life one actually
\emph{lives}. It is not a sight one sees, but an eye that sees. One need not be
conscious of one's view of life. For philosophy, especially for psychology, views
of life can only be objects of description and investigation. Philosophy's
relation to views of life is not exhausted by the remark that they are different
or conflicting. It has the task of positively analyzing them, testing their
consistency and finding their origin and their course of development. There is
here for philosophy not only a psychological, but also an epistemological and an
ethical investigation to carry through. It is important that philosophy should not
push this task away from itself. Not only does it thereby narrow its domain, but
when views of life are not taken up for objective treatment, they easily creep
forward by detours and back ways and exert an influence that can be all the greater
the more unconsciously it goes on. Philosophy always has the Socratic ``Know
thyself!'' as its motto. And to know oneself \opage{46}means to be conscious of
one's presuppositions, of one's standpoint, to as high a degree as is at all
possible.

The task of the present treatise is to investigate the psychological possibility
of a stage of life founded on the great humor. When psychology occupies itself
with the more concrete and complicated states of the mental life, it must be clear
that it is dealing with types, individual forms, peculiar mental formations of
wholes. The general viewpoints or hypotheses at its command can be used, in their
interplay, to throw light on such formations; but it will be impossible to prove
the necessity of these formations arising and persisting at all. One can only
investigate the conditions of a mental whole when it is actually there. Such an
investigation can furnish useful and often highly necessary correctives to the
inclination to apply schemes set up once and for all to the domain of real mental
life. ---

It is perhaps not superfluous either to remark that the present treatise is
psychological, not aesthetic. The concept of humor does of course play a part in
the history of aesthetics, and good examples can be taken from poetic literature to
illuminate it. But such examples are to be used in this treatise only where there
is likelihood that it is experience of life that lies at the ground of the literary
presentation. Already Henry Home, who in his ``Elements of Criticism'' (1762)
ch.~XII proceeds from the older meaning of humor discussed above, distinguishes
between writers who use the humorous form and those who are ``humorists in
character,'' and whose works therefore involuntarily take on the stamp of humor.
``Swift and Fontaine were,'' he says, ``humorists in character, and their
writings are full of humor. Addison was not a humorist in character, and yet there
reigns in his prose writings so fine a humor.''\footnote{Mary Suddard too comes to
this view of Addison in her above-mentioned paper. Humor was for him essentially a
form of presentation, in part a means of safety. --- Whether Swift is rightly to be
called a humorist, in the narrower sense, will be discussed later.} % the print has no reference mark for this note
A recent French literary scholar has remarked that while in ``classical'' works
one can avoid thinking of the author, the ``humorous'' works necessarily betray the
author's \opage{47}temperament and character. (Baldensperger: Leçon d'ouverture).
This can hold only where humor is something other and more than a form. Literature
springs from life, and if the stage of life called the great humor has any
significance, it is not because it has found expression in literature, but the
reverse: it gives literature heightened worth that a significant stage of life has
been able to find expression through it. And it is possible to come to know this
stage of life by other paths than occupation with literature. ---

In what follows we shall come to give contributions to the history of the great
humor. Conceptually it was first characterized by German aestheticians at the
beginning of the nineteenth century, and here Solger is to be named first (in his
dialogue ``Erwin,'' 1815, and his posthumous writings). For Solger there appears in
humor at once the feeling of the glory of existence and of its imperfection, by
which an inner connection between the tragic and the comic is conditioned, so that
nothing stands as valuable without its being able to become comic through its
limitation, nothing as comic without being recognized in its worth! (Erwin. New ed.
1911. p.~350--354).

Before Solger, Jean Paul is to be named (Vorschule der Aesthetik. 1804), later
especially Zeising (Aesthetische Forschungen. 1855), Lazarus (Leben der Seele.
1855), George Eliot (Essays 1883) and Th.~Lipps: Komik und Humor (1898).

In Danish literature the concept of humor has been finely developed on its
aesthetic side by Sibbern and Wilkens in their works on aesthetics. But in this
inquiry I shall keep especially to Kierkegaard's treatment of the subject, in my
eyes a contribution of the highest rank. Of his relation to Solger there will be
talk later. Kierkegaard describes humor (in the sense of the great humor) as one of
his ``stages'' (Uvidenskabelig Efterskrift. 1846. --- Samlede Skrifter
VII). It has been my wish to carry further what Kierkegaard indicated in a few
strokes, but also to show the unjustifiedness of the barriers he, from his
presuppositions, sets to the significance of this stage, as regards extent, as
regards content and as regards energy.

% ===== II. LAUGHTER AND HUMOR (pp. 48--58) =====
\chapter*{II.\quad Laughter and Humor}
\addcontentsline{toc}{chapter}{II. Laughter and Humor}
\markboth{Laughter and Humor}{Laughter and Humor}
\opage{48}12.~In connection with the older meaning of the word humor, Henry Home
already pointed out that a character is called ``humorous'' only when it worked
comically. In the later meaning of the word too it has something to do with
laughter, or at any rate with the smile. Let us begin the inquiry into humor by
considering its relation to laughter.

Laughter can be investigated from two sides. One can ask about the nature of the
outer experiences that awaken laughter, and one can ask about the inner mood that
finds expression in laughter. As a rule most weight has been laid on the first
question, and it is especially this that theories of the ludicrous seek to answer.
But for our inquiry the second question is of the greatest significance. By way of
introduction, however, we must touch on the first, especially as the two questions
cannot in the end be kept apart.

Two theories have appeared here in recent times that might seem to stand in sharp
contrast to each other. --- The one lays weight on this: that when something is
ludicrous, it is because it first appears with a certain trustworthiness,
self-evidence, importance, superiority, but then suddenly reveals itself in its
nothingness and its impotence. There is thereby called forth in the one before
whom it appears first a feeling of security or admiration or enthrallment and
tension, and then suddenly a disappointment or a gloating, or at least a relief and
a release of the tension. What at first stood so self-assured and imposing now all
at once takes on the opposite character. The ludicrous, on this view, rests on an
effect of contrast in a downward direction. In short, vigorous strokes this theory
was already put forward by Thomas Hobbes. In my Psychology \opage{49}(VI E, 9) I
joined it, only with a special stress on the different character laughter takes
on according to the fundamental mood of the one who laughs. James Sully (Essay on
Laughter. 1902) calls this theory the degradation theory. --- The other theory
lays weight on the ludicrous always presenting an absurdity, a discrepancy, a
self-contradiction. In Danish literature it was put forward by Kierkegaard. Sully
calls it the incongruity theory.

Both Sully and Ribot (Psychologie des Sentiments p.~344) think that the two
theories each fit their own group of cases, and that it is impossible to find a
formula that can fit both groups. Undoubtedly each of the two theories stresses
important viewpoints. And they both lead to an understanding of the significance
of laughter in life, for example to an understanding of how laughter could be
conceived as a criterion of truth (test of truth), as Shaftesbury and Henry Home
have done. Through the purgatory of laughter everything artificial and hung on
from outside is separated from a subject; it comes to appear in its whole
nakedness, whereby either its worthlessness or the worth it perhaps has despite
everything can appear all the more strongly. Holberg conceived his activity as a
comic poet as an endeavor to separate the shadow from the body.

In reality, however, the two theories are not opposed to each other. The
incongruity theory is evidently a particular form of the degradation theory.
Contradiction is, after all, a special form of contrast. To be revealed in one's
inner discrepancy and self-contradiction is, after all, a ``degradation.''
Purely logically the self-contradictory $= 0$. But psychologically the discovery
of self-contradiction works as a contrast in a downward direction. What appeared
as something reasonable and reliable now shows itself in its nothingness. This
produces the shock that is the condition for the release of laughter. Among all
cases of the ludicrous there is on this point an analogy grounded in a common
condition. The kind and degree of the tension can be very different. A puffed-up
person who is suddenly exposed as empty and thoughtless, like a blown-up bladder
that is pricked, naturally stands before us in a different way from a profound
thinker who, absorbed in his subject, commits an error of calculation or of
thought on a \opage{50}subordinate point, even if this subordinate point was a
necessary link in the train of thought. One can of course formulate such contrasts
as contradictions: fullness--emptiness; profundity--stupidity. But the
contradiction is here at the same time contrast, and only thereby does the
phenomenon become ludicrous. When Sancho Panza eagerly clings to the edge of a
ditch, which in the dark he takes for the brink of an abyss, this can of course
well be formulated as a contradiction, an ``incongruity.'' Sancho's great terror
undeniably does not fit the situation. But it is not the absurdity in itself that
works comically; it is its contrast to the minimum of reason we always presuppose.
If we see a person in mortal terror, we trust provisionally that there is reason
for it. We take everything provisionally at face value; and when the terror is
suddenly removed, there is a possibility of laughter. Before the resolution there
was also a possibility of tears; the tragic and the comic can lie close to each
other. Laughter presupposes that those who laugh are still safe and secure on the
side of reason and power, even after the ``degradation'' has set in. There must be
something that is not dissolved by the downward contrast --- at any rate the
standpoint on which the one who laughs finds himself. With the tragic the
relations are more involved, as we shall see later.

The feeling of the ludicrous is a marked example of a feeling of relation, i.e., a
feeling that is not only, like all feeling, conditioned by its relation to other
feelings, but is determined in its innermost being by the experience that calls it
forth consisting in a contrast. In the feeling of laughter this contrast is
downward; in the feeling of the sublime it is upward. In laughter we feel ourselves
great over against the small; before the sublime we feel ourselves small over
against the great.

The contrast that plays a part here need not be a contrast between two definitely
formed states of consciousness. A contrast can be contained in one and the same
experience, which suddenly asserts itself for us without any succession being
observable. An example of this from another domain has already been discussed
above, namely the mood Velázquez's pictures of dwarfs call forth. When --- to take
an often-used example --- the sight of a little child in a man's hat works
comically, the \opage{51}contrast is not between child's head and man's hat, but
between the child with its usual head covering and the child with the man's hat,
and the representation of the usual head covering need not come forward expressly
at all; it can work as a tacit, habitual presupposition. Effects of contrast can
very well arise in relation to something that is not itself above the threshold of
consciousness. There are quiet, secure expectations so ingrown that they can form
the condition for an effect of contrast, a condition we perhaps first notice when
we wonder at the phenomenon of contrast and seek to explain it to ourselves. It is
here as when we wonder that something stands before us as familiar, and then
investigate where and when we have seen or heard it.

13.~Laughter is like the breath that blows away a speck of fluff. The relation
between breath and fluff can be very different. A whole scale can be set up here.

a.~In childish laughter, ``the frolicsome high spirits and delight in nothing,''
there need be no fluff at all, no outer occasion. Overflowing vitality and feeling
of life lead to a playful and jesting attitude toward everything. Laughter is here
the involuntary utterance of the healthy and vigorous mind. There is neither any
consciousness of inner security and superiority over circumstances nor of an outer
occasion. This element of involuntariness never wholly vanishes from laughter as
long as laughter is possible at all.\footnote{\textit{Sully} (\textit{Essay on
Laughter} p.~289--293; 427--432) is anxious for the future of healthy, spontaneous
laughter on account of pessimism, fatigue and cynicism. The reservoir which in the
past supplied the stream of national gaiety, has certainly fallen and threatens
even to dry up. The merriment of Old England developed after the common struggle
against Rome and against Spain. But now class oppositions come to the fore. ---
Perhaps the struggles of a new time could have a regenerating effect by releasing
and gathering the forces.} All laughter presupposes firm ground under one's feet,
however limited or illusory this ground may perhaps prove to be on a purely
objective view.

b.~Closest akin to this spontaneous laughter is the one owed to the lightening of
a pressure or the release from a danger that has \opage{52}threatened the outer or
the inner life. Here there is express consciousness of the fluff, of the ground of
the laughter.

At a steamer's departure a lady stands on the quay with a little boy to take leave
of people going away. She holds the boy out over the bulwark and raises and lowers
him as a greeting to those departing, calling out: ``Gustav says goodbye, goodbye,
goodbye!'' She drops him in the water. General consternation and despair. But a
sailor jumps in and gets hold of the boy, and while he is still down in the water,
he lifts the boy up and calls: ``Gustav says how-do-you-do, how-do-you-do,
how-do-you-do!'' A liberating and jubilant wave of laughter replaced the tension
and the fear. There lay at its ground as well a certain gloating over the foolish
mother's conduct and a sympathetic joy that a human life had been saved. Thereby
this example also belongs under the two following types.

c.~A still greater part than in the laughter of release is played by the object of
laughter where laughter is the expression of scorn. The object here means a
resistance that seemed or thought itself to be overwhelming, but which is overcome
and falls to the ground. While in spontaneous laughter the object has no
significance, and while in the laughter of release it is something one moves away
from and perhaps soon forgets, in scorn the object plays a positive part. There is
a need to hold fast the representation, or perhaps even the sight, of the object in
order fully to enjoy one's own feeling of power. Here a reaction asserts itself not
only against one's own preceding tension but also against the object itself.
Contempt finds vent against it again and again. The fluff has here become a
mountain one has overthrown.

d.~In humor, the fourth type of the feeling of laughter, the object plays, as in
scorn, an essential part, only the relation to it is not antipathetic but more or
less sympathetic. Here again different forms can be distinguished, different kinds
of humor, whose difference rests on the significance the object has.

α.~Where the superiority over an opponent is great and deeply grounded, the work of
overcoming the resistance will stand as play, and there will be no need to dwell
on it again and again. Real superiority rests in itself, and its
\opage{53}self-esteem does not take on the character of contempt toward other
things and other people. The need to feel scorn or contempt and to give them
expression is precisely a sign that one is dependent on the outer object and can
feel oneself rightly only in contrast to it. Even if no positive sympathy is felt
for it, it can very well be recognized in its peculiar worth, and the smile or the
laughter can then be conditioned by understanding. A bridge can be built between
the one who laughs and the object of laughter.

β.~Since real superiority commands more energy than is needed to overcome the
resistance and secure its own persistence, there is here the possibility of a
magnanimity that not only understands the object but seeks to shield and protect
it, all the while its paltriness and impotence assert themselves. Magnanimity bears
with a smile not only its own burden but the burdens of others, adding its strength
to the slight strength that called forth the smile. The recognition of the worth of
the paltry and impotent is here more positive than where bare superiority rules.
There is here an inner bond, not merely an outer bridge, between subject and
object.

γ.~It is not only superiority that can rest in itself. Where the mind is taken up
with something valuable that is an object of memory or hope, the inner life can be
so rich that outer objects and circumstances stand as paltry and insignificant over
against it, without therefore being misjudged in their peculiar worth. Humor can
have wistfulness or longing as its basis. Wistfulness is directed toward the past,
longing toward the future. Against both, new events strive in vain to gain power;
but they are not therefore scorned or regarded with indifference. By virtue of the
expansion of feeling (cf.\ § 6) the individual's own mood can come to take the
interests and strivings of others under its influence, especially when these others
are themselves referred to the worlds of memory or of hope. Wistfulness and longing
do not repel, but can lead to a deep understanding of everything that misses and
longs.

δ.~Humor takes on a peculiar character when an understanding of the great laws of
life and existence forms the background. This understanding need not be strictly
scientific. A person who uses his reflection and has comprehensive experience will
\opage{54}naturally reach the conviction that there is a great order of things on
which all events, even those that cut most deeply into human weal and woe, are
dependent. And this conviction will furnish a background, a horizon, in relation to
which what occupies people's interest can stand as paltry and vanishing.
Superiority is here conditioned by understanding. But precisely because the paltry
and vanishing, just as well as the great, belongs within the lawful connection of
existence, precisely for that reason it can be regarded in its peculiarity and with
the sympathy of understanding.

ε.~The relation between subject and object in the domain of laughter becomes most
intimate when the subject harbors a positive sympathy for the object. Before the
smile or the laughter there can be given a bond that does not break because the
object is seen in its paltriness or its self-contradiction. Such a bond exists
between parents and children, between brothers and sisters, between friends, and
everywhere that love of mankind and a sense for human mental life rule. The
superiority is here the superiority of love, and the love is the magnanimous love
(Ethics IX, 1; XII, 2) that is akin to justice. The understanding is here the
understanding of sympathy (while the preceding form of humor could be
characterized as the sympathy of understanding). The stronger the bond, the more
freely smile and laughter will be able to stir. The forces that could otherwise
repel --- wit and the comic sense --- can here come to bear witness to the
firmness of the bond and perhaps tie it more firmly. To this kind of humor
Thackeray's and George Eliot's definition of humor as the unity of wit and sympathy
(wit and love) applies.

All these forms of humor can appear as momentary states (emotions or ripples) or as
more lasting states (moods) or as the expression of the real self (disposition).
This last form is what in this treatise I call the great humor.

In my Psychology (VI E, 9) I took account only of the form of humor that can be
designated as the unity of wit and love. James Sully (Essay on Laughter. p.~306)
objected to my description of humor that living sympathy will prevent laughter.
This \opage{55}cannot be denied. But whether it happens depends not only on the
degree of sympathy but also on the degree and kind of the contrast that is present.
There can, after all, be contrasts so sharp and bitter that they cause wounds that
cannot be healed. An overwhelming fate can rule, over against which all the power
that can be summoned is impotence. To this possibility I return in a special
inquiry that is of decisive significance for my whole view, namely the inquiry into
the relation between humor and the tragic (ch.~V).

For the time being it follows from the characterization given of the different
types of laughter that what a person laughs at depends on what kind of person he
is. A person's character reveals itself through his relation to the ludicrous.
Good observers have seen this at all times. It will not be without interest to
cite some statements on it.

Plato demanded that what one found ludicrous should be determined not by what the
senses show us but by what thought sets highest (Republic V, p.~452 D). % the Danish prints „forlangte,.at“

Thomas Hobbes says (Leviathan cap.~6): ``Much laughter is a sign of a small mind
(pusillanimitas). It is proper to a great mind, on the other hand, to bring help
and to free others from contempt, and to compare itself only with the greatest
men.''

``All my life,'' says Lichtenberg (Vermischte Schriften, I p.~173), ``I have found
that a person's character cannot be better known than by a jest he takes amiss.''

Goethe says (in his ``Maximen und Reflexionen''): ``People make their character
known by nothing more clearly than by what they find ludicrous. The sensual person
often laughs where there is nothing to laugh at. Whatever then makes an impression
on him, his inner pleasure shows itself. The man of understanding finds almost
everything ludicrous, the man of reason almost nothing.'' --- And in one of his
epigrams he says (which confirms Lichtenberg's statement and is of significance for
the psychology and ethics of humor): ``He who cannot make fun of himself surely does
not himself belong among the best.''

Carlyle says (Sartor Resartus I, 4): ``No man who has once heartily and wholly
laughed can be altogether irreclaimably bad. \opage{56}How much lies in laughter:
the cipher-key wherewith we decipher the whole man!\dots{} The man who cannot laugh
is not only fit for treasons, stratagems and spoils; but his whole life is already
a treason and a stratagem.'' % Carlyle's own words (Sartor Resartus I, 4); Høffding quotes them in Danish

14.~In what precedes, humor has been described as a single member of the whole
series of forms of laughter. And for the most part humor is treated under the
psychology of laughter; so too in the presentation of psychology I have given. But
it differs from the other forms of laughter by its basis, by having in a higher
degree than they the character of a total feeling. More than the other forms it
presupposes experience of life and reflection, a purgatory in which the contrasts
of life and the conflicting forces of existence have exerted their influence on
the innermost core of the mental life. The great and the small, the exalted and
the low, the bright and the dark, everything has left its trace, and reflection
has bethought itself of the general laws of life.

Whether the basis is the superiority that excludes contempt, or the magnanimity
that bears the burdens of others as a natural thing, or the wistfulness that is
the fruit of rich but painful experiences, or the longing directed toward a great
goal, or the understanding that sees everything ordered into a great connection,
or the sympathy that feels the pulse of life in the least as in the greatest ---
there is a possibility of a smile or laughter at everything that in reality or in
appearance comes forward as contrary to what occupies the mind. The contradictions,
disharmonies, disappointments and sorrows of life fade before the basis and the
background that determine the effect of the single experiences. A main thing for
the form of laughter that has the character of humor is that the relation to life
is positive, or that the one who laughs or smiles at any rate takes himself along
here. Before the great measure that is applied, the sharp contrast between the
great and the small vanishes in the end. In the great humor we ourselves, with our
whole personality and our whole striving, belong to what finds its justification in
the poet's words
\begin{quote}
``the least of things defends\\
its place in the wreath of life.''
\end{quote}
% Danish: „det Ringeste forsvarer / sin Plads i Livets Krans.“

\opage{57}An example of this side of the stage of humor can be taken from
Kaalund's poem ``Naturens Graad'' (``Nature's Weeping''). Nature longed for her
children, who had now shut themselves up in the great cities. They then streamed
out to her into the open. But the great mother was indignant to see these dressed-up
gentlemen and ladies, and her indignation found vent in a thundershower that drove
them back to the city in wild flight and amid all kinds of calamities. But the poet
confesses: ``I myself was there, I rolled in a cab!''

There is in the great humor a double movement. It is the expression of a
self-assertion that builds on self-understanding. The saying about mastery showing
itself in limitation holds here in the fullest sense. For only he can assert himself
who knows clearly where he stands --- what he is capable of and what not --- what
he is worth upward and downward. The quality Aristotle called megalopsychia, and
which consisted in being conscious of one's own real worth, expresses right
self-assertion from the positive side.

The barriers then follow of themselves, being determined from within, not from
without. They lie quite simply where the consciousness of one's own worth stops.
When the Romantics in their day tried to set up Tieck as Goethe's equal, Goethe
declared that it was as if he wanted to compare himself with Shakespeare. Goethe
was clear about his limit upward and downward. His saying that he who cannot make
fun of himself surely does not belong among the best can be used as a motto for the
great humor. Making fun of oneself is not one with scorning oneself. The humorist
scorns himself as little as he scorns others. He holds himself as well as them up
against the great background life has set before him, and the relation to it
determines both the place he assigns himself and the place he assigns them.
Everywhere he seeks the kernel among the shells. And since the shells do not always
betray the kernel, he will be inclined to think that it is our ignorance or our lack
of sympathetic understanding (or of the sympathy of understanding) that makes us
often think we find shells without kernels. Not all nuts, after all, can be cracked
so that one can get to see what is inside them. As Goethe wrote in a young student's
album: ``Our Lord created the nuts, but he did not crack them.'' There may well be
need for \opage{58}humor toward those who think they have cracked all the world's
nuts, although in reality they have cracked only their own teeth. The great humor
will be bound up with constant seeking, and it stands in contrast to all dogmatic
wisdom, whether this comes forward in the name of sound common sense, of science or
of religion.

It was in the time of Romanticism that people came to see that in a disposition
like the one described there lay the possibility of a type of life, that is, of
something that was other and more than a special form of aesthetic feeling. The
justification of this still needs to be maintained and grounded. Hermann Lotze, for
instance, denied that humor could be recognized as a comprehensive mood of life. In
his eyes it can only be a certain kind of comic mood, and a kind that has its danger
at that, since it easily works dissolvingly toward everything exalted. (Geschichte
der Aesthetik in Deutschland, p.~389). And although Kierkegaard describes the great
humor with sympathy and understanding, for him it cannot designate a highest stage;
that he reserves for the religious way of taking life.

There is therefore every reason for a closer investigation of this subject.

% ===== III. IRONY AND HUMOR (pp. 59--74) =====
\chapter*{III.\quad Irony and Humor}
\addcontentsline{toc}{chapter}{III. Irony and Humor}
\markboth{Irony and Humor}{Irony and Humor}
\opage{59}15.~Whether one regards humor merely as a peculiar kind of feeling of
laughter or as a type of life, there is no contrast to it so characteristic as
scorn. But there are several intermediate forms between these two extremes, and we
shall consider some of the most important here.

Scorn presupposes, as we have seen, a relation of contrast that is not present in
purely involuntary laughter. Scorn is antipathetically taken up with its object; it
brings with it a need to occupy itself with it again and again in order to bring out
and paint its impotence, its insignificance, its stupidity or its baseness. It is a
decidedly reactive feeling, although Nietzsche thought that it was precisely suited
to high superiority. Well-grounded superiority leads, as stated above, to humor
rather than to scorn; it rests in itself to such a degree, taken up with its rich
self-unfolding, that no need can arise for side glances or downward glances, and
where it stands in relation to other beings, it has at any rate no ground for
antipathy. It is a sign of insecurity and of a lack of absorption in the unfolding
of life when it is necessary again and again to dwell on the paltriness of others.
It was the tragic thing in Nietzsche that his constant need to utter his antipathies
did not give him the rest and strength to develop positively what nevertheless lay
nearest his heart, the great power and self-affirmation of life. Again and again he
breaks off the working out of his main work to let scorn and antipathy have the
word. The work then was never finished either. But the contradiction appears most
glaringly in him when the feeling of distance (das Pathos der Distanz) toward the
great mass is to be an essential element in the disposition of great men. He did not
see the sharp conflict that \opage{60}exists between the pessimism that ``the pathos
of distance'' necessarily presupposes (since ``the great mass'' is only to be able to
be means, hindrances or copies!) and the absolute optimism his gospel was to
proclaim.\footnote{See my characterization of Nietzsche in \textit{Moderne
Filosofer}.}

In great satirists one can study this relation, and in no one better than in
Swift. His absorption in and enthrallment by what he rejected and fought comes out so
glaringly that the whole figure has often been darkened by it for posterity, just as
Swift's own mind was darkened by it. He constantly seeks new expressions for the
contempt he felt for the mass of men and their doings, as he thought he had come to
know it in his participation (or attempt at participation) in political life. He
revels in images and allegories that could express this feeling. He has his Gulliver
travel to Lilliput so as to be able to ridicule the culture of Europeans in this
people of Tom Thumbs. And when Gulliver then comes to the giant people of
Brobdingnag, there is occasion, through the account Gulliver gives there of life in
Europe, to depict Europeans as ``the most pernicious race of little odious vermin
that nature ever suffered to crawl upon the surface of the earth'' (Gulliver's
Travels. Part II. Chap.~6). % Swift's own words; Høffding quotes them in Danish
The high point is reached, however, later, in the
description of the people called the ``Houyhnhnms,'' an ideal people, which of course
does not have human shape but the shape of horses. In this land there lives, besides
the main people, a hideous, bestial, base-minded race that has human shape and lives
a human life. In these ``Yahoos'' disgust at the human has found its most violent
expression.

The great satirist has found in these depictions of lowness and baseness an
expression not only of his moral indignation but also of wounded self-esteem, of
inner sufferings and of tragic circumstances of life. It is, to use Leslie Stephen's
expression in his book on Swift (p.~184), a wounded soul that speaks to us, not
merely a prophet chastising wretchedness.

Perhaps every prophet has a wounded soul. But on posterity he will nevertheless work
most by his ability to make the great and the true prevail. It is, after all, only
the prophetic \opage{61}figures of modern times whose inner life we know. In the old
times there was no sense for individual psychology, and in any case it was only
insofar as the prophet, in indignation and enthusiasm, unfolded the powers of a great
mind that he made an impression on the nearest posterity, from which the picture we
have of his figure stems. When he had then gone over into being one of the great
exemplars of the race, the purely individual features easily fall away, especially
those owed to the conflict between inner, conflicting drives and tendencies. Had
they been wholly taken up with inner strife, they would hardly have been fit to be
an example for others, but would have had enough to do with themselves. It is a
great, often insoluble task to find the psychical connection between the mighty
polemic and what went on in the prophet's soul.

16.~The ironic form can belong among the means of satire. It then comes forward
with apparent acknowledgment and earnest, although its real endeavor is to attack and
annihilate. An indirect form of attack then takes the place of the direct. Irony is
here only a means; it does not designate a particular type of life. It is what one
could call \emph{the little irony}, while \emph{the great irony} is not merely a
form and a means but a stage of life. The difference will consist among other things
in this: while with irony as a form or figure of speech it must be aimed at that the
others understand that it is irony, that the earnest is not earnestly meant, it is
otherwise with the great irony. But before we investigate it, we must dwell somewhat
on the history of the concept of irony.

In Greek usage and Greek thought this concept underwent remarkable changes, which
have been shown in an interesting way by Ribbeck (Über den Begriff des εἴρων.
Rheinisches Museum. N.~F. XXXI p.~386).

The word appears first in Aristophanes' ``Clouds'' as a term of abuse, the ironist
being placed here alongside the liar and the braggart. For the Greeks, who both
presupposed and demanded that the inner should be expressed in the outer, irony, in
which the inner stands in contrast to the outer as jest to earnest, had to be a bad
thing. When Socrates was called an ironist, it was an expression of people's ill
will and exasperation at the way he lured those he talked with onto slippery ground
by apparently acknowledging \opage{62}them and going along with their opinions. It
is, as Ribbeck says, an irony of fate that precisely the most truth-loving of all
Athenians should be classed with liars and braggarts. But precisely the personality
of Socrates, as Plato depicted it, also brought about a change in the conception of
the concept of irony. In Aristotle the ironist is set in contrast to the braggart,
although, like him, he designates a defect in contrast to the truthful man (Eth.\
Nic.\ IV, 13). The ironist's behavior is explained by his reluctance to make a
show. But it is not yet stressed here that the great difference between the ironist
on the one side and the hypocrite and the braggart on the other consists in this:
while the braggart makes the outer appearance his end, because he precisely wants to
make himself greater than he is, and while the hypocrite uses the outer appearance
as a means to attain selfish ends, for the ironist the appearance is a means to keep
his inner life undisturbed, or perhaps even, by going along with the presuppositions
of others, to get them to gain in clarity. It is in this way that we now conceive
Socrates' irony, guided by the indication Plato has given in the depiction of his
master. In a single place Aristotle gives a hint in this direction, when he says of
the man who has a great soul (ὁ μεγαλόψυχος), and who not only thinks himself worth
something but really is worth it, that he behaves ironically toward the multitude.
(Eth.\ Nic.\ IV, 9). % the Danish prints "Eth. Nrc."
He has, after all, something within him that he cannot communicate without more ado.

Later Greeks again stood as strangers to the contrast between the inner and the outer
that showed itself in Socratic irony. The Epicureans conceived irony as an expression
of arrogance and braggartry,\footnote{\textit{Theophrasti Characteres et Philodemi
de vitiis liber decimus.} ed.\ Ussing. p.~4. 57. 175.} and Cicero's teacher, the
Epicurean Zeno, even called Socrates an Attic buffoon (scurra
Atticus).\footnote{Cicero: \textit{De natura deorum}. I, 34.} ---

Socratic irony is great irony, if Socrates' stage is exhausted by the concept of
irony at all, which will be considered more closely later. In the great irony the
polemical relation to the surroundings falls away. There can, as already indicated,
be two tasks solved by this irony. Irony can be the refuge of \opage{63}a mind that
wants to withdraw from its relation to the outer, to other people, to the given
forms of life, guided by the need to assert itself and work at the unfolding of its
inner life, a need that must naturally become the stronger the richer the content
that is working its way forward in the personality. A great inner fullness can make
isolation necessary, and by going along outwardly with and acknowledging the given
forms of life, people's opinions and customs and the ruling authorities, one saves
all one's strength for inner work and inner growth. But irony can serve not only
self-assertion. It can also, and at the same time, be a method, a pedagogical means
to make possible the understanding of other personalities, an understanding that is
in turn the condition for help, education and collaboration. Through apparently
going along with the presuppositions and assumptions of others, through
acknowledging the place and the standpoint on which they stand, one more easily gets
them to open up, to unfold the consequences of their thoughts, and with that the
cooperation in which all real help and all true education consists has already been
made possible. While satire pushes its object away, or holds it fast only to push it
away again and again, irony can be a way by which it becomes possible to penetrate
into the object and set up a community with it. Even where irony serves
self-assertion, it is not a repelling but a withdrawal, a way of keeping the
sanctuary of the mind free from the intrusion of the surrounding world.

When the ironist arouses anger by his conduct, it may be because people do not like
anyone wanting to educate them, to lead them beyond that in which they securely
rested. As soon as they discover that this was the end of the apparent acknowledgment
of their standpoint, their self-esteem is wounded. But the reason may also be that
people do not like to see anyone withdraw from community with them. It causes
insecurity, merely that there is someone who will not join in. One wants a common
confession. Not only in the domain of religion is there dogmatism. If I am right,
how can another then have a right to go other ways than I? And especially when I do
not even know what other ways he goes, I perhaps begin to feel uneasy.

In both its forms, both as self-assertion and as \opage{64}education, the kind of
irony described here, after hints in the historical Socrates\footnote{See more
below § 44.} but in a further working out, marks a great turning point in the life
of the spirit. It is the first form in which the principle of personality, the
understanding of and respect for individual human personality,\footnote{On this
principle in its more special elaboration see my Helsingfors lectures, published
under the title \textit{Personlighetsprincipen i Filosofin.} Stockholm 1911.}
one's own and that of others, breaks through. And Socrates stands here as the
pioneer. In his irony, so far as one can see, both the forms named asserted
themselves.

Kierkegaard in his brilliant youthful work (Om Begrebet Ironi, med særligt Hensyn
til Sokrates. 1841) constructed an absolute irony, by which he meant an irony of a
purely negative kind. It was the little irony he made absolute, instead of making
it great. He regarded irony merely as dissolving, as devoid of every positive
content. He denies it both inner fullness and the sympathetic element. The jest
hidden in ironic earnest is supposed to be directed at everything between heaven
and earth. Everything is dissolved into nothing --- with the exception of the
selfish. The ironist takes everything as vanity --- except himself. (Skrifter XIII
p.~332). Socrates dissolves Greek spiritual life (ib.\ p.~338) without putting
anything whatever in its place. But --- Kierkegaard thinks --- this was justified
in Socrates, for the time of Greek culture was essentially over. In modern times,
on the other hand, such an irony (as it appears, for example, in Fr.~Schlegel) is,
he thinks, altogether unjustified: for now ``subjectivity'' (the principle of
personality) has triumphed. It has carried through its demands in state and church,
in social and religious life, and it constantly reveals itself through these forms.

Kierkegaard proceeds in the youthful work mentioned from the speculative theology of
his time, developed under Hegel's influence, which was represented among us by
Heiberg and Martensen, to whom he also refers. He was thereby led, at this stage of
his development, to misjudge the positive character of Socratic irony. Even later
(1846) he calls the ironist (Papirer VII, 1. p.~27) a buffoon. It recalls the
Epicureans' scurra \opage{65}Atticus. For the rest, he corrected his view in later
writings, as will be shown later.

Not only Socrates but also Fr.~Schlegel Kierkegaard did wrong in his youthful work.
The circumstances are supposed, according to Kierkegaard (and Hegel), to have been
entirely different when Schlegel appeared than in Socrates' time. ``Subjectivity''
had now, after all, triumphed in state and church! --- To this it must be said that
when Schlegel appeared it was hardly unnecessary to maintain the right that the free
life of the spirit has over against the traditions crystallized in institutions and
dogmas. There is something here, after all, that cannot be settled once and for all.
There is an eternal strife between freedom and norm. Externality, triviality and
spiritlessness are precisely often at home in the forms for which the earlier fresh
upsurge of the life of the spirit laid the foundation. One rests content with
celebrating the great, spirit-filled beginnings and thinks oneself, despite all
stagnation, akin to them. There seems, especially in Germany at the beginning of the
nineteenth century, when Schlegel appeared, to have been much of this kind to
combat. Thus Lotze remarks in his ``Geschichte der Aesthetik in Deutschland''
(p.~371): ``the mood then prevailing in the German people could, in aesthetic
respect, rather be called empty than rich. In sluggish tradition and narrow-hearted
custom, receptivity for the beautiful had been lost to such a degree that it could
stand as the task first and foremost to restore the unprejudiced life of impulse by
rising against countless barriers and testing and disputing countless prejudices.''
One of the best men of the time, the great idealist Fichte, even declared in a
letter that the time was characterized by ``absolute putrefaction of all ideas.''
The time when Fichte and Schlegel appeared hardly had any advantage over the time
when Socrates appeared. On the contrary, Socrates' appearance was bound up with one
of the highest cultural periods of world history.

And despite all the objections that can be made to Fr.~Schlegel as a person and as a
writer, there was in him a real need for freedom and honesty as the first condition
of a worthy life. It was this that won him friendships with men like Fichte and
Schleiermacher. And instead of merely relying, with Kierkegaard, on an analysis of
Schlegel's novel ``Lucinde,'' one ought to \opage{66}hear how Schlegel himself gave
an account of his concept of irony in his essays in the periodical
``Athenäum.''\footnote{I have already remarked this in \textit{Den nyere
Filosofis Historie,} II note 38. --- The statements of Schlegel cited in what follows
are found in \textit{Athenäum} (1800). I,2 and III, 1.}

It is, according to Schlegel, the impossibility of complete self-communication that
asserts itself in irony. The most precious thing man possesses, inner satisfaction
itself, has in the end its root in something obscure, which is nevertheless what
bears the whole, and which would lose its force the moment one tried to dissolve it
into pure understanding. In irony there appears the consciousness of an eternal
mobility (Agilität), an infinite chaos that can never be wholly systematized. Only
apparently can the personality be exhausted in a single relation or in a single
work. There is in every personality, as in each of Leibniz's monads, a universe, an
infinite fullness that makes it both possible and necessary to transport oneself,
not only with understanding and fantasy but with one's whole soul, into the most
different domains, although none of them can wholly suffice.

These statements show that Schlegel himself conceived his standpoint as positive, and
that he was right in this. But it also appears probable that such a standpoint could
not be definitive. It was still too chaotic. There was a fullness that pressed
forward, but it stood as arbitrary which part of this fullness one took up in
particular, and the arbitrary play, the capricious shifting, evidently had a great
attraction in itself. And Fr.~Schlegel was unable to give shape to what pressed
forward. He lacked the energy of Socrates and the Socratics to work in the free field
of personality, where all good plants grow. His sense for the great and the infinite
was paired with a lack of ability for steady and faithful work and for becoming,
through it, master of the chaotic content. Therefore he ended in
dilettantism.\footnote{Cf.\ \textit{Dilthey's} fine characterization of Fr.~Schlegel
in his \textit{Leben Schleiermachers.} Second book. Ch.~4. 12. 14.} % the Danish prints „Schleiermaches“

There came in him a reaction against the capricious shifting, and a need stirred for
definite forms of the life of the spirit. He was thereby led to Catholicism. Already
in the ``Athenäum'' \opage{67}he longed for the Orient. ``Zunächst rede ich nur mit
denen, die schon nach dem Orient sehen'' (III, 1). It was from the first ``the
eternal Orient'' he sought; but then he thought he found it in the mighty
crystallization of the old Church.

17.~Shakespeare too some have wanted to make an ironist. A.~W. Schlegel
(Fr.~Schlegel's brother) constructed a Shakespearean irony, which was supposed to
show itself partly in the characterization of the single persons in his dramas,
namely in the way great and small motives clash in them, partly in the action of the
drama as a whole, through the impartiality with which the course of fate is
depicted. Characters are depicted in their whole nakedness, not in the garb of
enthusiastic admiration, and fate is exposed in its whole reality without the poet's
heart being moved by it.\footnote{\textit{Vorlesungen über dramatische Kunst und
Literatur}\textsuperscript{4}. II. p.~197--199.}

Against this one of the finest recent Shakespeare scholars\footnote{\textit{Dowden:
Shakespeare. His Mind and Art.}\textsuperscript{5} (1880). p.~345.} rightly objects
that Shakespeare's impartiality toward persons and fates does not mean indifference
and remoteness, but springs precisely from his interest in his subjects, from the
firm will to let each of their sides come into its own; his enthusiasm is not
weakened by the barriers to which human actions and fates are subject.

There is an artistic irony that is closely connected with the great task of art,
that of giving concrete and individual pictures of characters and fates, not
abstractions and utopias. Everything in art must be subject to the law that springs
from this task. Life is so rich, and works its way forward both within the single
personality and within the race from so many sides, that a limitation, a polishing,
must take place if everything is to come into its own. Both as regards the individual
characters and as regards the events in which they are placed, a negative element
necessarily asserts itself, in that the single trait or the single event is
subordinated to the whole connection and is not allowed to unfold in the way it could
if it had been the only one. \opage{68}In addition, the artist's task is to give the
feelings and images that rule him such an expression and such a form that his work
can acquire worth and win recognition also beyond the circle that most nearly
surrounds him and perhaps fully shares his feelings and is familiar with his images.
Therefore what moves within the poet must be shaped into clear and definite forms,
have the marked character of objectivity.\footnote{Interesting remarks on this point
in \textit{Yrjo Hirn: Konstens Ursprung.} Stockholm 1902 p.~95--102.}

But it would not be right to lay the main weight on the negative, limited element in
artistic shaping. The essential thing is, on the one side, the glowing mass that is
to be formed, and, on the other, the plastic figures that come forth through the
shaping. The struggle, the often bitter struggle, and the work, the often hard work,
lie between these two positive elements and are probably most often known only to the
single artist himself. The artistic ability is the ability to individualize what
presses forward involuntarily in feeling and fantasy. In art as everywhere the
individualized and fully marked proceeds from a more indefinite basis, an unorganized
fullness. The difference between the artist and the rest of us consists essentially
in the ability to transpose what is obscurely imagined and felt into individual
figures and events. And this ability means something positive, not merely a
limitation. When the work is finished, the casting successful, there is no contrast
between the involuntary fullness and the shaped figures; the former has come fully
into its own in the latter.

18.~The kind of irony that has the character of a positive going along with the
presuppositions of others stands close to humor. There is a subterranean connection
between the humorist and his object, as there is in pedagogical irony. The difference
between irony and humor has often been expressed thus: in irony there is jest behind
earnest, in humor earnest behind jest. Earnest is everywhere conditioned by the
relation to a value: earnest sets an end and sets the will in motion in the direction
of this end. In irony one apparently \opage{69}goes along with values and endeavors,
but the ironist makes this movement only to keep his inner self free, or --- in
pedagogical irony --- to release psychical energy in the one who acknowledges certain
values and ends. That there is jest behind the earnest need only mean that the
ironist has his earnest in other directions than those around him. Therefore irony
can be an outer form for humor, so that we get the following scheme (where the sign
$<$ means ``means'' or ``expression of''):

(Earnest $<$ Jest) $=$ Irony $<$ Humor $=$ (Jest $<$ Earnest), or more simply

\begin{center}Earnest $<$ Jest $<$ Earnest.\end{center}

The jest in which the humorist hides his earnest consists, then, precisely in the use
of irony. Or more simply: behind the jest that in irony is hidden in earnest there
lies earnest hidden again.\footnote{In a corresponding way the little humor can
enter the service of irony, so that we get the scheme: Jest $<$ Earnest $<$ Jest. The
earnest that in irony is the outer form for jest is then expressed humorously, i.e.,
clothed in jest. Formally, psychical processes like those indicated here could
always be carried further: Earnest $<$ Jest $<$ Earnest $<$ Jest \dots{} But where a
definite type of life, a real self, has developed, there will in the end be values
and ends that cannot be taken in jest by the individual himself.} In this way the
great irony, as a form of self-assertion or of educating care, passes over into
humor.

One might ask why earnest should be expressed in jest at all. Fr.~Schlegel (of whom
Novalis rightly said that he was really a humorist, not an ironist) has already given
the answer. The answer is that there can be so intimate and so comprehensive a
spiritual content of life that every positive form becomes unsatisfying,
incommensurable. And this will be the case especially where experience of life has
set us face to face with great contrasts whose sting we have felt. The value in whose
service earnest stands will prove to stand in contrast to other values and thereby
points beyond itself. Here, as in artistic shaping, a negative element comes forward.
Even the highest values must be shaped and limited through their relation to other
values. This experience makes the \opage{70}humorist prefer to speak jestingly and
indirectly; he thereby guards his inner tasks, reserves for himself the continued work
on them. Then there is in addition the relation of value to reality, and with it a new
series of tasks. Reality can, as experience shows, conflict with our values, or be
altogether indifferent toward them. We need not therefore give up the faith in the
validity of values that shows itself in continued work to realize them; but we become
more reserved in confessing, and the jesting or indirect form then offers itself
naturally. By using it we do not shut ourselves up absolutely; the experienced and
knowing will be able to read earnest out of the text of the jest. The jest itself,
however, is not merely a refuge, a way out. Through it the mind proves its
superiority over the contradictions and contrasts it has to struggle with in its
intimate holding fast to what makes life worth living. This feature comes out in all
the different forms of humor (to which we shall return in the next chapter).

Of greatest significance here is the contrast between the world of values and the
world of reality.

If one takes a purely empirical standpoint, the world of reality stands as far more
comprehensive than the world of values. The valuable is only a part, perhaps a
vanishingly small part, of the world of reality. This came forward in grand style
when the old picture of the world, with its cozy surveyability within a closing frame,
was dissolved, and it dawned on consciousness that it is impossible to fix limits for
the world. Angels could now no longer ascend and descend from the heaven of values to
the earth of reality, and no absolute center can any longer be shown, either as
regards value or as regards reality. Like a little oasis the realm of values now
stands in contrast to the great desert of reality. Or let me use another image, taken
from an entry in a diary.

``It is cold and raw today. A sharp wind sweeps through the leafless trees. But this
same wind produces in the Aeolian harp on the house I live in tones full of mood, a
long, deep, soothing sound. So the realm of values can reveal itself in the midst of
the raw, cold world of reality.''

\opage{71}There are two elements here, which both images indicate. The oasis lies
within the desert itself; the same wind that causes the dreariness also calls forth
the comforting sound. Value, then, is both a contrast to reality and a part of it.
Through this doubleness a mood of life like humor becomes possible --- a disposition
that does not lose its sense for the valuable and its faith in its validity because
not everything can be derived and explained from it, or because everything else in
the world seems indifferent toward it. Humor will fasten on this: that it is by
virtue of the same laws that something valuable exists in the world and that the
indifferent or hostile reality comes to be. The wind shakes the trees, but sounds
through the Aeolian harp. It stands as a victory that in some places in the world
there has nevertheless been a breakthrough of something great and good, even if it is
a pure fantasy that everything should be arranged so that such breakthroughs could be
possible. A humorous mood can come forth before the representation of the
overwhelming power of the enormous masses, which nevertheless cannot prevent the
sound of values from being heard. For it is only the materially or quantitatively
great that appears in the masses of the world; in the world of thought, and of values
in general, there expresses itself a greatness of another kind, a greatness of
inwardness, with which all extension in space and time is incommensurable. A feeling
of defiance can arise when the contrasts come forward in their sharpest form; but the
more the mind immerses itself in the values themselves and in the causal series within
the world of reality that produce values, the more the great humor will come to rule.
And it is, after all, also such a disposition that
is the basis of energetic work to produce and hold fast values.

A special occasion for humor toward the course of things in the world arises from
the experience that the valuable is often produced by causes that seemed to work in
an entirely different direction, so that they, to speak mythologically, against their
will or without their will come to accomplish what stands for us as beautiful and
good. It is the same thing that, when human feelings and actions are in question, is
called the shifting of value. (Cf.\ Psychology VI C; Ethics XIII, 4). Natural forms
and natural scenes that \opage{72}appeal to our aesthetic feeling have nevertheless
not developed in order to produce aesthetic values, but according to definite causal
laws. Human forms of society and culture to which we ascribe ethical value may
nevertheless have developed under the influence of the blind need for
self-preservation, nourishment and propagation. It is often manifold and very
diverse motives that bring about the great and valuable results in history. This
holds in small things, in people's good advice and friendly helping
hand,\footnote{Cf., for example, the complex of motives that (in \textit{George
Eliot's The Mill on the Floss.} Book I. Chap.~3) led Mr.~Riley to advise
Mr.~Tulliver to put his son in a particular boarding school. It is a piece of
psychology that cannot be summarized.} as it holds in decisions of world history.
What came into the world as a mere means, perhaps an enforced means, later acquired
value for itself and could be set up as an end. Thus, for example, the great
conquerors have often worked in the service of humanity by leading beyond national
and confessional boundaries, although their own end was the satisfaction of lust
for power and ambition. And even where it was an end valuable in itself that was
set, it often proved to be only a means, a station on the way to a greater goal.

One may grieve that to the great ends there do not always correspond great and noble
means, but that it is often very elementary or even purely egoistic causes that lead
to the realization of the great values. But one can also see it as significant that
even such causes come to serve the valuable in the world. It is then not the values
that are degraded by being realized through elementary causes, but conversely the
causes that are ennobled by the laws according to which they work containing the
possibility of the arising of valuable wholes. There is room, then, for smiles as
well as for tears, and both moods can unite into humor, with different timbre
according to the relation between the two single feelings.

A little diary entry can be used here too, partly as example, partly as symbol. It
arose from an observation made on the steam ferry across the Great Belt. --- ``I
stood looking at the gulls that followed the ship, and admired their quick and
\opage{73}daring turns, the perfectly formed movements, the graceful curves their
flight described. I enjoyed the splendid sight of the white-gray birds in the
radiant sunshine. And then I came to think of the cause of all this glory --- what
it was that made the whole sight possible: the need for nourishment, the striving
for food, the struggle for sustenance!'' --- So far as I remember, the sight did not
lose its luster through this reflection; but the mood became somewhat more composite.
Perhaps the first joy was nevertheless pushed aside for a moment, but then returned
strengthened.

With all ideal values shifts of value play an important part. As regards art, Yrjo
Hirn has shown in a very interesting presentation how a whole series of practical
(``utilitarian'') factors furthered the development of the different art forms and
led to the growth of the sense of beauty. Such factors are: the need to exert magical
influence or erotic enchantment, to spur on to war or work, and not least to
communicate one's thoughts to others.\footnote{\textit{Konstens Ursprung}
p.~173--294.} Something corresponding holds, as I have shown in my Ethics, for
ethical ideas and feelings. ---

The great humor sees life both as great and as small, both as favorable to value and
as hostile to value, both as tragic and as comic. Solger was the first to describe
humor in detail as the unity of comic and tragic feeling, of irony and enthusiasm. He
is somewhat wavering in his use of the concepts irony and humor. Most clearly he
describes humor in his dialogue ``Erwin'' thus: ``The tragic and the comic here still
lie undivided side by side\dots{} Everything in humor is in one stream, and
everywhere, as in the ordinary world, the contrasts pass over into one another.
Nothing in humor is ludicrous and comic that does not have an element of dignity or
cannot awaken wistfulness --- nothing is exalted and tragic that could not, in its
shaping in the low world of time, sink down into the insignificant or the
ludicrous.'' (Erwin. Berlin 1907. p.~354). Only in longing can the thought of the
highest, as raised above vanity, dwell; but the humorist must regard this
\opage{74}longing as worth far more than the limited form in which fantasy is able to
shape the highest (Erwin. p.~358).

This Solgerian standpoint Kierkegaard does not let come into its own in his treatise
on ``Begrebet Ironi.'' He looked at it with theological-Hegelian eyes. In the work
mentioned (and in notes contemporary with it) Kierkegaard conceived the relation
between irony and humor as one of opposites. Irony is supposed to be egoistic, and to
be placed in constant struggle against the world. Humor, on the other hand, is
supposed to be freed from the world; it stands as the expression of a religious
earnestness of life and of the consciousness of belonging within a great whole. Irony
is supposed to be aristocratic and polemical, but humor is reconciled with the world
and rests in a faith that is held fast despite all absurdities. (Skrifter XIII
p.~35--119.)

Later, on the other hand, after the definitive break with speculative theology,
Kierkegaard no longer regards humor as a religious stage. It now stands for him as the
last human stage before the religious stage --- as the expression of the limit of the
highest purely human. It is characterized more closely as essentially conditioned by
wistfully smiling recollection. It is directed backward, or rather, time no longer
plays any part for it. Religion, on the other hand, looks forward. At the religious
stage life is, from first to last, in small things and great, determined by the
representation of God, in relation to which everything else loses its worth, so that
the highest becomes self-annihilation before God.\footnote{The writings belonging
here are: \textit{Studier paa Livets Vej} and \textit{Uvidenskabelig Efterskrift.}
--- Cf.\ on Kierkegaard's stage of humor \textit{Danske Filosofer}.
p.~160--162.} --- A test of this definitive view of humor of Kierkegaard's will be
undertaken below, in the chapter on the tragic and humor.

% ===== IV. MAIN FORMS OF HUMOR (pp. 75--87) =====
\chapter*{IV.\quad Main Forms of Humor}
\addcontentsline{toc}{chapter}{IV. Main Forms of Humor}
\markboth{Main Forms of Humor}{Main Forms of Humor}
\opage{75}19.~The single feelings that later, in union, produce a total feeling can
come from very different sides. Here especially, as on so many other points, the
relation between self-assertion and devotion is of significance. Humor thus takes on a
different timbre when it develops on the basis of superiority and loftiness of mind,
or on the basis of wistfulness and longing, than when its essential elements are
sympathy and understanding. Moreover, within these different forms it again depends on
what has been experienced of sorrow and joy, of want and enjoyment --- what experience
and reflection have opened the eye to --- how great a wonder life has called forth ---
what disappointments it has caused --- and what understanding has been won from the
whole. Here, as on so many other points, the sporadic character of the mental life at
its earlier stages of development, so strongly stressed by Sibbern, appears. Now when
both the basis of the experiences and the experiences themselves are different, there
must naturally occur an extraordinary number of forms of humor. There will moreover be
a great difference in how easily and quickly the fusion or the organization takes
place. Often there will be hard crises to go through, inner strife to fight out,
before it reaches a certain firmness and closedness in one of the main forms of humor.
Such periods of fermentation belong among the psychological experiences Sibbern has
described best.\footnote{\textit{Danske Filosofer} p.~101. --- \textit{Sibbern}
first arrived at his idea of sporadic development by the psychological path. Later he
found analogous relations in nature and in the circumstances of society, so that that
idea could enter as an essential member into his view of the world. Since he very
often uses development within embryonic life as an example, I assume that he was
influenced by the French embryologist \textit{Serres,} who (1829), in contrast to the
development from center to periphery assumed by Harvey, Wolf and Haller, maintained
that embryonic development proceeds ``par des noyaux latéraux qui se soudent
ensuite,'' something he showed in particular for the formation of bone.
\textit{Cuvier: Histoire des Progrès des Sciences Naturelles.} Bruxelles 1837--38. II
p.~431--434.}

\opage{76}However prominent the sporadic character of the mental life is, it is
still a more prominent peculiarity of it that every psychical element has a tendency
to enter into connection with other elements so that a whole can be formed. It is
Kant's immortal merit --- greater than the one he himself laid most weight on --- that
he maintained the need and ability for gathering activity as the innermost
peculiarity of all mental life. This need and ability he found in sense perception and
in the association of representations, and especially in the special forms that
scientific research presupposes. He also found in it the explanation of the human
spirit's constant yearning to reach out beyond the domain of experience and science,
since it could not, after all, from the first know its own limits. With this way of
looking at things Kant set the science of the mind its proper task.\footnote{\textit{Den
nyere Filosofis Historie.} II. p.~27--29; 42--46. --- Cf.\ above § 2 and § 8.}

What is to express the peculiarity of the mental life in all domains must naturally
be of a formal kind. What Kant called ``the hidden art in man's interior'' is not
bound to any definite content, precisely because it sets its stamp on all content
taken up into the soul. Kant himself applied his concept of synthesis predominantly to
the domain of the life of thought. But if he is right in his view, something
corresponding must also hold for the other sides of the mental life. There must thus
be a certain tendency, a need and an ability, to gather the joys and sorrows of life
into a total feeling, which one might perhaps call a synthetic feeling, just as in the
domain of the will there must be a tendency to gather all endeavors into one great
striving.

What are called, often somewhat pretentiously, psychological laws state the different
forms and conditions of the gathering activity. As regards the special formations of
wholes, \opage{77}one cannot, as already remarked earlier, expect that they could be
derived from the knowledge we have of the sporadic utterances of sensation,
representation, pleasure and pain, involuntary striving, drive and wish. There are so
many possibilities of cooperation and counteraction, of strife and union, of
connections of very different degree and kind, that we are referred to studying the
wholes after they have actually formed, and then afterwards finding how the sporadic
development may have led to such results. We are placed here, in the psychological
domain, in a way similar to our position before all formations of wholes in
nature.\footnote{Cf.\ on this: \textit{Den menneskelige Tanke} p.~218--237;
338--345. In a special paper I hope to illuminate this point more closely from an
epistemological viewpoint.}

The great humor, as I am trying here to describe it, has a quite peculiar possibility
of being a concluding total feeling. We still lack thorough psychological
investigations of the different possible total feelings that might tend to take on a
concluding character for a period of life, or perhaps for a whole life. The
description and analysis of a single total feeling cannot furnish proof that it should
be the most significant or the noblest. I shall only draw attention to what I called
the great humor and investigate the conditions it has for being a concluding total
feeling. It comes to be through the experience of the contrasts of life, through the
clash of single feelings, with the cooperation of reflection and of striving for ends
held fast with energy and faithfulness through the years. And it has the possibility
of keeping receptivity open to what new experiences may bring.

The psychological possibility of such a total feeling, the dangers that threaten it,
and the conditions to which it is subject will be investigated in the following
sections of this treatise. For the present we go more closely into some of the main
forms already mentioned (see § 13) in which the great humor can appear.

20.~The limiting phenomena appear here in such cases where it is superiority or
loftiness of mind that forms the basis. \opage{78}Against such a basis even very strong
and penetrating experiences can glance off without lasting effect. They do not shake
the firm center of the mental life and therefore release neither scorn nor fear. But
precisely because of the inner firmness there is a possibility of a jesting attitude
toward experiences, the unfavorable as well as the favorable, as well as of sympathy
with people. The more freely a person stands in his own self-assertion, the more he
masters it, the more he can take part in the struggle of others for self-assertion.
There is a secure feeling of power present here that brings no need to make
superiority felt toward the small, but can precisely give rise to the need to find an
outlet for energy through helping work. True superiority does not need to compare
itself with those who stand lower, but presupposes precisely that one has won a
thorough self-assessment through comparison with the greatest. Often the way goes
through struggle and conflict. As long as no resistance is felt, no hindrance stops,
no expectation is disappointed, no thorough total feeling will form. The strength is
not put to a decisive test. As with all experience, it holds here too that nothing
but positive instances permits no secure inference. Just as logic demands that we
inform ourselves not only about the cases in which a phenomenon occurs but also about
those in which it does not occur, so too experience of life is won only when one works
not only with the current but also against the current. Only thereby is the ground
color obtained that according to Plato was necessary for the proper color of character
to become fast. Even where the development of personality proceeds continuously,
reliability and firmness come into it only when it has been tested through
disappointment and want, and when it has learned to hold a great goal fast despite the
hindrances it meets and despite the distance at which attainment lies. A storm, says
La Rochefoucauld, makes a bonfire blaze more strongly, while it puts out a candle. If a
feeling can maintain itself although its object is far away and no sensuous presence
or vivid image supports it, it will become all the more intimate and strike root all
the more deeply. A great sorrow can release energy that would otherwise not come into
play. \opage{79}It already takes energy to grieve, to endure and live through the
tension between the feeling of the worth of what is lost and the certainty that it is
now really lost. The easiest way to be freed from this tension would be to push one of
the terms out, either to forget what is lost or to imagine that what is lost has no
worth. It takes superiority and loftiness of mind not to seize either of these ways
out. Or rather it is the other way round: through the courage that seizes neither of
these ways out, the way goes forward to superiority and loftiness of mind.

Stoic loftiness of mind did not go this way to the end.\footnote{Cf.\ on Stoicism
\textit{Mindre Arbejder.} First series (the paper on pagan seekers after truth) and
Third series (the papers on Epictetus and Marcus Aurelius).} It was dearly bought with
fear of the involuntary movements of the mind. Every bond that ties us to someone or
something brings the possibility of being drawn into joys and sorrows that come
independently of our will. Therefore, the Stoics maintained, if one wants to keep
one's spiritual freedom and balance, one must keep to what is in our power and regard
nothing as good or evil except insofar as it furthers or checks our self-mastery. The
Stoic stands at the play of life as a cool spectator who attentively follows the
course of the play but takes good care not to let himself be carried away by it. Above
all he fears coming ``to sigh within himself.'' The inner sanctuary is to be kept
wholly untouched by what happens in the world.

For Stoicism, therefore, more than for any other direction in the Greek conception of
life, humor is a foreign feeling. Stoicism seeks precisely to push out and exclude the
elements by whose connection humor could arise. It shuts the contrasts of life out
instead of working its way through them. Humor is possible only when self-assertion
need not be so anxious and convulsive that it cuts itself off from being gripped by
what is experienced in the course of life. Humor presupposes that so much energy can be
released and put at the disposal of self-assertion that one can take the blows the
world gives, and be a witness to the stupidity and sluggishness that spread in the
world, without being shaken in one's inner confidence and firmness --- indeed, even so
that one has a smile to spare for it.

\opage{80}When one has learned that the highest happiness consists in striving and
work, one will find the most natural explanation of all moral evil in its springing
from inertia, from a lack of ability to develop beyond the stage on which one has
come to stand or rest. What is lacking is a sufficiently strong influence and impulse
to release the bound forces. From such a viewpoint the struggle between good and evil
presents itself differently than when one, with the theologians, derives evil from an
incomprehensible fall that was supposed at one stroke to have brought the highest
sin, defiance of the clearly known good and right, into the world. And when one then
also has the experience that precisely the struggle against evil can release ideal
forces to combat it, call forth great and beautiful endeavors that would otherwise not
have arisen, the feeling toward life and existence becomes more composite. It may
recall the mood that in an old Easter hymn leads to the outcry O, felix culpa! No
formula and no dogma suffices here, here where the greatest contrasts, dependent on
each other, are united in what one might call a cosmic feeling of life. Nowhere does it
show itself more than here that our deepest and most comprehensive experiences of life
can be expressed only in the symbolic form of poetry.

21.~Where there is superiority and loftiness of mind, there is a firm ground, and
the question becomes how great a mobility and receptivity is possible despite this
firmness. Where humor has developed with wistfulness and longing as its basis, the
question becomes the opposite, namely whether despite the mobility and receptivity
there can be the firmness that no total feeling can do without.

Wistfulness looks back, lives in memory. Longing looks forward, lives in hope and in
yearning. There is an element of sorrow in both, in the one because a value has gone
out of life, in the other because the distance from the goal is so great. But there is
joy in both, because one lives and breathes in something that has value, whether it
belongs to the past or to the future. It has rightly been maintained\footnote{%
\textit{Alexander Shand:} Foundations of Character. p.~395.} that the real, hostile
opposite of joy \opage{81}is not sorrow but loathing and disgust. Joy and sorrow can
very well enter into a fusion and an organization. Wistfulness and longing too can go
together; one can long back to the past when it becomes really alive in memory, and
there can be wistfulness in longing when the thought of how distant and difficult the
goal is predominates. The presupposition of both these feelings is that there is
sufficient psychical energy to live not only in what is sensuously present but in a
certain independence of it. In dull and depressed times there can stir, as a last
flaring up of psychical energy, a longing, perhaps a wistful longing, for the time when
one could still long.\footnote{Cf.\ the chapter on ``Sehnsucht'' in
\textit{Holzapfel's Panideal} (Leipzig 1901).} These feelings lose this significance
as an expression of psychical energy where the relation to the present falls away
entirely, and past and future stand in no connection whatever with the present. Life
then becomes a life of dream or fantasy. Romantic humor took on this character in a
high degree. According to Steffens the best thing in the time he lived in was the sense
for the old, vanished times.\footnote{See \textit{Danske Filosofer} p.~43.} Every
view of life that has its ideal in the past has a great danger to struggle with here.

It was by awakening a great and living expectation, an intimate longing, that original
Christianity worked so strongly in its first period. Only such moods as can directly or
indirectly be derived from this great longing bear the name of Christian with
historical right. This longing was directed toward a great goal; therefore it showed
itself so violently and so intimately. But it was not directed toward any distant
goal. In the very nearest future, which the people of original Christianity could hope
to live to see, a great breakthrough of heavenly forces was expected, which was to
bring about an entirely new, ideal order of things (``the kingdom of heaven'' in its
definite, historical meaning). Longing here became a force of world history. And there
is perhaps no better proof of the significance of longing than this: that even after
that expectation had paled and its goal had been pushed out into the blue, into the
indefinite --- which happened when the Church had settled down comfortably in the
present --- even then the moods \opage{82}it called forth, and the representations to
which it gave occasion, work as essential elements in the entirely new religion that
modern Christianity in fact is.\footnote{Cf.\ the section ``Urkristendom og moderne
Kristendom'' in my \textit{Religionsfilosofi} (IV.~C).} ---

Where memory or expectation of great things rules, much of what could otherwise have
stood as important, far-reaching and fateful will now call forth only a smile. What
moves the half-experienced so that balance is lost fades now before the mind's
absorption in the great. And when it nevertheless tries to exert its power over the
mind, the smile is drawn forth --- not the bitter smile of satire or scorn, or the
hidden smile of irony, but the quiet and bright smile of humor. There is, after all,
access to a world far greater than the one the new experiences belong to. Yet that
world is not jealous of this one. It has room open for it, though perhaps not the
prominent place the new experiences would in and of themselves claim.

Where the relation to the present is wholly lacking, or is only negative, wistfulness
and longing do not give rise to humor. Thus it has probably been rightly remarked of
Shelley that he lacked humor.\footnote{Browning: \textit{Essay on Shelley.} London
1903. p.~60--64. --- \textit{A.~C. Bradley: Shelley's View of Poetry.} (Oxford
Lectures on Poetry. London. 1909. p.~163--167.)} His overflowing life of mood found
vent in mighty or ethereal fantasies that expressed his indignation or his longing, and
that bear the stamp of pure light or pure darkness, but do not let the light be broken
in any prism. The whole real world stood for him, in its finitude and limitation, as a
misunderstanding, an illusion.

The basis for humor is often found, rather than in people of fantasy, in people whose
life has passed in practical occupations and amid manifold changes that have brought
experience and understanding with them. Deeds and fates need not be in the grand style
to open the mind to the secrets of life. George Eliot speaks in a letter of ``the
poetry and pathos, the tragedy and comedy, that may be hidden in a human \opage{83}soul
that looks out through dull gray eyes and speaks in a voice of quite ordinary tone.''
One may meet people of this type on one's way, but uses the opportunity all too seldom
to get to know them more closely. One of the minor characters in Mrs.\ Gyllembourg's
``To Tidsaldre'' (``Two Ages'') is a fine example. It is said here of the skipper's
widow Madam Lyng: ``She had been through much in life and was one of those never
sufficiently appreciated people in whom goodness and gentleness of heart take the place
of upbringing and cultivation, and whom the memory of their own sufferings and wrongs
teaches fineness in social intercourse in spite of any outward position in life.'' ---
This characterization contains essential conditions for humor, even if this quality is
not directly ascribed to the good woman. It is especially the ``fineness'' that must be
noticed. It is the ability to take a considerate and understanding attitude toward what
life brings, without the outward differences it presents being able to impress. Much
can here be taken with a smile that in the untried would awaken unrest, pain or scorn.
Sorrow has had, so to speak, a vaccinating force. Thereby the plain woman could be so
much to the story's unhappy heroine and take her fate in her hand. If she was not in
reality a humorist, it was because her disposition first and foremost deposited itself
in action (see ch.~VII), not branching out into moods. ---

Wistfulness and longing can lead not only to outer events and the conduct of others
being taken with humor; our own faults too we can take in the same way, especially when
we can see them as the faults of our virtues, as members or parts of our own
personality. When others then attack, we may perhaps (like Viggo Stuckenberg in
``Bekendelse,'' toward a friend's attempts at conversion) maintain the positive worth
of our sins, because without them we would not be what we are. We need not therefore
underestimate the purgatory of remorse. Our personality as it now is, with its
wistfulness and its longing, has presumably gone precisely through this purgatory; but
in memory the single action fades into a member of the whole development of the
personality, perhaps even into a felix culpa, and in the longing forward there is not
\opage{84}time to dwell on what no longer has practical significance.

22.~The types of humor depicted in what precedes stand nearer to self-assertion than
to devotion. Superiority and loftiness of mind, wistfulness and longing can
nevertheless form the basis for a jesting with one's own and others' conduct and fate
that is bound up with a positive valuing of their acting and suffering. In contrast to
the course of development that leads from the original need for self-preservation in
the direction of egoism (Psychology VI C, 1), these types are the results of a
development that has excluded captivity in the own isolated self and its interests. And
it is by virtue of great energy that this liberation has been reached in the school of
experience of life. There can then be formed one total feeling toward the contrasts and
changes of life that is not conditioned merely by these having intervened in the
person's own isolated existence.

This liberation from the contrast between one's own and another's fate becomes still
more complete where humor develops on the basis of sympathy (in the sense of immediate
joy with and sorrow with). Here there is from the very first a connection between the
person and what sets his feeling in motion. Set against the feeling of unity % Danish: „I Modsætning til“ (so printed)
or wholeness that here once and for all joins subject and object, much
that would otherwise stand as an object of scorn and ill will can be taken in a jesting
way, because it is not for a moment forgotten that behind it there lies a value that
perhaps even shows itself indirectly in the ludicrous itself. It is here too the deep
basis that makes humor possible. The contrast to scorn is here the greatest possible,
and yet the possibility of jest does not vanish; it is precisely this that makes humor
a total feeling.

As I have shown in my Ethics (IX, 1), it would be wrong to set sympathy without more
ado in contrast to self-assertion. Devotion can be a continuation of self-assertion. It
need not be a purely passive state, but can be the expression of a force able to
embrace something other and more than what self-preservation requires. Here the great
saying about love can hold, that it bears all things, endures all things, overcomes
\opage{85}all things. Loftiness of mind need not bear the stamp of love, but in love,
as it can be, there is loftiness of mind. It is love of this kind that can be the basis
of the great humor. It may have developed by different paths: through a shifting of
motive from self-assertion --- through development on the basis of a natural instinct
--- or in the school of sorrow and pain.\footnote{In \textit{King Lear} IV, 6, 225
Gloster says to Edgar: Now, good sir, what are you? --- And Edgar answers: A most poor
man, made tame to fortune's blows; who, by the art of known and feeling sorrows, am
pregnant to good pity.}

The active, energetic character of the sympathy presupposed here will also show itself
in this: that the larger circle of values and interests to which the person attaches
himself in sympathy opens the possibility of sorrows and cares he would otherwise have
been free of. The person could have got through life more easily if he had not widened
the circle of what could set his feeling in motion.

One must think here not only of the relation to other people but also of the relation
to life or existence as a whole. There can be a fellow feeling that extends to nature,
to all life that unfolds and strives forward. A person can attach his interest to
something that reaches far beyond the domain of individual life, and whose advance or
decline determines the worth he ascribes to reality as it is. The humor that forms on
this basis becomes a cosmic feeling of life, a fundamental mood raised above the
contrast between optimism and pessimism, since it does not close its eyes to
disharmonies and misfortunes, but does not, over the suffering in the world and the
downfall of so much that is great and beautiful, lose the feeling of unity with the
stream of value that constantly makes its way despite all hindrances. It is properly
only at this point that humor becomes a view of life.

23.~This last consideration leads over to a presupposition that is in a higher or
lower degree common to all the forms of humor depicted in what precedes. In the
description of \opage{86}these types weight has been laid especially on what I call
the \emph{basis}, i.e., the life of feeling that forms the presupposition. But the
intellectual presuppositions cannot be excluded from the description, whether these
presuppositions are owed merely to the person's personal experiences or also to the
scientific insight into the course and laws of nature and history that he has been
able to win. When these intellectual presuppositions come forward with special weight,
they form a \emph{background} that works together with the basis to determine the
timbre of the whole type. The background is intellectual and objective, the basis
emotional and subjective. The background exerts a more indirect influence than the
basis, just as air and light gain influence on life only when it already exists in
germ. On the same objective background the most different standpoints or views of life
can unfold, according to the different basis that is present. Background and basis will
not be altogether independent of each other. It may depend on the basis how firm and
comprehensive a background develops, and the intellectual background can on the other
hand check or further the development of the basis.

Whatever it may be that we can feel ourselves bound to, were it the greatest and most
comprehensive that could be thought, it is nevertheless part of the great swell of
existence. However far we go back into our innermost, and however high we rise in our
thoughts, we still do not get beyond the great, lawful connection that proves to rule in
all the regions of existence. When this insight has once been won, it will afterwards
color all our experiences. The desire and ability for bold striving forward are not
necessarily checked by it. Experience teaches us that without great ends and great hope
nothing great is accomplished. Without them life stands still or sinks into passive
vegetating. But we also know that even the greatest ends are only members of the series
of values, which can just as little be thought of as concluded as the series of causes.
The goal that is last for us, which forms the horizon of value at which our sight stops,
is perhaps nevertheless the beginning of a new series of values. But it is in
limitation that mastery shows itself, and even if the trees cannot grow up into heaven,
their growth can nevertheless be straight and beautiful.

\opage{87}Barrier and resistance meet not only the great striving; they appear also
where strength barely suffices to keep one's head up. Suffering finds its way also to
those who lead the quiet life and do not challenge fate. As it rains on the just and the
unjust, so pain strikes both the quiet and the striving. But just as the great humor
does not lose its smile because earnest striving must clash with barriers, so its smile
is not lost either because suffering also intrudes on the path of quiet happiness. The
disharmonies of life fade in the end before the great force with which life is carried
on. The humorist does not flee to the desert because not all dreams of blossoms ripen.

The consequence of the reckoning with life that the great humor presupposes can be a
new youth. Humor does not come naturally to first youth. Youth lets its strength stream
out without thinking of the limits, which can, after all, be known at first hand only
through one's own experience. It is impatient and indignant at every stop and every
resistance. The new youth presupposes that one has formed firm plans for the use of
one's strength --- that one has made clear to oneself that one lives in a finite world
--- but that one nevertheless, with cheerful courage --- perhaps with what Sibbern
called a cheerful irony of life --- always takes up life anew. For the fact too that
life constantly offers repetitions one finds a place in one's view of life. The same
energy that, despite all barriers, holds fast the values of life one has won is also
able to
preserve freshness and inwardness despite the dulling and leveling tendency of
repetition.

% ===== V. THE TRAGIC AND HUMOR (pp. 88--111) =====
\chapter*{V.\quad The Tragic and Humor}
\addcontentsline{toc}{chapter}{V. The Tragic and Humor}
\markboth{The Tragic and Humor}{The Tragic and Humor}
\opage{88}24.~As already remarked above, no total feeling can be proved to mark a
complete conclusion of the life of feeling. As long as new experiences come, there
will also be a possibility of change in the fundamental mood in which what has been
experienced has deposited its effects. A complete formation of a whole, whether as
fusion or as organization, is an ideal. Complete conclusion can just as little be
reached in the domain of the life of feeling as in that of the life of thought. We
stand in the midst of a development within which there is constantly the possibility
of unexpected experiences.

The great humor can therefore not be proclaimed as a mood of life concluding once and
for all. Yet it has many conditions for approaching this ideal. Its peculiarity is
that the opposed experiences and moods of life can come into their own in it. It
presupposes that the mind has been open to all sides of life and has not, in
blindness or cowardice, avoided what did not please. It knows both times when one must
live despite a deficit and times when one lives out of a fullness that seems
inexhaustible; but it comes out of these oscillations with a surplus that makes
superiority over the course of fate possible, a surplus that is not, like the child's,
merely a gift of nature, but is in good part self-acquired, is saved-up work. What has
been done and suffered in the past one works together with the new that comes. The
different experiences do not stand on the same level. It is the most earnest
experiences of life about which all other experiences of life gather without being
able to break them, but gather precisely so that by contrast with them they appear as
jest. In the certainty of the validity of values one can smile at reality's attempt to
be everything. One can smile at the constant \opage{89}shifts of value, at having got
something entirely different out of life than one had thought. Wilhelm Meister's
``Lehrjahre'' ends, after all, with the young man being compared with Saul, the son of
Kish, who went out to look for his father's asses and ended by being anointed king and
counted among the prophets. Perhaps one has also had the experience that the outer
forms of life we build up so laboriously as a defense for it, as a bed for its stream,
became precisely only hindrances to its free course. In short, the world has been
overcome in earnest and can therefore now be overcome in jest.

Humor strives to unite the tragedy of life with its comedy. The two elements are united
in it, now so that they are joined into one, now so that they make up an ordered whole.
When language formed the expression ``tragicomedy,'' it thereby designated especially
the preponderance of the comic, and therefore the word does not fit humor, where the
jest moves on the ground of earnest. The word has moreover a decidedly scornful or
belittling meaning. Or perhaps Frater Taciturnus is right, and the word designates
something one can neither laugh nor weep over, because it is neither comic nor tragic.
The word ``comitragic'' would fit humor better.\footnote{\textit{Kierkegaard: Samlede
Skrifter.} VI p.~392.}

Kierkegaard maintained that tragic art needs a historical background more than
comedy.\footnote{\textit{Samlede Skrifter.} VI p.~402 f.} Insofar as this is right,
it must presumably come from this: that earnest first fully asserts itself when we
stand before reality in its whole severity, while the comic can better romp in an
invented world. But in great comedy the background of reality cannot be lacking; it
concerns life, after all, just as tragedy does. Bergson has even assigned the comic a
quite special significance for reality by stressing its social character and its
social function.\footnote{\textit{Henri Bergson's Filosofi}, p.~46--51.} Kierkegaard's
statement\footnote{When Kierkegaard refers to Lessing, it must be noted that
\textit{Lessing} (\textit{Hamburgische Dramaturgie.} N.~87--92), in contrast to
Diderot and Hurd, maintains that in comedy as well as in tragedy the characters are to
be individually marked, whether they are taken from history or not.} must in any case
not be understood as if individual and private life, life in the narrow circle, did not
\opage{90}have its tragedies. It is the merit of the so-called realistic literature to
have maintained this through definite examples.\footnote{Cf.\ \textit{Etik}. XXX, 3.}
Every suffering can take on this character where it befalls a significant personality.

I cite an example from real life that can be said to mark an intermediate form between
the tragic and humor.

The young Swiss woman Adèle Kamm lay chained to her sickbed by an incurable illness and
was thereby prevented from working for the ends she had set herself. She felt with
anguish that she was shut out from life. But through resignation she reached a new joy
of life. Since she could not ``realize her ideal,'' she said, she would ``idealize her
reality.'' She had the idea of organizing a correspondence among young girls who, like
her, were chained to a hopeless sickbed and thereby excluded from taking part in
ordinary life. The organization became very extensive and reached beyond all
confessional barriers. Since the healthy can never properly understand the sick, this
fellowship became of great significance. The young patients communicated their moods
and experiences to one another and exchanged their thoughts on what kept them up. They
thereby supported one another in their struggle still to get something out of life. The
feeling of loneliness vanished, and Adèle Kamm now felt herself to be one who took part
in the great struggle between good and evil. Thus she kept herself ``joyeuse dans
l'affliction.'' ``Happiness,'' she said, ``is more precious when it is mixed with
despair; peace is more complete when it follows the storm, light more radiant after
darkness.'' Relying on her own experiences, she amused herself sometimes by teasing
doctors and clergymen when they came with their prescriptions.\footnote{\textit{Paul
Seippel: Adèle Kamm.} Lausanne 1912.}

The tragic in suffering is that it hinders the work for what at first seems bound to
bring unity and meaning into life. The comic can come forth in that suffering can
nevertheless show its impotence before a new spiritual upsurge that it has itself
indirectly called forth. And a second time the comic lies in others' vain attempts to
treat soul and body by traditional formulas.

\opage{91}That there is pain and suffering in the world is a fact to which no earnest
view of life closes its eyes. It is a fact that sets us before great problems. Is there
any meaning in there being suffering? Even very pious people can deny that there is
such a meaning, and declare that when they believe in God, it is in the certainty that
he cannot be the cause of evil.\footnote{See, for example, \textit{Souvenirs de Marie
Flournoy-Burnier.} Genève 1910. (Not available in bookshops). Painful experiences let
this highly gifted and deeply religious woman find her only comfort in the thought that
all suffering was owed to nature, which God had not created. With a consistency more
possible for a Calvinist than for a Lutheran, she maintains: ``Si Dieu était
véritablement dans le monde, il serait partout le vainqueur.'' --- She did not know that
other theists had reached the same conviction, personalities as different as Plato,
Holberg, Rousseau, Voltaire, Stuart Mill and William James.} It is, to use
philosophical language, the contrast between value and reality that comes forward here,
and in particular the difficulty that not all reality can be understood from value. The
realm of reality is far more comprehensive than that of values. Only the most elementary
values, those immediately tied to the persistence of individuals and races --- and
hardly even these --- can be understood purely biologically. And yet, things being as
they are, pain marks a higher stage of development when compared with a state that is
painless only because it is limited. Pain is often tied to the transition from one form
of life to another, and it can be a symptom that full harmony among the tendencies of
life, or between life and its surroundings, has not yet been reached. That it should be
absolutely necessary for such states of transition and crisis to be bound up with pain
cannot be proved. It must be taken as a blind fact. But as circumstances now are,
suffering is a sign of life, a testimony that the struggle between value and reality,
or rather between valuable reality and valueless reality, is still going on. One could
escape much suffering by lowering one's level and reducing one's demands on life to a
minimum.

It is this actual connection that the great humor looks in the face. It does not
engage in speculations to \opage{92}justify the sufferings in the world, just as little
as it acknowledges theological explanations of the fact of suffering. But it takes this
fact up into its view of life, into its reckoning of moods. It is no small work to fit
the bitter and the dark too into one's picture of existence as a whole; but this work
can still leave strength over, so that the small and absurd can be met with the smile
it deserves.

25.~The great humor is an art of living. It leads beyond all aesthetic play by its
clear insight at once into the disharmonies of life and into the values of life. It does
not dwell, like art in the narrower sense, on single pictures of life. Poetry and the
visual arts can only give examples of how things may go in life, but cannot decide
whether these examples mean rule or exception. This holds both when idylls are depicted
and when tragedies are depicted. Of course one can, through investigations of the
subjects an artist chooses, as well as of the way he treats them, come to know his own
mood of life. Here is art's open place, where it is most intimately connected with life.
There is an art of living behind poetry and the visual arts, an art of living in which
the experiences of life are gathered in a practical intuition. It means no scientific
reckoning, although the scientific knowledge available may be woven into it. The great
humor points beyond every limited experience that would determine the whole life of
feeling, as well as beyond every picture fantasy is able to form. Over against such
limitations it uses jest as a liberating means. But it is not itself jest. When the
downfall of something valuable is experienced --- even if what went under had until then
stood as the highest value --- the great humor finds new confirmation that no single
value can comprise all value. By and in its very downfall the single value points beyond
itself.

When Goethe stood dismissive toward the concept of humor, it was surely because he was
offended by romantic poets' burlesque and arbitrary play with everything, a play they
themselves perhaps regarded as an expression of humor; but it was not the great humor
that \opage{93}lay at its ground. Goethe thought that humor had no hold or law in itself
and therefore easily, sooner or later, degenerated into melancholy or ill humor. (Letter
to Zelter. 30 October 1808). He even declares that one for whom life is bitter earnest
cannot be a humorist; the humorist is, he thinks, more taken up with his own mood than
with his object. (Unterhaltungen mit Kanzler von Müller. 6 June 1824).\footnote{Yet
Goethe admired Sterne. He found his humor inimitable. But he adds: ``Not every humor
frees the soul.'' (Sprüche in Prosa).} --- What Goethe is speaking of can evidently not
be humor as a total feeling, wrought through experiences of an active and passive kind.

In contrast to the concept of humor Goethe presupposes stands Solger's and
Kierkegaard's concept of humor as a view of life or mood of life developed with earnest
as its basis and understanding as its background, which does not let the exalted be
taken from it because it is connected with the low and paltry, but is able to unite the
smile at the latter with enthusiasm for the former. Kierkegaard expressed his view in a
very clear way by giving humor a place in the series of ``stages'' after the ethical
stage and before the religious. It was for him the culmination of paganism that it had
the strength of spirit to see the comic and the tragic joined into unity. Humor was for
him the last stage of life before faith. Naturally he thought only of what has been
called the great humor in what precedes, not of what he designated as the immature
humor, which is only devil-may-care, an aesthetic refinement that skips over the
ethical.\footnote{\textit{Skrifter}, VI, p.~393, 405, 414. --- VII, p.~249 f.}

There is an ``aesthetic'' way of taking life that does not have the character of great
humor. One then places oneself outside life as a pure beholder and spectator, perhaps
as a passionate spectator, since it can become a passion to look at things without
entering into personal relation to them and without heeding whether they cut into the
weal and woe of other personal beings. One perhaps feels precisely a need to withdraw
from every practical relation to things in order to live in fantasy. This kind of
aesthetic feeling is akin to intellectual interest, only it aims more \opage{94}at
living in images, while the intellectual need seeks concepts and laws. The question is
whether such an aesthetic-intellectual interest can fill the whole life of the
personality without narrowing and drying it up. Goethe's protest against what he
understood by humor is motivated by this question. In both scientific and artistic
activity it is only a part of the personality that is at work. But how do thought and
fantasy relate to the joy and sorrow, the enthusiasm and indignation, the exaltation and
dejection that are awakened by one's own personal fate or by the experiences one has of
the fortunes of the interests that make life worth living? Here a comprehensive
synthesis is demanded, the synthesis of the personality itself, within which the
intellectual and aesthetic syntheses are to find their place.\footnote{Cf.\ my paper
\textit{Aandelig Kultur.} (Mindre Arbejder. Third series).} The total feeling that
becomes ruling in a person, even in the greatest scientist or artist, is determined
quite as much by what was experienced outside the particular intellectual or aesthetic
sphere as by what was experienced within it.\footnote{In ch.~IX we shall return to this
point.} Experiences of a more individual kind will help to determine the ruling total
feeling, and besides, both science and art have their open places where they point back
to other sides of personal life and take their place within it as a whole. It is the
human being, not merely the scientist or the artist, who finds expression in the type
of life or the stage. The great humor is one of the possible types of life (stages,
syntheses of personality).

It is not in itself necessary that there be a great distance between the human being
and the scientist or artist. In the greatest they can stand close to each other. But
this closeness is not a matter of course and can presuppose great work. It is in a
certain sense human, ``all too human,'' to arrange things so that intellectual work
becomes something by itself, which one pushes aside when one has shut the door of one's
study in order to take part in life and take what it brings. But when that door is
shut, the new work remains of weaving the work and its fruits into life as a whole.
Otherwise it easily becomes a pitiful type of life \opage{95}that comes to correspond to
an intellectual work perhaps significant in itself. One does not escape having a
standpoint, but it may be a standpoint to match.

Some psychologists think that the feeling of pleasure in the tragic necessarily
presupposes that one does not forget that the whole is not earnest. One is conscious
that one is not part of it, but sits securely in one's seat in the stalls and watches
what goes on on the stage. In a picture by the French painter Vibert, ``Le Récit du
Missionaire,'' one sees a missionary showing his scars and telling of what he has
suffered for his faith; some fine and distinguished clergymen listen with great
interest to his story; the whole is evidently for them an aesthetic enjoyment that
captivates them while they sit softly and comfortably on their sofa. One does not get
the impression that the captivation concerns their whole mind. A complete and conscious
feeling of security must exclude the tragic effect. A clear representation that the
whole is not earnest will prevent the living into the course of the action, the
participation and suspense without which neither catastrophe nor resolution acquires any
significance. Even when we immerse ourselves in memory in our own earlier misfortunes and
dangers, we must (in this Spinoza is evidently right)\footnote{\textit{Ethica} III, 19;
47 Schol.} live through them anew as if they were present, so that the release in memory
can then call forth a feeling of pleasure.

In any case --- for the great humor we are describing here, what occupies the mind is
earnest. If it is a play, it does not take place on ``the boards that mean the world,''
but in the real world itself, the stage on which we ourselves play along, whether we
will or not, and where we can neither survey the course of the play nor know its end.

26.~Every total feeling can be threatened with dissolution, either by essential
elements falling away or by new experiences demanding greater energy or changed
organization. The great humor too must often fight for its life. The bright and
energetic harmony can, for example, be threatened by the clammy fogs of hypochondria,
which are called forth by inner organic dispositions or outer depressing experiences. A
depressed mood favors or summons all \opage{96}representations of imperfection,
misfortune, paltriness, inadequacy. Nothing in us or outside us seems to have worth.
Even if a world of thought is otherwise at our disposal, in the hypochondriac state it
is as if it had never existed, or as if it were an illusion that we lived in and for
it. A gray monotony rules, and the fruitful contrast between the great we looked up to
and the small we smiled at, this contrast that is the condition of humor, is blurred.
Now it is the great that cannot be made present, now it is the small that cannot be
seen in its worth. The tasks and events of daily life stand as indifferent, gray,
vexatious, obstructive. Even if the thought of the great has not vanished, the small
things are nevertheless not brought into relation to it, neither in the comic or tragic
relation of contrast nor in the intimate relation of sympathy and understanding, let
alone in all these different relations at once. Should a gleam of light fall into this
fog, it still brings no joy.

At its height hypochondria becomes the acedia, the inaccessibility to vigorous spiritual
life, that in the Middle Ages was counted among the deadly sins. Dante placed the
acediosi in hell, where they stir in the mire, and has them say (Inferno VII,
121--124):

\begin{quote}
Sullen were we\\
where the sun shone mildly on the meadows,\\
for sloth smoldered in our soul;\\
now we must hang dismally in the mire!
\end{quote}
% Danish (the translation Høffding quotes): Lidet glade / var vi, hvor Solen skinned
% mildt paa Enge, / ti Mismod ulmed i vor Sjæl den lade; / nu maa vi drøveligt i
% Dyndet hænge!

The old mystics and saints knew this state well and regarded it as the devil's work.
Even though they otherwise wished for suffering so as to show their courage and faith
through it, this kind of suffering was nevertheless only evil, because it was
unfruitful; it did not bring them, as they expressed it in their mystical-erotic
language, ``to give birth.'' It makes action impossible.\footnote{\textit{Joly:
Psychologie des Saints.} Paris 1898, p.~187 f. --- \textit{Vie de Ste.\ Terèse,
écrite par elle-même.} Chap.~XXX.} Boccaccio describes it in his commentary on Dante
thus: ``Such a person is unable to take any initiative, and if he is forced to it, he
can still accomplish nothing. His life flows away as if \opage{97}he were not living it;
his thoughts grow ever darker and darker; he shuns the company of others\dots{} First he
feels indifference, then reluctance, and at last disgust at everything good and
beautiful.''\footnote{Quoted after \textit{A.~Herrlin: Renaissancens Etik,} p.~50.
(Printed as manuscript).}

Not only can no initiative be taken in such states, but such states are so painful
precisely because of their indefinite character, since no definite cause appears that
one could try to take hold of. And it can go with the hypochondriac, if in his good
times he is a humorist, as it went with Harlequin, who asked a doctor's advice for his
melancholy and got the advice --- to go and see Harlequin. The jest that has helped many
a time to get round a sharp point does not help now, for the good reason that it cannot
make headway against the current. The hypochondriac could perhaps be cured if he could
get a free moment in which his hypochondria itself could be taken humorously --- seen as
a dark shadow in the landscape of life, which fades when one sees it from the height,
but overwhelms when one is down in the midst of it. But the presupposition for this is
that the height is clear, and --- that one can keep to the height. ---

This everyday example does not even point to the greatest shadows. They appear only when
the tragic threatens to tear the connection of life apart forever. How does the great
humor stand before this possibility? Can the tragic in life not become a testimony to
its impotence? --- Before we take up this earnest question, we shall see how the greatest
of all tragic poets stands before it.

27.~Shakespeare's mind was open to all sides of life, to all the contrasts it holds, and
he knew how to give his experiences expression in individual figures and fates. If we
have not ourselves had occasion to have experiences, we can come to stand before them in
him. He knew, as Spinoza knew, that nature is rich enough to let the most different and
conflicting individuals and fates come forth. As Cordelia and Goneril have the same
parents, so goodness and wickedness spring from the laws to which existence is once and
for all subject. And fortune \opage{98}and misfortune go their way through the world
according to the same laws. Wickedness can become so great that one is tempted to lose
faith in goodness, and misfortune can seem so overwhelming that everything seems to be
in dissolution. There can be occasion to cry out with Cordelia, when she hears of her
sisters' conduct toward the old father: Let pity not be believed! (King Lear IV, 3, 30),
and with Kent and Edgar when the old mad king comes carrying Cordelia's body: Is this
the promis'd end, or image of that horror? (King Lear V, 3, 203).
That Shakespeare's heart is in all this is felt through the whole mood and through single
moving words. His depictions are like crystallizations that have formed under the
influence of the deep movement his experiences caused in his mind. We stand here at
art's open place, where from the historically given poetic work we come to see into the
poet's own mental life, as it was tuned by what befell him or by what he saw in his
circle. Shakespeare scholars think that in the period of his life in which his mightiest
tragedies appeared, and which became the high point of his production, shattering events
took place for him or in his circle. It went with him as with Dante: he had been in
hell.\footnote{The women of Verona pointed at Dante when he went by and said: ``There
goes he who has been in hell.'' The Danish translator of \textit{Vita nuova}, Joh.\ Dam,
adds: ``This was not irony; he had been in hell, and everyone knew it.''} Of what kind
these events were one can perhaps form a representation through the two main tragic
motives that appear in his work.\footnote{\textit{A.~C. Bradley: Oxford Lectures on
Poetry}, p.~305.} The one motive is the relation of beautiful and noble natures to a
harsh surrounding world. The other is the relation of violent and demonic natures to a
surrounding world they seek to bend to their will. If the characterization of
Shakespeare's personality that A.~C. Bradley has given in his paper ``Shakespeare the
man'' (Oxford Lectures on Poetry) is right, he himself belonged to the open and fine
natures whom the blows of the world strike so easily, and who are also strongly affected
by what they experience of collision and dark fate around them. But he did not succumb
to the \opage{99}overwhelming, and perhaps his art was his best help here. He was at any
rate able in his presentation to let the characters, the evil as well as the good,
unfold in their natural consequences, and to let the fates, the unhappy as well as the
happy, go their quiet and great way. He was not prevented from revealing to us again and
again how the small hangs together with the great, the grotesque with the tragic, the
ludicrous with the exalted. Nor has he concealed that the tragic can come forth not only
through the collision of character with the world but also, as in ``Hamlet,'' can spring
from conflicts in the inner world of the mind between contending forces. Above all, the
horrors he looked in the face, or whose possibility his fantasy showed him, did not
prevent him from showing us the nobility of the intimate and deep mind in misfortune and
downfall itself. As has been said:\footnote{\textit{A.~C. Bradley: Shakespearian
Tragedy}, p.~324.} what happens to a being like Cordelia is not the decisive thing; the
decisive thing is that she exists! Philosophically we could express this thus: Shakespeare
does not lose his faith that there is value in the world, despite the fact that the
valuable so often goes under.

In Shakespeare's last great work, ``The Tempest,'' the feeling of peace and wistfulness
gain the upper hand. The tragic is softened and the comic vanishes. But still at the very
end, when Prospero thinks he has overcome all evil, so that he gives himself up calmly to
festive visions, quite forgetting the human beast Caliban and his conspiracy, he is shaken
through by the suddenly emerging thought of the disharmonies that still remain. He has
not been able to make the human beast into a human being. And like the festive visions
that had to be driven off so suddenly, the whole world we know, with its harmonies and
its disharmonies, will one day dissolve into a chaos! But from the thought of tenacious
wickedness and of the fleetingness of the glorious his soul is then saved through
magnanimity toward those who had acted wickedly against him.

{\small\textbf{Note.} In several places in this treatise, in agreement with various
Shakespeare scholars, it has been assumed that from the great poet's works one must be
able to infer back to what he himself had \opage{100}experienced or seen at nearer or
greater range. \textit{Georg Brandes} also proceeds from this assumption in his book on
Shakespeare, whose chief merit, for the rest, lies in its fine analyses of the content of
the single works. On the other hand this assumption is decidedly denied by Shakespeare's
most learned and most critical biographer, \textit{Sidney Lee.} ``The biographical
facts,'' he says (\textit{A Life of William Shakespeare.} New edition. 1915, p.~417),
``give no support to the assumption of a prolonged personal experience of tragic
suffering. Nor is this nebulous theory supported by the whole tendency of his literary
work\dots{} To seek the key to Shakespeare's triumphant conquest of the peaks of tragedy
in his personal experience, which must after all necessarily have been limited, is to
underestimate his creative power and its magic force.'' --- But the question is precisely
whether high artistic creation is psychologically possible without experiences of life
lying at its ground. And if this must be denied, the assumption that the poet built on
such experiences will by no means lead to thinking less of his art. Rather one would have
to admire that art could shape the earnest subjects in so perfect a way.\par}

28.~Is there an absolute tragic? Are there fates in which pain and downfall become the
last word? If so, we stand here at a limit of the great humor. And if there is a limit,
the humorist must himself acknowledge it, and it must increase his wistfulness and make
his smile less bright, if there still remains any possibility of a smile at all.

In a letter to Zelter Goethe says: ``I am not born to be a tragic poet, for my nature is
harmonizing (`conciliant'). Therefore the purely tragic case cannot interest me, since it
must properly in its ground be impossible to harmonize.'' In his conversations with
Chancellor v.~Müller he comes back to this thought. ``Everything tragic rests on a
contrast that cannot be evened out. As soon as evening out sets in, or is seen as
possible, the tragic vanishes.'' ``Tragic I call a situation from which no way out is
possible, no reconciliation (`composition') conceivable.'' (Unterredungen mit Kanzler
v.~Müller. 6 June 1824 and 2 July 1830). He criticizes Hegel's conception of the tragic as
resting on a conflict between two powers each justified in itself. Hegel's favorite
example was ``Antigone,'' whose main subject in his view was the conflict between family
and state. Such collisions take place, after all, within one and the same ethical whole;
family and state are both members of the life of mankind, and since \opage{101}this more
comprehensive whole is maintained precisely through the downfall of one-sided endeavors,
tragedy leads to a clearer understanding of life. The properly tragic is thus for Hegel
not the catastrophe but the conflict, and the conflict always has the necessary function
of maintaining the mutual limitation of the different ethical powers.

Hegel's bright and energetic optimism would, if it could be held fast, widen the limits of
the great humor as a view of life. But Goethe is unfortunately right against Hegel. There
can be, and there is, absolute tragedy. For Hegel the fate of single individuals is
indifferent, so long as the great ideas constantly prove their power. But even if a
person's worth rests on the cause he serves, should it then be indifferent whether he
himself, whose soul is filled with this cause, persists or goes under? Tragic suffering
cannot be indifferent, however great the weight one lays on the often (not always)
fruitful conflict between different values that have not found their limitation. The
ideas (the values) are, after all, represented in the world of experience by personal
beings. And, as we have seen (§ 27), it is precisely the way in which the tragic hero in
his very downfall reveals and maintains his exalted disposition that conditions the proper
working of tragedy on our mind.

On the famous passage in Aristotle's Poetics, commented on countless times, where he finds
the effect of tragedy to consist in its awakening terror and pity in such a way that these
feelings are thereby ``purged,'' there has been disagreement especially as to what is
understood here by ``purging.'' It is perhaps most correct, with Erwin Rhode
(Psyche\textsuperscript{2} II, p.~49), to think of an analogy to the way in which music
can at once awaken strong feelings and --- precisely through the awakening and the
discharge --- bring it about that they do not rule over the mind in a painful and
inhibiting way.
Of more recent music one could think here of Beethoven's Fifth Symphony, ``the Fate
Symphony,'' where one first hears fate's heavy blows on the gate and, by being permeated
with the mood thus awakened, is prepared to be affected by the at once cheerful and wistful
impression of the symphony's last part. But what Aristotle at any rate says nothing about
is what \opage{102}character ``terror and pity'' take on when they have been ``purged.''
A.~W. Benn has therefore called him a poor psychologist.\footnote{\textit{Aristoteles'
Theory of Tragic Emotion.} Mind. 1914, p.~89.} But the point is perhaps that the Greeks
did not have the understanding of the metamorphoses of the mind and the shifts of feeling
that was won later. What Aristotle could not express in his psychological language may
have been precisely what Benn himself maintains: that pity toward heroic characters passes
over into admiration, and toward gentle and self-sacrificing characters into love. And
terror must be assumed to undergo a corresponding change.\footnote{Benn interprets
Aristotle as meaning that pity passes over into terror, and that this is then calmed by
the artistic mood (plot-interest) gaining the upper hand through one's knowing oneself to
be in safety. I cannot find anything about this in the text. It is both terror and pity
that are ``purged.''} Terror and pity do not vanish, but are ``purged'' by becoming
elements in a more complex feeling.\footnote{Recently the ``psychoanalysis'' founded by
\textit{Freud} has maintained that certain forms of mental illness can be cured by a
painful experience in which the real cause of the illness lay being drawn clearly and
definitely forth into the patient's consciousness, which is thereby ``purged.'' The
method is accordingly called catharsis, the same expression Aristotle uses.}

29.~In the tragic, humor has one of its limits. The humorist does not bend himself
together over the wound like the tragic hero. His wounds are not mortal and need not drive
away the smile. The tragic hero can show greatness even in downfall, but pain and sorrow
exclude the possibility of the smile.

The humorist knows that life has its tragedies as well as its comedies. His experience of
life, both what is owed to his personal fate and what is owed to experiencing the fates of
others, can be rich and comprehensive. And his striving aims, consciously and
unconsciously, at letting the great contrasts have their justified influence on the timbre
of the total mood. But he must admit that where the tragic is of such a nature that there
is no strength or room for letting the bright elements of the mind have influence, there no
organization or fusion is possible. The cosmic harmony in which the humorist consciously
or \opage{103}unconsciously believes collapses where there is an absolutely tragic fate ---
where suffering and downfall are one. Only the tragic hero, not the humorist as such,
maintains the nobility of mankind in such a catastrophe.

The humorist is willing to grant the tragic hero the highest place. But he does not give
way to anyone else either. And especially he does not give way to those who spend their
lives talking of what tragic heroes have done and suffered without themselves having
anything tragic in their existence. Moreover, the humorist protests against the tragic
being declared, in defiance of experience, to express the innermost being of life. That is
only romanticism. The tragic hero himself will have no time or sense for such assertions;
he feels himself standing precisely out at the limit of all human relations and all human
ability; all his energy, although it surpasses what the humorist as such commands, goes
into standing where he stands. The highest pathos makes one silent, or at least sparing of
words, however talkative those may be who know this pathos only at third or fourth hand.

The tragedy of life that lies beyond the great humor comes as an earthquake comes. It is
perhaps announced, like this, by preceding tremors. But it does not depend on ourselves
whether we experience anything of the kind. Perhaps by omitting an action that our ideal
of life demands we could avoid a tragic fate. But then the ethical weight lies on the
omission of this action, and this point belongs to the discussion of the relation between
action and humor that a later chapter will bring. Just as action marks the limit of humor
from the active side, absolute tragedy marks this limit from the passive side. The humorist
thinks he has honestly and fairly sought and followed the right way without having failed
any task that could have brought about tragic conflicts. He admits that this way of his has
led through mild regions, and that the tension his stage, like every earnest stage of life,
has brought is to be regarded as harmony in comparison with tragic pain. He perhaps dares
add that he would not draw back if a tragic fate could not be avoided without harming
his soul. His admiration for the tragic heroes is grounded in understanding. Partly
\opage{104}he can very well put himself into situations that for his own part do not
go beyond the world of possibilities, partly he recognizes that tragedy reveals a limit
to the harmony between value and reality. But he does not renounce being a hero because
he is not tragic. He fights his fight to assert himself through holding fast his
conviction of the coherence and worth of life despite all disharmonies in detail. This
fight can lead him right up to the limit of the tragic. If this limit is to be crossed,
the question becomes whether he can undergo the metamorphosis without which the tragic
hero does not come to be.

In Gunnar Gunnarson's magnificent book ``Borgslægtens Historie'' it is said of a young
man who wanted to give up everything to wander about as an ascetic and prophet (IV,
p.~117): ``It did not for a moment occur to him that he had already suffered a defeat ---
that this life he had resolved to lead had no root in his own soul, no presupposition in
his earlier life. He did not know that suffering is a good of fate that is granted to
few people in so full a measure that they become wholly able to look away from their own
weal and woe and live and think only for others and feel satisfaction in it. He did not
know that only in pain's elect does all self-love die.''

The relation between the tragic on the one side and the harmonious-human stage of life
that humor represents on the other cannot be better expressed. The suffering that becomes
tragic is not arbitrarily produced or sought out. It comes when fate knocks at the gate,
without this knocking (as in Beethoven) being the introduction to a symphony.

30.~For Kierkegaard humor was a stage that has the religious stage above it. And all true
religiousness was for him of a tragic character. In the person who is to be in a state of
religion, the representation of God must, in Kierkegaard's view, work transformingly on
his whole conduct of life. No human relation can have value in itself; it has value only
when it is consciously subordinated and transformed from the viewpoint of eternity and
infinity that the representation of God leads one to adopt. From the \opage{105}contrast
between the infinite demand and human finitude springs a suffering that cannot be
harmonized. In Christianity --- mind you, original Christianity, primitive Christianity
--- the whole of life is led with a view to the beyond, or rather to the one place in
time and space where the divine has appeared. There is, according to primitive
Christianity, only one value, the one that was put into the world with the God-man.
Everything else is forsaken by value or hostile to value. Primitive Christianity is a
funeral march. For the true Christian life is suffering, a suffering called forth not
only, as in all religiousness, by the absolute demand of the representation of God, but
by there being value at only a single point in existence, while everything else lies in
darkness.

I shall not go more closely here into whether Kierkegaard is right in this
characterization of primitive Christianity. I shall only repeat what I have said in
several places (in my book on Kierkegaard and in my Philosophy of Religion), that
primitive Christianity was not a funeral march but a victory march, since the mind was
wholly filled with the expectation of Christ's near return, the real founding of ``the
kingdom of heaven.'' Since, in contrast to this great, living expectation, everything else
had to pale, all other interests and tasks vanish, suffering cannot, as Kierkegaard
thought, have been ruling. Before the great images in which the expectation found
expression, not even the pain of tearing oneself loose from earthly and human relations
could assert itself. ---

The tragedy that Christianity was, according to Kierkegaard, far surpassed for him
everything that even the greatest tragedian could write. ``O my friend,'' he says (``Sygdommen
til Døden.'' Skrifter XI, p.~236 f.), ``what have you really attempted in life?
Strain your brain, tear aside every covering and bare the entrails of feeling in your
breast, raze every fortification that separates you from the one you read about --- and
then read Shakespeare, and you will shudder at the collisions. But the properly religious
collisions even Shakespeare seems to have shrunk back from. Perhaps too they can be
expressed only in the language of the gods.'' ---

When Kierkegaard calls the religious stage, as he conceives it, tragic, it is a tragedy
that in his \opage{106}view it is man's duty to enter into. It is willed tragedy. It does
not come over man, but man is to come to it. It is not owed to collisions that fall only
to the lot of some, not of all. But it is the highest form of human spiritual life.

In the two pseudonyms Frater Taciturnus (Stadier paa Livets Vej) and Johannes Climacus
(Uvidenskabelig Efterskrift) Kierkegaard presented representatives of what he understood
by humor. They both stand hat in hand before the religious stage. This is consistent from
Kierkegaard's own standpoint, according to which suffering in the end became the decisive
measure of the height of spiritual life.\footnote{Cf.\ \textit{Danske Filosofer,}
p.~165--67.} But it does not agree with the concept of humor as it has been developed in
what precedes. According to this, a tragic relation is a fate that can bring about a
serious trial, but not something it is everyone's duty to experience.

For humor all value is not contained at any single place in time and space. There is a
constant possibility of discovering and experiencing new values in the world, at the same
time as there is a clear consciousness that if the proposition of the persistence of value
holds, it is at any rate not the empirical values that persist. If there are eternal
values, they make themselves known through the laws for the arising and passing away of
empirical values, could we find these laws. Eternity and time are for the humorist not
absolute opposites, so that there would necessarily have to be a painful relation of
tension between them. Eternity is not something awaited or lying far out in time. It is an
ever present, it designates the great lawfulness that makes all development in time
intelligible. We constantly presuppose it, whether we look back or go forward.
Kierkegaard's proposition that we understand backward but live forward is not entirely
right. For that great lawfulness is of significance for our striving forward no less than
for our looking back; through the memory of the past what we can do in the present to
prepare the future is determined. Humor does not see reality as a falling away from
eternity, but sees precisely every reality as a single stretch of an eternal stream.
Kierkegaard's humorists see, in romantic fashion, the world of memory \opage{107}as a
world of dreams, over which one can forget the ever onflowing existence in time. But there
is an intimate bond and a constant interaction between memory and life. The world of memory
is constantly peopled from the experiences that the life of action brings, and the life of
action is animated by what is preserved in memory. In the discovery of new values memory
always works along.

The possibility of tragedy is not to determine the whole of life or the view of life. When
humor means faithfulness to life, because it gives the paltry and the obstructive its place
as well as the great and the helpful, it also leads to going to meet the tragic fate if
this is really a consequence of having been faithful to reality and its tasks. Here, then,
the contrast between humor and the tragic falls away, and there appears the possibility of
an energetically carried through, continuous transition from the one stage to the other.
There is then no leap, or, if there is a leap, continuity is not broken when the leap
really carries over to the other side. A leap does not, after all, without more ado mean a
fall. The free man, whose thought busies itself not with death but with life, will
nevertheless go to his death so as not to deny what alone makes life worth living and
thinking about.

The humorist may well be dismayed, like Prospero when, absorbed in his visions, he had
forgotten Caliban and his conspiracy. (How many have not fared like Prospero in the year
1914!) But behind the great humor there always lies faith in the possibility of the
release of new values. If all the sources of value in the world were exhausted, there would
no longer be any possibility of the total feeling we call humor --- indeed, of any total
feeling whatsoever. The same would be the case if existence should prove to be so
tragicomic that there were bound (potential) values enough, but no possibility of releasing
them. This the humorist could not take humorously. He would then have to admit that tragedy
got the last word. Indeed, if it were even tragedy, with all the inwardness and loftiness it
can call forth! But if the sources of value of existence ceased to flow, a progressive
weakening and dulling of human spiritual life would set in, and the last act of the drama
of human life would hardly bear the stamp of the sublime.

\opage{108}The great humor stands and falls with the presupposition that what passes away
are only single values given in experience. Their passing away is the tragic element in the
experiences from which the great humor proceeds. But this element does not become sole ruler
as long as new sources of value are discovered. In the very magnanimity that can reveal
itself when a value passes away, a new value appears. And after that life is taken up anew.
As Kent, when he is declared an outlaw, ``shapes his old course in a country new,'' so the
fight for the values of existence is taken up after a defeat and carried on perhaps in new
forms. The dying Hamlet with his last breath gives his voice in the election of the king to
the young, bold Fortinbras, whose energy and daring will lead him to other fates than those
that fell to the lot of the melancholy dreamer.

31.~Whether Kierkegaard is right in setting the sharp contrast between humor and religion
depends of course on how one defines the concept of religion. If Kierkegaard is right in his
definition, the humorist takes a dismissive attitude toward religion and refuses to
acknowledge it as a higher stage of life. But it must always be acknowledged as a great merit
of Kierkegaard that he so definitely and positively acknowledges and grounds the existence of
an independent, high and noble human spiritual life outside (in his view prior to) the domain
of positive religion. This acknowledgment was closely bound up with his life's cause, that of
letting the Christian appear in its strict purity. One cannot exactly say that this side of
his thought has been the object of much understanding since his death, either in the Church
or outside the Church. In particular the level in this respect seems to have sunk within the
Church, whose preachers have an increasing inclination to think that everything the words
spirit and soul can mean comes under them. It was nevertheless their natural duty and task to
draw attention to the independent and high spiritual life that has been lived and can be
lived outside the Church, so as to avoid confusing and mixing up Christianity and humanity.
But it is naturally easier to preach against ``reason'' and the evil will than to enter into
stages of life that are built up on \opage{109}a wholly different basis from their own. And
yet many will have had the experience that the more they have freed themselves from dogmatic
presuppositions, and also from the antipathetic and polemical moods that the work of
liberation easily breeds, the more they are convinced that there is an independent human
spiritual life that has developed historically under influence from many sides, from
religion and from science, but first and foremost from and in life itself, a human life that
in turn can appear in very different types. One of these types is the special subject of
this treatise.

As for the humorist, it will hardly matter much to him whether one wants to call him
religious or not. He will perhaps take a humorous attitude toward the various attempts to say
what religion ``really'' is. Yet he will (so I must of course think) have no reason for humor
toward my attempt in the philosophy of religion, since the definition I have proposed is
expressly set up partly as a hypothesis in the history of religion, partly as a
psychological hypothesis. He will perhaps even feel a certain spiritual kinship with my
conception of religion as \emph{faith in the persistence of value}. For the great humor
becomes possible only through a holding fast to the fact that value in existence does not
vanish because single values given in experience vanish.

It is a fact of the history of religion that people's faith is determined by what they have
learned in their lives to ascribe value to. The gods stand as the powers that maintain these
values, and heaven is the sum of the values one knows. One can of course in one's
investigations of the historically given religions keep to the outward forms of cult and the
dogmas without asking what psychological presuppositions they rest on. One can be a very
learned man without being a psychologist. But it is at any rate worth an attempt whether one
might not be able to penetrate to the psychological basis. It goes without saying that the
further back we go to primitive forms of religion, the less information one will be able to
get about the way in which the organized forms of cult came to be. Even so elementary a
religion as the Central Australian is known only as a finished whole. And yet it must have had
its history of development like everything else in the world, \opage{110}and every attempt
to give an account of it leads back from sociology to psychology.\footnote{See my paper
\textit{Sociologi og Religionsfilosofi.} (Critique of Durkheim's book ``Les formes
elementaires de la vie religieuse'') in the periodical ``Vor Tid.'' 1914--15.} In the
religions we call the higher, the psychological basis comes plainly forward, since the divine
beings in whom one believes are conceived as champions of truth and goodness. If there lies a
fear at the ground of religion (a proposition that cannot be set up as general), it is a fear
of losing what --- at the given stage --- gives life worth. The proposition of the persistence
of value lies at the ground of all religion in a way similar to that in which, for example,
the causal principle lies at the ground of practical common sense when it seeks means to
attain its wishes, or when at striking events it involuntarily looks about for what may have
gone before. And just as little as practical common sense needs to know the abstract causal
principle or what lies in it, just as little does religious consciousness need to know the
abstract proposition of the persistence of value, its consequences and the problems it poses.

Whether the need to hold fast the faith in the persistence of value will maintain itself when
it is not supported by being clothed in the forms of cultic life and dogma, only experience
can decide. It is a psychological task to investigate the conditions for the persistence of
such a need, just as it is, for example, a psychological task to investigate the conditions
for the arising and persistence of the great humor.

It is the relation of tension between value and reality that conditions the most far-reaching
human experiences and thereby in turn determines the persistence of total feelings. Insofar as
the humorist takes an interest in the spiritual life of the past, he will be able to use the
proposition of the persistence of value as a working hypothesis. With every dogma and every
cult he will ask what values one has here felt, or feels, the need to maintain or see
maintained. In the worship of Kurotrofos, Demeter, Isis and the Madonna he thus traces the
need to maintain mother love as divine. In the series of different representations the Greeks
formed of Zeus at \opage{111}different times he will trace an increasing tendency to regard
justice as the highest quality. But not only for the understanding of the past, also in
action in the present, that working hypothesis is needed. In all striving for an end it must
be presupposed that existence contains the possibility of the release of values. Otherwise the
hope that must bear all action falls away. In a world forsaken by all possibilities of value,
action is impossible.

This consideration forms a natural transition to investigating the relation of the great
humor to understanding and to action.

% ===== VI. UNDERSTANDING AND HUMOR (pp. 112--123) =====
\chapter*{VI.\quad Understanding and Humor}
\addcontentsline{toc}{chapter}{VI. Understanding and Humor}
\markboth{Understanding and Humor}{Understanding and Humor}
\opage{112}32.~The background humor presupposes is conditioned in good part by the
understanding a person has of himself, of the conditions of his life and of the existence
to whose laws he is subject. The understanding in question here can appear in a great many
degrees and forms. But if it is to acquire significance, it must satisfy two demands.

It must be so comprehensive and so coherent that the person can have the representation of
a great order of things in relation to which the disappointments, contradictions and
sufferings that life brings can stand as something that nevertheless also belongs to
existence and has its right to be, while at the same time it is unable to break the trust
in the great connection within which the values of life unfold. The resistance one meets
perhaps surprises for the moment, but is nevertheless fitted in as a member of the series
of experiences and helps to determine the timbre of the total feeling. Luther, according
to the legend, was disturbed by the devil during his work on the translation of the Bible;
but when he had discovered whom he had before him, he only said: ``Stay there as long as
you like!'' The devil had, after all, his place in his order of the world and could be
taken humorously. Every view of life has its devil --- rebel, tempter, accuser --- whom it
is a matter of putting in the place that belongs to him. There always remain scattered
elements from the chaos out of which the development of character proceeds, and in the
course of life new, unorganized and unfused elements can be added that can put both total
feeling and understanding to the test. (Cf.\ § 26.)

The second demand the understanding must fulfill is that it must be able to assert itself
in an immediate way and with \opage{113}gathered force. It must appear as a practical
intuition, as a certainty that is the fruit of the person's experience and reflection.
There need be no express consciousness of the sources from which this intuition stems. If
attempts are made to give reasons for the intuition, it will as a rule turn out that these
reasons were found or constructed afterwards and have nothing to do with the causes a
purely objective investigation discovers. But however this may be, the understanding, if
it is to be the background of humor, must have the character of concentration, must stand
as a firm coast against which experiences break without being able to break it.

Contributions to the understanding in question here can be made from very different sides.
It may be owed now more to the head, now more to the heart. It may be science that for the
most part forms the intellectual background, or it may be a religious circle of
representations, speculations or poetic visions and images, or perhaps merely the
involuntary survey of things of sound common sense. The decisive thing is that a firm
connection of thought has formed in one way or another. Spinoza rightly says that feelings
firmly tied to a coherent train of thought (affectus secundum ordinem ad intellectum
ordinati et concatenati) are harder to dissolve than vague and incoherent feelings. Yet it
is of significance that the work of reaching a firm connection of thought can always be
taken up anew. If the devil cannot be understood within the connection we have so far
found, his unexpected appearance can perhaps give occasion to see that we have worked with
too narrow a horizon.

The intellectual background is only a single condition of humor. This is seen from the
fact that even where this background is the same in several different individuals, it is
not said that the great humor develops in all of them. In the treatise's last two chapters
I shall go more closely into this question. Here I shall only remark that there will
nevertheless be a certain kinship between the intellectual background and the basis
(superiority, loftiness of mind, wistfulness, longing, sympathy; see ch.~IV) that the
great humor presupposes. Both this basis and that background lead, \opage{114}each in its
own way, the mind beyond the narrow horizon within which life so often moves. (Cf.\ § 13 c
and § 23).

33.~However different the worlds of thought that form the intellectual background of the
great humor may be, they must all be closely connected with the way in which, at a given
time and in a given place, one distinguishes between imagination and reality. The measure
by which this distinction is made, that is, the criterion of reality, is always the
starting point for the development of the world of thought. And at every standpoint this
criterion consists in the firm connection one thinks one can find among the observations
at one's disposal --- in the possibility of a total connection into which the single
experiences can be fitted. It never consists in anything altogether single and unique;
even if a single event or a single observation should have made such an impression,
awakened such a feeling, that everything else loses its significance in relation to it,
there will still always be an endeavor to find a relation between that single experience
and all the others. These will be interpreted by it when it cannot be interpreted by them.
The mystics of the Church interpreted their ecstatic experiences by their relation to the
connection maintained by the Church's dogmatics;\footnote{Cf.\ my
\textit{Religionsfilosofi} § 62.} where ecstatic states occur in psychologically
experienced people in modern times, on the other hand, a definite contrast will be
maintained between what is experienced in such states and the dogmatic representations in
which current religious life rests.\footnote{\textit{Th.~Flournoy: Une mystique moderne.}
(Archives de Psychologie XV. Genève 1915), p.~99 ff.; 122 f.; 134 f.} The relation between
experience and interpretation is altogether different at the different stages of
development, as Auguste Comte showed in his doctrine of the different stages in the
development of knowledge. Scientific interpretation consists in letting one experience
illuminate the other, until they all at last prove to stand in a great inner connection,
so that one can perhaps even derive the following experiences from the preceding ones.

Already the task set to every person, that \opage{115}of finding his bearings in the world,
of distinguishing between what is observation, what is memory and what is fantasy, will
lead him to see himself as a member of a great connection. What he calls reality then
becomes the whole that he consciously or unconsciously forms from the observations and
memories at his disposal. That a reality exists for him he shows through the actions he
performs; they presuppose, after all, that he thinks himself standing before something
other and more than his own representations and fantasies. Here the old simple mark holds:
Show me thy faith by thy works! --- How large the whole is that makes up reality for the
single person, and how firmly it is joined, can be very different. That is why the human
pictures of reality are so different. But the law by which they are formed is always the
same, from the animistic conception of nature of savages to the law-bound cosmos of modern
science. In the endeavor to work together as many observations and memories as closely as
possible, the great concept of truth makes itself known, which under different forms and
in countless degrees lies at the ground at all stages of human spiritual life.

The humorist is a realist. His elevation above the single has the great connection of
reality as its background. On this it rests that humor (especially Shakespeare's humor) can
be called faithfulness to reality. Faithfulness rests on the presupposition that the one or
the thing I am faithful to remains one with itself under all changes, develops in
continuity with itself. No single observation is enough here. There must be formed,
involuntarily or with full consciousness, a representation, a conception, that can
correspond to all the different observations in which the object of faithfulness comes
forward, such a representation as I have called in my Psychology a typical individual
representation. Faithfulness presupposes that this representation will continue to hold,
that the law by which the object has so far developed will be the expression of its
development further on. And faithfulness itself then consists in a continuity in our own
conduct corresponding to the continuity in the object's being. It is neither for the object
nor for faithfulness itself unchangeableness that is presupposed; the
\opage{116}presupposition is only that the changes take place in the same direction and
according to the same law.\footnote{Cf.\ \textit{Psykologi.} V B, 9; VI C, 6. ---
\textit{Personlighetsprincipen i Filosofin,} p.~18--27.}

Faithfulness to reality will then show itself in our trusting the laws that have so far
proved to hold for our experiences, and being willing to take them as they are, with the
goods and evils, the tasks and labors, the means and hindrances they bring with them.

One might object to the close relation maintained here between humor and what stands as
reality that Falstaff must surely be called a humorist, but that he lives, after all, on a
great lie. To this I would answer that Falstaff's comic effect on us rests precisely on
his appearing, within the reality he constructs for himself, a reality in which everything
takes on the opposite meaning from what it has in the ``real'' world --- on his appearing
within this world with faithfulness and consistency. He explains his fatness by sorrow and
adversity blowing a man up; he reduces the concept of honor ad absurdum because honor
cannot set an arm or a leg on again; and like the old toper he is, he defends his drinking
by saying that young people must live too! With full earnest the whole system from which he
passes his judgments is built up. The world order he constructs lets precisely individuals
like him stand as justified. His humor is faithfulness to the picture of reality he has
formed and which he holds fast undauntedly against his friends' criticism and the
intervention of outer circumstances.

When Falstaff not only appears himself as a humorist but is also taken humorously by us, it
is because his picture gets its place and its significance through the effect of contrast
within the whole picture of the earnest time and its struggles that Shakespeare unrolls in
``Henry IV.'' He must be seen in this whole connection to be able to put us in a humorous
mood. Sidney Lee thinks that when Falstaff works humorously, it is because the contrast
between his old age and his wanton way of life brings in the melancholy element that cannot
be dispensed with in humor. (Life of Shakespeare p.~245.) But such a moral consideration
would quite \opage{117}surely make the humorous effect impossible, as would happen with
Mephistopheles in Goethe's ``Faust.'' One must take God's standpoint in the Prologue and
say that there must, after all, be such fellows too.\footnote{Just as God regards
Mephistopheles humorously, so too the reverse. Mephistopheles' humor toward God (``Von Zeit
zu Zeit seh' ich den Alten gern'') is conditioned by his seeing God as an element in his
own, Mephistophelian world order; from this he could on his side have said of God: ``Es
muss auch solche Käuze geben!'' --- Such mutual humor can also occur on a smaller scale. An
old schoolfellow of mine, who had become an ardent politician, once said to me, when we had
been talking about my scientific interests, that, to be honest, he looked humorously on
everything of that kind and valued it only by its significance in the political market.
With the bonhomie he always kept, he added that he could well understand if I, from my
standpoint, also took his political striving humorously. I believe, however, that neither
of these quite came off.} It is by virtue of this principle that Falstaff makes his
humorous effect. That is why it jars and cuts when Prince Henry, as soon as he has become
king, casts off his old crony with harsh words. As Bradley aptly says in his essay ``The
Rejection of Falstaff'' (Oxford Lectures p.~259): ``Shakespeare had here created so
wonderful a being and set him so firmly on his intellectual throne that the attempt to
dethrone him again could not succeed.'' Shakespeare was not faithful here to the picture
of real life he had created, to the reality that offered not only characters like the
care-ridden king, the impetuous Hotspur and the heroic Prince Henry, but also had a place
for old Jack, his wit and his merits.

34.~Our picture of the world can and must be changed in many respects with advancing
experience, and even if we trust that the changes follow the law we have so far traced,
the concept of reality is nevertheless an ideal concept, and the concept of truth --- as
the greatest possible connection among the greatest possible number of observations ---
contains a summons to ever continued research.\footnote{Cf.\ \textit{Psykologi} V D. ---
\textit{Den menneskelige Tanke}, p.~257--262; 289 f.} An absolutely firm background is a
possible assumption only \opage{118}as long as uncritical security and dogmatic
confidence rule. The world of experience cannot be concluded in time and space; the causal
series cannot be carried through completely; events often appear that come so suddenly and
surprisingly, and are so different from earlier experiences, that it becomes a great task
to find the connection anew. But the great humor presupposes, after all, not only an
understanding of single experiences each by itself, but also the possibility of all
experiences being able to enter as elements into a picture of the whole. And thereby an
infinite task is set. And when can it be completely solved? Moreover, the single and
individual is itself a little world whose inner connection must be unraveled. And if it
stands as a torso or as a monstrosity, one must find what caused the destruction or
checked the growth.

There can further, when the concept of value is applied, appear a glaring contrast between
the general order of things, which calmly fulfills itself, and the fate of individual
beings, a contrast that can work now tragicomically, now comitragically. Even before the
modern doctrine of evolution with its idea of the struggle for existence appeared, Tennyson
wrote his famous lines on nature (In Memoriam 56):

\begin{quote}
So careful of the type she seems,\\
So careless of single life.
\end{quote}

Here it is aptly expressed how difficult it is to apply concepts of value to a unified
picture of existence.

First, the series of values cannot be concluded. This already follows from the constant
possibility of shifting: what at first has value as a mere means can acquire independent
value as an end, and what at first stood as an end, perhaps as a last end that bounded the
horizon of wishes, can acquire significance as a means to still more remote ends. Where then
could we find an absolute fundamental value on which a concluding valuation could be built?
And is it given that, even if we found a last measure, all single values would be
commensurable with it?

Second, value and reality do not, after all, coincide. The concept of value springs up
within the psychical world, but this \opage{119}seems to be only a vanishingly small part
of existence in relation to the material world with its enormous unfolding of masses and
forces. The world of values then stands, it seems, in a lonely place in existence, like the
broom bush in Leopardi's poem (La Ginestra) that grows on a single spot in the desert of
lava. In his bitter mood the poet holds this picture up before all those ``who have grown
used to praising life'': ``let them come here and see whether the weal and woe of living
beings really lies near nature's heart!''

All these difficulties do not, however, necessarily make the intellectual background of
the great humor vanish. The consequence is only what properly lies in the nature of the
case, that no intellectual background for human life of feeling can be set up once and for
all. A drama need not, if one does not maintain the demand for ``unity of place,'' be
played before the same backdrop through all the acts; the backdrop can change according to
the course of the action. As long as there is a connection among the changing scenes, unity
is maintained. A fixing of the intellectual presuppositions once and for all would even
make the great humor impossible. For the great humor turns not only against the single
experiences and lets them fade before the background. It turns also against the background
itself as subject to the law of change. Every background that would lay claim to absolute
validity would call forth the humorist's smile. Modern epistemology too assumes nothing
absolutely first or last, either as regards the forms of thought or its principles or its
results. We seek to build trains of thought that we can carry through over against all
present and coming observations. Where the applicability of such trains of thought proves
to have a limit, we investigate whether the limit might not be a consequence of the train
of thought we had begun being itself only a member of a more comprehensive train of
thought.
Light vanishes in the evening and comes again in the morning; but this vanishing
and arising means no leap in nature; we need only pass from the train of thought formed
only by observations from the earth's standpoint over to the train of thought that can
be built up when one proceeds from the earth's rotation and its position relative to
\opage{120}the sun. Galileo already said: ``Think the earth away, and there is no
sunrise and no sunset, no horizon, no meridian, no day and no night!'' What seems to be a
limit for thought is thus only a sign that thought has ended its work too early. --- This
holds for the series of concepts of value just as well as for the series of grounds and
causes.\footnote{Cf.\ \textit{Den nyere Filosofis Historie} I, p.~176. --- \textit{Den
menneskelige Tanke.} p.~196. (Critique of Renouvier's and Dühring's setting up of ``the
law of determinate number'').}

Understanding is a spiritual work that must always be taken up anew. The great humor would
build on sand if it overlooked this. The relation between understanding and humor falls in
a way under the relation between action and humor, which we investigate in the next
chapter. The humorist must take the claim to have reached an absolute background as the
expression of thought having grown tired and confusing its resting place with an absolute
conclusion. He knows that with advancing experience two things can happen: the results we
have won so far can fall away --- or the problems that have so far stood can fall away. Two
of our older Danish thinkers each seized one of these possibilities here. Holberg thought
that when we found no more problems in the world, it was because our experience was
limited; with widened experience one would see how illusory many of our solutions were.
J.~S. Sneedorff thought, on the contrary, that with widened experience many of our present
problems would fall away, because they are owed only to our narrow horizon.\footnote{%
\textit{Danske Filosofer}, p.~5 and 18.} They are both right, and the great humor can use
both views. It jests both with the absolute solutions and with the absolute riddles. And
this jest is grounded precisely in the great earnest with which the great humor embraces
the work of thought. Innermost it rests on a faith in truth as that which demands ever new
work to widen experience and to think through the experience won.

Just as the contrast between earnest and jest, sorrow and joy, tragic and comic, greatness
and paltriness, reason and absurdity is dissolved in humor as a total feeling, so it is too
\opage{121}with the contrast between understanding and ignorance. The humorist's standpoint
is aptly stated in the expression docta ignorantia, as this was set up in the early
Renaissance by Nicolaus Cusanus.\footnote{\textit{Den nyere Filosofis Historie} I,
p.~81.} The humorist's stage is conditioned intellectually by the relation of contrast
between the rich fullness of existence and the full understanding available at any given
time. Neither of these opposed terms can fall away without the whole stage being changed.
If one throws out the rational shaping because it cannot hold everything, one ends in
mysticism, which in turn quickly falls back into dogmatics.\footnote{When Nicolaus Cusanus
had written his remarkable book \textit{De docta ignorantia} (1440), an interesting
discussion arose as to whether one needed to keep up the insight into the conditions of
knowledge once one had seen that knowledge was limited. Some thought that the consequence
must be pure mysticism. \textit{(Autour de la docte ignorance.} Par
\textit{Vansteenberghe.} Münster 1915.)} % the Danish prints „[1440)“

35.~In a similar way the humorist stands before the fact that people, even the most
acutely thinking people, differ so greatly from one another in their view of the world.
History shows that different sciences have at different times laid claim to being the only
or the highest science. Thus for Hobbes physics was all in all, for Moleschott chemistry,
for Hegel logic, for Bergson vitalistic biology, while for Leibniz and Lotze a concluding
view of the world could be reached only through the application of psychological analogies
to existence as a whole. The very conflicts that arise among the different attempts to
interpret existence from one of its single sides the humorist takes up into his view of the
world as a testimony to how we human beings are placed, once and for all, before existence.
He can find in this, with Sibbern, an example of all development, that of thought too,
proceeding sporadically. He sympathizes with every thorough and honest striving to make a
general viewpoint prevail, all the while he holds fast the greatness and infinity of the
task. This contrast conditions, as already said, an essential side of his total feeling.

I once tried, for my own part, to describe \opage{122}a form of this total feeling. Since I
cannot give a better description, and since the one I gave at the time is moreover the text
I have sought to develop more fully in the present treatise, I set it down here (after
``Den menneskelige Tanke'' p.~377 f.).

``There is a standpoint on which, through life's mixture of tragedy and comedy, of victory
and defeat, a fundamental mood has been won that has at once the character of earnest ---
by virtue of the great that has been experienced and the pain it has often cost --- and of
irony toward every attempt to express in images what in its fullness and through its
contrasts bursts every form one tries to fix. The view of life or feeling of life that thus
arises holds fast the conviction of a great unity in existence, as regards values as well as
all other subjects and elements, and can therefore turn jestingly against the events of
life, because they are powerless against that certainty, as well as against the attempts
to say the unsayable. The jest nevertheless takes on the character of wistfulness, because
a contradiction is constantly felt in the great being so often tied to the paltry, and in
not succeeding in saying a concluding word about the core of life and the conditions of
life. Behind the irony and the wistfulness there will lie a great resignation, now more
practical, now more contemplative --- now bitter, now cheerful.'' ---

I have for my own part learned to regard the great humor as one of the ideal human types of
life, and it is the one I most sympathize with. But there can, after all, be both an unhappy
and a happy love. (Cf.\ § 26 and § 32). And both unhappy and happy love of something can
lead to understanding.

That the great humor, as it has been described here, should in and of itself be the highest
stage of life cannot, as said earlier, be maintained. The great practical striving on the
one side, tragic suffering on the other, can both be higher stages of life than humor. As
we shall see later (ch.~IX), the joy of knowledge too can develop into a total feeling and
become the central thing in the soul, and it will not be the least minds in whom this
happens. Something corresponding can also happen with the artistic need. Here, then, as with
practical striving, it is a great task that takes the mind captive and draws all energy to
\opage{123}itself. But the great tasks and the tragic fates do not fall to everyone's lot.
Even one who has worked in earnest, and who has had his sorrow, can still be clear that his
striving was as fluff compared with that of those who were placed in the great strife, and
that his sorrow was as play compared with that of those to whose lot the deepest sufferings
fell.

% ===== VII. ACTION AND HUMOR (pp. 124--134) =====
\chapter*{VII.\quad Action and Humor}
\addcontentsline{toc}{chapter}{VII. Action and Humor}
\markboth{Action and Humor}{Action and Humor}
\opage{124}36.~Only purely theoretically can a distinction be made between feeling and
will. Every feeling of pleasure or pain is connected with a need, an involuntary movement in
one direction or another. When the need is joined with the ability to reach a suitable
result, we call this connection of need and ability instinct. The need becomes a drive as
soon as there is joined with the feeling of pleasure or pain a representation of what can
preserve and increase the feeling of pleasure or avert and lessen the feeling of pain. The
active side of a state of feeling is not something that comes in only after the feeling has
arisen, but is tied to it from the first. The question from the first is not how need or
drive can arise from feeling, but how they can be checked so that the feeling can lose its
active character, at least so that it does not lead to immediate action. Only in the later
development of the life of feeling does a contrast perhaps arise between resting and active
feelings.

What holds of all feeling holds also of a feeling like humor. The need and drive that make
themselves known in it will be able to go in two directions. There will stir a need or drive
to immerse oneself in the great, to maintain it, to fight for it; but there will also stir
a striving to secure even the least its place in the world, perhaps even to hold one's hand
over stupidity and wickedness, if they do belong and ought to belong --- by virtue of the
principle felix culpa, or because evil itself is thoroughly removed only by the way of free
development.

There lies, as we have also seen (§ 8, end), an ethics in germ in every total feeling, and
this comes from there lying \opage{125}a willing (the word taken in the widest sense) behind
every total feeling, a willing that leads to valuation and action. The earnest in humor
rests, like all earnest, on a value, an end, standing as guiding and determining over against
the single moods and the single actions.

Many psychologists trust to such a degree the schematic survey of the qualities
(``elements'') of mental states that a provisional analysis can furnish that they regard the
concrete, actually given and very complicated mental states only as mechanical connections
of such elements. Just as, on a frequent popular conception, the elements in nature exist
first, perhaps ``in the beginning,'' and then the various compound substances arise later,
so one would build up the mental world from schematically set up mental elements. One
confuses real coming-to-be with the presentation in a textbook, where one begins with the
elements. It cannot be denied that in the mental domain sensations, representations and
feelings, in their abstract and isolated form, are the simplest we can think of. Ordinary
usage has already made the analysis that leads to setting them up. A childish psychology now
operates with these abstractions. But they have been torn out of the connection with need and
striving in which they originally occurred. It is only a special development, a
differentiation owed to definite conditions, that leads to a certain isolation of the
different sides or qualities of the mental life. Just as pleasure or pain are from the first
tied to need and drive, our sensations are from the first conditions for the release of
reflex and instinctive actions, and our representations from the first closely tied to
expectation and drive. Even scientific research begins as a striving to find means to
definite ends; only later can the thoughts thereby called forth acquire immediate interest.
It is not action that is a special case of sensation, representation and feeling. It is,
conversely, these that, in the isolated character psychological analysis gives them, rest on
a checking of the involuntary urge to action. Even after differentiation has set in, all
elements of cognition and feeling keep a relation to action. They keep the stamp
\opage{126}of being transitional members within a striving in a certain direction, or of
being called forth by such a striving.

Genetic psychology, which seeks to follow the mental life in its actual unfolding, corrects
in a useful way the misunderstandings that can arise from confusing the elements found by
analysis of the differentiated mental life with the forms of the beginning mental life. What
at a later stage we can conceive as separated and as to a certain degree independent of need
and striving we have no right without more ado to read into all the stages of the mental
life. A psychological antinomy asserts itself here, owed to the fact that reflection comes
into play only after the mental life has long been in luxuriant unfolding; therefore it has
difficulty understanding its own past and believes that the differences it can now find have
been present at all stages of the mental life whose late fruit it is itself.\footnote{See my
\textit{Psykologi} ch.~IV and \textit{Den menneskelige Tanke,} p.~10--12. --- The
significance of the genetic viewpoint over against the purely analytic has recently been
stressed emphatically by \textit{Edward Thorndike (Animal Intelligence.} 1911. ---
\textit{Elements of Psychology} 1905). See above § 3.} --- One could very well present the
whole psychology of sensation, representation and feeling as chapters within the psychology
of the will. Even logic one could find a place for there, namely in the chapter on
deliberation.

37.~When we now return to humor, it must be admitted that it can pass over into a purely
contemplative state, an aestheticizing, in which the mind revels in thoughts and images
while the personal basis (see ch.~IV), especially the element of sympathy, is blurred. It
can then happen that one arbitrarily tumbles the objects about and tentatively places them in
different combinations so as to enjoy one's own power and one's own whim. Or one looks at
reality as something distant, as through a reversed opera glass. Or one merely delights in
the play of forms and colors, indifferent to the real existences that lie behind. Thus the
young man in Ibsen's ``Paa Vidderne'' (``On the Heights'') looks from the mountain waste at
what goes on down in the valley; the real subjects of the picture do not concern him,
although it is his own \opage{127}mother's house that is burning and his own sweetheart who
is riding to church as another man's bride:

\begin{quote}
All this together I stood and saw\\
from up on the heights of life;\\
a higher light lay over the sight ---\\
but see, none can understand that now\\
who sits down below in the swarm.
\end{quote}
% Danish: Det Alt tilhobe jeg stod og saa / oppe fra Livets Vidder; / et højere Lys
% over Synet laa --- / men se, det kan nu Ingen forstaa, / som nede i Sværmen sidder.

Something of this kind Goethe probably had in mind when he called humor heartlessness, and
when he also thought that humor has no law or hold in itself and therefore easily degenerates
into arbitrariness. (Letter to Zelter. 30 Oct.\ 1808.)\footnote{In \textit{Kierkegaard's
Papirer VI} p.~38 there is a piece headed: ``Replik af en humoristisk Individualitet''
(``Rejoinder of a humorous individuality''). What he depicts here is a blasé indifferentism
that does not agree with the way in which he, at about the same time, characterized humor in
``Uvidenskabelig Efterskrift.''} He is thinking of many Romantics' capricious play with
everything between heaven and earth, which left no time for earnest artistic working out.
And yet, he says, ``the highest, indeed the only operation of nature and of art is shaping
(die Gestaltung), and within form again individualization (die Spezifikation), which leads
to every thing becoming and remaining something peculiar and significant.'' --- Rather one
ought surely to say that when humor does not lead artistically to the shaping of individual
figures, it is either because the humorist is not an artist or because he is not in reality
a humorist (in the sense in which Goethe himself takes the word when he speaks of Sterne).
The great humor, as it has been described in what precedes, will, where there is artistic
ability, lead precisely to busying oneself with interest with the single until it stands
fully marked in its peculiarity, the peculiarity that precisely gives it worth and determines
its place within the great whole. The great humor cannot be thought without a double
interest, an interest both in the whole and in the single.

38.~But does there not lie a danger at all in a total feeling like humor, in which so many
different, in part mutually \opage{128}opposed elements are united? Will not such a union of
contrasts easily lead to a leveling or evening out that weakens --- to a middle way being
taken, so that the fruitful energy the contrasts can release precisely in their extreme form
falls away? --- This remark is owed to a young friend with whom I have often discussed the
subject of this treatise. I believe that the coming-to-be of every total feeling can bring a
danger, in that its elements are joined so that they are no longer released each by itself.
It is the timbre of the fusion or the law of the organization that now determines how one
reacts to a new experience. There is, after all, in general a danger in every result
attained, every concluded state, whether it bears more the stamp of feeling or of cognition.
Instead of being regarded as a standpoint for new orientation and new work, such a state
easily becomes a conclusion of life and activity. But in and of itself it does not lie in
the concept of total feeling that it must be a definitive conclusion. It can very well mean
only a station on the way, a practical reckoning that is necessary to be able to continue
life so that the experiences of the past can benefit the future. But there is in addition a
circumstance of general psychological interest that is of significance in this connection.

A total feeling has essentially the character of a disposition, not of an emotion. It is
owed, after all, to a gathering together of a series of experiences that each by itself
have called forth more or less deep-going emotions and have thereby left lasting traces in
the mind. And it marks a firm direction that the life of feeling has now taken, even if
oscillations or ripples can arise. And when movement in this direction is furthered or
checked, emotions can constantly be released that are then often of a wholly different
nature from the disposition (cf.\ § 2 and § 7). Love for a cause one has made one's own can
give rise to combativeness where it must be defended. And the fact that this special emotion
has been released in the service of the cause can then in turn react on the total feeling
toward the cause itself, change its timbre or bring a new element into the whole
organization. Instead of the quiet occupation with the object of feeling \opage{129}there
can now appear a mood of armed watchfulness toward what may threaten. But it may also be
that when the situation that called forth the emotion has ceased, the mind sinks to rest in
its old bed.

In our conception of people there is a distinction to be made between the character as a
whole and the particular way in which it shows itself in a definite situation. An older
aesthetician (Hurd, quoted in Lessing's ``Hamburgische Dramaturgie'' No.~73) already
remarked that even Shakespeare's most strongly drawn characters do not express their
essential or ruling qualities in every utterance or action, but only when the situation
gives a natural occasion for it. Yet I would rather turn the proposition round and say that
when the situation gives occasion for it, actions or utterances can come forth that cannot
directly be said to lie in the ruling qualities, although they are not incompatible with
them. This bears witness to Shakespeare's deep understanding of human nature. Every
character is like an orchestra that is interrupted by the solos only in single episodes.
The disposition is an orchestra, the emotion a solo. One does the dramatic poets wrong if
one overlooks this difference. Aristotle already (Poetics, ch.~15) commits such a wrong when
he finds an inconsistency in Euripides' Iphigenia in her appearing as another person when
she begs her father for her life than when she later goes willingly to her death. The
explanation is, however, quite simple. It has, after all, dawned on her in the meantime
that her sacrificial death is necessary for the sake of the Greek army, and that her refusal
would expose the chivalrous Achilles to destruction.\footnote{See \textit{Gertz'}
translation of \textit{Ifigeneia i Avlis}. The note to v.~1368.} A similar misunderstanding
criticism I have also heard against Euripides' Medea, because she lamented in private, but
appears with royal dignity before the women of the city. Most people do, after all, give way
more to their emotions in private than when they appear before other people. This
difference need not have anything to do with character at all, and it contains no
contradiction. A similar criticism is directed against Sophocles' Antigone, who defies the
king's prohibition against \opage{130}burying her brother and thereby shows strength and
courage, but later, when she is to suffer death, laments her early death. There is,
however, no contradiction in her now, after having done what she regarded as her duty and
thereby shown a bold defiance, appearing as the young Greek girl she in reality is. There is
no contradiction in the orchestra sounding again after the solo motivated by a definite
situation has ended.

Only experience shows that in the more harmonious times of the disposition it is not always
easy to understand what one performed involuntarily in a sudden emotion, when need and
ability were awakened at once. The moments of action present the greatest concentration the
mental life commands; here is the most intimate gathering of all the powers of the mind in
one single direction. In the moment of decision the concentration can rise to ecstasy. No
wonder, then, that one later does not understand the zeal, the enthusiasm, or the hotness,
or the strength one showed in such moments. One can in memory grow dizzy at what did not
cause dizziness during the experience or the performance. In the moment of action emotions
were to the fore that stood in contrast to the otherwise ruling disposition and for a while
laid claim to all energy.

An eager mountaineer told Ludvig Feilberg (Samlede Skrifter p.~736) of a dangerous descent
from a high, icy slope. ``When I afterwards,'' he added, ``thought back on the time when,
pushing myself carefully downward on my backside, I tried to get near enough to the edge of
the abyss to be able to look down, a feeling of terror came over me. It was really
remarkable, for in the moment itself there was no question of any terror\dots{} One had
other things to do.'' --- The anthropologist G.~H. Schneider reports a similar trait from
his own experience. On a botanical excursion in Crete he was attacked by some armed shepherds
and could save himself only by a breakneck flight down a steep slope. ``From block to block
I sprang, or rather plunged, often in dangerous leaps, in mortal terror down the mountain,
and so reached the nearest village\dots{} Later I could only be astonished at my unerring
sureness in the most dangerous different leaps from stone to stone, following quickly on one
another, and yet there \opage{131}could be no question of deliberation.'' (G.~H. Schneider:
Der thierische Wille, p.~203.) --- A delightful and at the same time touching example occurs
in Goldschmidt's ``Min Onkels Tømmerplads'' (``My Uncle's Timber Yard''). The old merchant
was in suspense and fear for one of his ships, whose good captain he was especially fond of.
Suddenly the shop door opens and the captain walks in. ``Uncle looked up --- the next moment
he had vaulted over the counter and was round the neck of the newcomer.'' When the emotion had
subsided, the old man was seen measuring with a thoughtful look the height and breadth of the
counter he had so unawares vaulted over.

One cannot, then, infer from the apparent calm of the total feeling the absence of energy in
momentary situations. When the strength of a feeling is spoken of, one must remember that it
can express itself in two ways, in the form of inwardness in the disposition and in the form
of vehemence in the emotion, and that these two ways can stand in interaction with each
other.

If from the world of memory it can be hard to understand the energy with which one acted in
the past, it is often still harder, from the present state, to think oneself into the state
a future action demands. Here the distance between the greatness of the task and what one
believes oneself able to accomplish can stand as so great that it cannot be evened out; every
possibility of approximately solving it pales. Amiel's journals (Journal intime) have here
their constant, again and again varied theme. % the Danish prints „stadige.“
``The ideal,'' he says (I, p.~22; 102), ``poisons every imperfect possession for me\dots{} For
me it is all or nothing. Irony, on account of a comparison with the glimpsed or dreamed ideal,
takes neither itself nor reality seriously.'' It was Amiel's misfortune that he could not
become a humorist, could not come to see the small --- also the small he himself could
accomplish --- not only in its contrast to the great but also in its connection with it.
Hence his want of will, his distrust of himself; he could not take himself seriously; at most
he could feel pity for himself.\footnote{Je n'aime que le sérieux, et je ne puis prendre au
sérieux mes circonstances ni moi-même\dots{} je me prends perpétuellement en pitié au nom de
tout ce qui est beau et admirable. \textit{Journal intime.} Genève 1908. I. p.~239.} % the Danish prints „moi-rmême“
\opage{132}Humor joins the greatness of the ideal and the paltriness of action, irony at one's
own limitation and admiration for the perfect. It maintains the positive significance of the
widow's mite. It makes strenuous work possible, all the while it gives room for the thought
that, when the highest measure is applied, we are nevertheless unprofitable servants. Irony
can here enter the service of humor, in that in the moment of action one takes oneself
seriously, although one knows very well that even the greatest work is only jest when the
ideal is set high enough; --- behind this jest there lies an earnest that is precisely the one
that maintains the ideal measure. (Cf.\ § 18).

The disposition in which humor consists is present, scattered and indirect, in work when this
is done without self-importance, without sentimentality, without oppressive melancholy,
without garrulous dwelling. The presupposition is that even if the ideal that stood before us
is not fully realized, and even if in the course of life much that is valuable vanishes, the
spring of value nevertheless constantly flows in secret, and it is only a matter of digging
deep enough to make it bubble forth. In the very mind that confidently sets an imperfect
action into the world of changes and illusions there is greater worth than in the soft
dreamer's spirited reflections and fantasies. Just as the tragic hero shows his loftiness
precisely in his downfall, so he who is not deterred by the imperfection in everything finite
from giving what he can shows that there is loftiness in his mind.

39.~In his profound work ``Sygdommen til Døden'' (Skrifter XI), which bears witness to
penetrating psychological experience, the finest of all he produced,\footnote{On this work see
\textit{Danske Filosofer}. p.~163.} Kierkegaard has depicted the relations of contrast we
have been occupied with in this chapter. I name his expressions for the most important
contrasts and the relation of humor to them.

``The despair of infinitude, which lacks finitude'' arises when one will not put up with the
limitation to which all realization \opage{133}of something great is subject. ``The despair
of possibility, which lacks necessity'' arises when one is caught in empty wishes and hopes,
or else in fear and melancholy. From these forms of despair the great humor frees. And the
same holds against the two opposite forms --- ``finitude's lack of infinitude'' and
``necessity's lack of possibility.'' For humor points out toward wider horizons and thereby
toward more possibilities than the narrow view, for which all sounds so easily seem to be
closed, can discover.

However profoundly Kierkegaard depicted those contrasts in their extreme forms, he would
nevertheless definitely have denied that humor could ever remedy them. In his youthful work on
the concept of irony (XIII p.~379) he calls Solger's standpoint, which essentially coincides
with what I call the great humor, a kind of contemplative devotion and thinks that it leads
only to living in an unreal reality, and that it stresses only the negative side of the great
and does not reach up to a higher reality, which in his view is reached only by adopting
theological presuppositions. And in his main philosophical work (``Videnskabelig
Efterskrift'') (VI.\ p.~525) he regards humor, which in ``Stadier paa Livets Vej'' (VI.\
p.~393) he had nevertheless designated as an expression of strength of spirit, as a revoking, a
withdrawing from existence, a backward perspective, a standing still, a life in the world of
memory.

This does not agree well, however, with the ``stage'' he designates as humor lying later in the
series than the ethical ``stage.'' According to Kierkegaard a following stage is not to deny
the preceding one but to show its superiority by having taken up into itself everything true
and justified in it. Humor would not be a higher stage than the ethical if it denied this. It
would then itself fall under an ethical judgment, instead of precisely showing richer
conditions for the ethical, as the great humor described in what precedes does. But the secret
is that all the ``stages'' that lay between the stage of enjoyment (the ``aesthetic'') and the
strictly religious lost more and more of their significance for Kierkegaard. ``Humor'' too was
squeezed in and crushed to death between the gospel of enjoyment and the gospel of suffering,
and \opage{134}Kierkegaard really took back what he had described as a high, purely human way
of taking life. He made fine contributions to the psychology of the great humor, but in the end
he could not let it come into its own as a type of life grown up on the ground of experience of
life.

% ===== VIII. HISTORICAL CONDITIONS FOR HUMOR (pp. 135--144) =====
\chapter*{VIII.\quad Historical Conditions for Humor}
\addcontentsline{toc}{chapter}{VIII. Historical Conditions for Humor}
\markboth{Historical Conditions for Humor}{Historical Conditions for Humor}
\opage{135}40.~The psychical process by which a total feeling comes to be cannot go on under
all circumstances. Every total feeling is conditioned by social and historical circumstances,
takes on a historical character, and it becomes the object of a special investigation whether
the conditions for its arising are present at a certain given time. It is one of the questions
in psychology where it must use the historical (sociological, comparative) method. Neither
self-observation, experiment nor analysis can be decisive here. It goes with humor here as it
goes with other total feelings, for example national feeling, the feeling of duty, religious
feeling, the feeling for natural beauty, etc. One sometimes hears the opinion that it is harmful
to science to admit that such feelings in fact exist, because a certain validity is thereby
admitted for what is their object, or for the representations to which they are tied. But in
the acknowledgment that a feeling exists there lies absolutely nothing as regards the validity
or the worth of its object. And if it proves to be tied to definite historical conditions, it
is clear that if these conditions fall away, the time of the feeling will also be over.
According to the view of the religious problem I developed in my Philosophy of Religion (and
later in the concluding section of ``Den menneskelige Tanke''), the decisive question in the
philosophy of religion is precisely whether there will always be a possibility of religious
need and feeling after the objective, scientific validity of religious representations has been
given up. The psychological-historical method here takes the place of the speculative or
metaphysical. And this method must also be used as regards humor.

\opage{136}Single feelings of joy or sorrow, of the exalted or the ludicrous, of the tragic or
the comic may be possible without the peculiar complication becoming possible by which a total
feeling like humor develops from the interplay of these feelings. This has at any rate, as we
have seen (§ 10), first been observed and described in modern times, especially in the last
centuries. One has, of course, no right to infer from this without more ado that it did not
exist before. There is a difference between the condition for the existence of a mental state
and the condition for its being observed and described. But for that very reason we have no
right either to assert that what we can now observe or see described has also been possible at
earlier stages of development. There are, for example, differences and distinctions that we
easily make after development has reached a certain stage, but which cannot without more ado be
applied to earlier stages. And something corresponding holds for total states, which rest on an
interplay of manifold elements and conditions that are not always present.

Psychologists and aestheticians are fairly agreed that humor is a modern phenomenon. Cervantes
and Shakespeare are often named as the first writers in whose works this total feeling is
traced as the fundamental mood. Antiquity's conception of characters and fates was simpler;
one kept to great and pure lines, and very composite personalities and events are not
presented. The ancients certainly had a lively feeling of the law of limitation to which
everything human, indeed everything in the world, is subject. By virtue of a ``divine envy,''
which corresponds to what modern Romanticism has called ``world irony,'' great fortune and
sudden downfall are closely bound together. To that extent there was here a possibility of a
total feeling into which opposed moods could enter as elements. But the Greeks seem to have
stopped at the successive relation between these contrasts of light and dark, good fortune and
ill fortune, without the opposed terms coming to assert themselves at once and thereby being
able to enter as elements into a total feeling (cf.\ § 6).

Plato speaks in the dialogue ``Philebus'' of several mixed feelings, without the way in which
they have arisen being stated more closely. Not only as feelings at a drama, but also in ``the
\opage{137}whole tragedy and comedy of life'' do we experience such mixed feelings (p.~50 B).
Tragedy and comedy stood for Plato (in one period of the history of his thought) in close
connection. But unfortunately all that is said in the famous ending of the ``Symposium'' is that
it is the business of the same man to write comedy and tragedy.
In Plato one might perhaps presuppose the great humor when one sees how he, precisely in the
``Symposium,'' sets the contrasts of life together and through jest and irony lets the
greatest thoughts unfold. But it did not come to full development as a total feeling, partly
no doubt because of the pessimism that marked his last years, partly, and perhaps especially,
because he was most taken up with developing his thoughts in their systematic form. To this
last point we shall return in the next chapter.

That Plato did not lack an eye for how a total state can develop through the alternation of
different feelings is seen, as discussed above (§ 2), from his description of how firmness of
character (``courage'') comes to be. But for the total feeling we call humor he had no eye. He
did not stand in the living relation to life that his great master did. His psychology was
not able to present the total state that lay behind the latter's life work, and which is
precisely the greatest example of the great humor. This does not conflict with the
proposition that while irony is an ancient feeling, humor is a modern feeling. For Socrates
stands precisely at a breakthrough where the ancient points beyond itself, and it is precisely
thereby that he can stand as the representative not merely of a single time but of all human
striving in general. I return to this point too in the next chapter, but remark the following
for the present.

Hegel and Kierkegaard showed that irony was the natural form for the breakthrough in antiquity
when personality for the first time sets itself in contrast to the given, traditional order
of things and strives to become a little world for itself. For the present it is only a matter
of maintaining the right to be oneself, to think one's own thoughts, to take life in one's own
way. But as long as a new comprehensive world of thought and a new social order cannot yet be
developed, it is natural that \opage{138}the single person apparently goes along with the
traditional and tries to take the standpoint of others, although what is given in the
existing thoughts and arrangements cannot in earnest express his real standpoint. He fights
for a formal right; such a right is what the principle of personality expresses (see § 16);
--- but a real content is lacking for the present. Often, then, jest will poke its head out
behind the earnest in a teasing way. What one overlooks when one regards Socrates merely as an
ironist is that the jest that lies behind the apparent earnest can itself in turn be the
expression of a deeper-lying earnest (see § 18). The new thoughts and arrangements that could
replace the given culture could for the present express themselves only in a purely subjective
way, through the personality's return to itself and through its liberation from what had
been shaped so far. Everything objective begins as subjective, in the form of subjectivity.
Therefore a man like Socrates is not only the great ironist but (without knowing it) the great
herald. The jest that lies behind the outward earnest is to that extent not only the jest of
release (§ 13), but stands in the service of an earnest that cannot yet be shaped in a positive
way. Irony becomes an expression of humor. The more the subjective basis of irony is deepened,
the more it becomes humor. In a certain sense humor can be said to be more subjective than
irony, because it presupposes a deeper layer in the personality than irony in itself needs to.
Thus Schopenhauer declared that irony is objective, humor subjective, and he finds in this the
explanation of the masterpiece of irony having appeared in antiquity, while humor is a modern
feeling. (``Welt als Wille und Vorstellung.'' II. ch.~8). But if it is the task of the
principle of personality to penetrate both the life of thought --- by laying weight on all
knowledge and all conviction being won by one's own work --- and social life --- by
maintaining the freedom of the single person everywhere it does not infringe the freedom of
others --- then there is heralded in the great humor an objectivity of higher rank than the
one irony seeks to free itself from. The new world that lay as a germ in Socrates' principle
of personality needed millennia to unfold, and it is still fighting for its life.

\opage{139}41.~Total feelings are not only more composite but also more individually
nuanced than single feelings. For the single feelings that enter as elements into a total
feeling can stand in the most different relations to one another within this new whole, and
significant differences in the timbre of the total feeling in different individuals are
thereby conditioned. Differences between individuals as regards the life of feeling rest
rather on such different timbres than on any single feeling being wholly lacking in any
person. Conversely, of course, the very possibility of a total feeling arising will be
conditioned in high degree by the original individual peculiarities.

The different temperaments will not, as regards humor, offer equally good soil for its
arising. There is, as I believe,\footnote{See my \textit{Psykologi} VII C, 1. --- Cf.\ also
\textit{Alfr.\ Lehmann: Hauptgesetze des menschl.\ Gefühlslebens.} p.~390.} reason to
distinguish eight main temperaments, which are determined partly by the speed and strength
with which one is affected and reacts, partly by the predominance of the feeling of pleasure
or of pain. Two of these temperaments will furnish good soil for humor, namely the dark, strong
and slow (the ``melancholic'') and the dark, weak and slow (which has no current name), since
they from the first, each in its way, let the dark sides of life come to play a prominent
part. The bright temperaments, in both their strong and weak, quick and slow forms, will not
so easily let the great shadows appear in the picture of the landscape of life that forms
involuntarily in consciousness. And the other dark temperaments, both the strong and quick
(the ``choleric'') and the weak and quick (which is anonymous), will not easily make possible
such a lingering state as humor presupposes; ``quickness'' means, after all, an inclination
to pass over into action or into other states. --- Together with temperament there work the
experiences that are had and the reflection that is exercised.

Of course the dark temperaments named can make humor impossible when they do not let
experiences come fully into their own. In humor, after all, earnest must be able to have jest
as its form of expression, and must be able to be united with mildness and understanding of
\opage{140}what life brings, even if it does not without more ado fit the ruling fundamental
mood. The bright sides of life must also have their effect.

Of the two opposed moods of life, optimism and pessimism, optimism nevertheless seems on the
whole to be a greater hindrance to the arising of the great humor than pessimism. Optimism has
an inclination to move on the surface, both as regards thought and feeling. It arranges things
for itself in the easiest possible way and does not take the theoretical and practical
problems as thoroughly as the great humor presupposes. But pessimism too can make humor
impossible, especially by its disregard of experience.

In his characterization of George Eliot Leslie Stephen remarks, what every reader of her works
will confirm, that she really appears as a humorist only in her first books (Adam Bede, The
Mill on the Floss, Silas Marner). Later her increasingly melancholy mood made the bright sides
of life be pushed back in her depictions. But the reason lay also, and in good part, in her
having built in her earlier works on her own experiences and memories, while later she fell
into a direct constructing of persons and events. --- Here the interaction between
temperament and experience comes out plainly. ---

More important than temperament, however, is perhaps the difference mentioned earlier between
the once-born and the twice-born (see § 2) --- between natures that develop in a continuous way
and natures whose development leads through breaks and leaps. The twice-born will experience
the opposed sides of life more strongly, because their experiences are successive, not
simultaneous. All else being equal, successive contrasts work more strongly than simultaneous
ones that come as parts of the same experience. The twice-born have the possibility of
experiencing stronger contrasts than the once-born as a rule experience. It can even be very
hard for such twice-born, whose life is split into periods between which even the bond of
memory breaks, to reach fusion or organization of the highly opposed feelings that
successively rule the mind. But if it nevertheless happens that the connection succeeds, the
twice-born will become \opage{141}humorists in the highest style. In the once-born the
contrasts do not so easily penetrate into the innermost mental life; their development has
the character of expansion, of self-unfolding. Yet experiences can give even the character
most harmonious in itself such an insight into the depths of life, into its involved being,
that it gets the conditions for a total feeling like the great humor. If Shakespeare, as much
may suggest, belonged to the once-born, he may have got the conditions for his humor in this
way.

That humor is especially a modern feeling is connected with its marked individual character,
which in turn is a consequence of its composite character, which offers the possibility of a
host of different timbres. Ethical and religious feeling are also total feelings, but they
have a support, a support common to all the individuals of a social group, in the norms of
society and in the forms of cult, so that individual nuances do not so easily come forth, or
at any rate are not so easily discovered --- and if they are discovered, often stand as
something to be pushed out. Religious cult has for millennia been an important basis for total
feelings. Through it the single individuals are led into alternating moods of fear and hope,
sorrow and joy, suffering and glory, which are determined point by point by objective, vivid
images and traditions. ``Good Friday was a bitter day, but fair was Easter morning.'' This
motto fits many forms of Oriental and Greek cult as well as the greatest festival of the
Christian Church. The individual life of mood here moves in prescribed rhythms. Humor, on the
other hand, is a result of life that can be reached only in a way peculiar to each single
individual, since it depends on the individual's relation to the experiences that fall to its
lot, and perhaps fall to no one else's lot. In the ethical and religious domain too
individualization comes more and more forward in modern times, with the consequence of a more
complicated character of ethical and religious problems.

It is not quite right, what Jean Paul says (Vorschule der Aesthetik. § 76), that there is an
earnest for all, but a humor only for a few. For earnest too is more and more individualized.
And the greatest collisions take place where earnest people do not understand one another. ---

\opage{142}Humor thus presupposes a complication, a subjectivization and an individualization
of the life of feeling that belong essentially to modern times, although Socrates stands as
the great pioneer of this whole development. Like all great minds he is at once a great
beginner and a symptom of what was beginning, and of what was later to be developed through
manifold historical and social changes. On one significant change in the intellectual domain
we must still dwell a little.

42.~Greek spiritual life developed under a limited horizon, as regards both the knowledge of
nature and historical experience; the contrast between the great connection of existence and
the limited human ideals and human fates could not assert itself in earnest. Existence stood
as a closed whole that could be conceived as the home of everything need and thought could
shape. The Middle Ages preserved the limited picture of the world. Greek thoughts were used,
more or less forcedly, in the intellectual shaping of the dogmas. And after the living,
ecstatic expectation of primitive Christian times had given way to more patient thoughts, it
became the task to fit the life of feeling under the frame formed by antiquity's picture of the
world and the dogmatically shaped doctrine of faith together. Within this frame love and
enthusiasm, fear and jubilation could unfold securely and richly. Sometimes, from the need of
the life of feeling, the foundation could be laid for a reshaping of cult and through it for
the shaping of new dogmas, as happened, for example, with the worship of the Madonna. But on
the whole the ecclesiastical authorities saw to it that the dogmatic limits were not
overstepped, not even by the boldest flights of mysticism. A total feeling grounded on
firsthand experience of life had no conditions for developing.

Here the Renaissance brought a change, partly by its discovery of natural man, partly by its
widening of the picture of the world. Here came new spiritual conditions of life. Human life in
its purely natural unfolding was seen against the background of an unsurveyable order of the
world. This contrast led partly to \opage{143}a bold self-assertion as the basis of all personal
development, partly to stressing, as a quality of character and a principle, the high mind that
shows itself in life with all its interests and tasks being seen as a vanishing episode within
an unsurveyable order of things. --- The word magnanimity (magnanimitas or sublimitas) was a
watchword in the literature of the Renaissance. The quality it designates stood as a natural
unfolding of the free strength with which one was now minded to take hold of life and its
tasks. In various forms this concept appeared in the moral philosophers of the earlier
Renaissance,\footnote{\textit{Fiorentino: Il Risorgimento Filosofico nel Quatrocento.} Cap.~IV.
--- \textit{A.~Herrlin: Renæssancens Etik.} (Printed as manuscript.)} and gained especially
far-reaching influence on the following period through Telesio's ethics, according to which
magnanimity stands as the highest form of self-assertion.\footnote{\textit{Den nyere Filosofis
Historie.} I. p.~96.} In Bruno, Descartes and Spinoza there then unites with the purely
psychological analysis of magnanimity the influence on the mind of the new picture of the world
in its unsurveyability. What in a nature like Pascal awakened terror, mankind's vanishing place
in the universe, awakened in the great thinkers of the Renaissance precisely a feeling of
nobility and dignity.

It was not only the greater extension in time and space that affected the mind. The widening of
the picture of the world was followed by an ever more thorough understanding of the lawfulness in
existence. The picture of the world showed more and more an inner connection among everything in
the world, the least as the highest, through general laws.\footnote{\textit{Den nyere Filosofis
Historie.} Second book.} The grain of dust and the globe move according to the same laws.
Nothing stands isolated and accidental. Even the most insignificant has its definite place in
existence, and this by virtue of the same law that assigns the great its prominent and
conspicuous place.

In the same direction as the conception of nature there soon worked also the historical
understanding that developed in the eighteenth and nineteenth centuries. Within human life,
which was an episode in the life of the world, the life of the single people, the single
generation and the single individual came in turn to stand as episodes. Thereby \opage{144}a
nearer background was formed, which had to work as a determining factor in the single person's
reckoning with his experiences of life. Here too the great and the paltry stood in inner
connection for the understanding eye, and from this viewpoint too the task was set of uniting
sympathy with the small and the sense of and reverence for the great.

A last intellectual condition was brought about by critical philosophy (epistemology). It showed
how all our science develops from the need to distinguish between imagination and reality, and
that this distinction can be carried through only by the use of the laws of human thought
itself. % the Danish reads "denne Skæbne" (this fate); read as "Skelnen" (distinction)
The picture of the world and the concept of reality are at any given time the work of human
thought. And neither that picture nor this concept can ever be concluded as long as the spring
of experience trickles and as long as there is life in thought. Reality is an ideal after which
we must constantly strive, and which makes every attempt at absolute conclusion illusory. (Cf.\
``Den menneskelige Tanke.'' IV B). But this must become of decisive significance for the total
feeling that forms as the result of experiences. The struggles and victories of thought, and the
conditions and limits that prove to be set for it, are facts of great significance in this
reckoning. (Cf.\ §§ 34--35).

--- If all the intellectual presuppositions indicated here can become of significance for the
formation of a total feeling like humor, it is understood that it --- quite apart from its
complex, subjective and individual character --- could develop in its whole fullness only late
and slowly, and that it must be able to appear in very different forms.

The intellectual conditions are not in themselves enough. (Cf.\ § 23). They will, as the
following chapter will show, acquire different significance according to the different
individual circumstances under which they appear.

% ===== IX. HUMOR AND PHILOSOPHY (pp. 145--175) =====
\chapter*{IX.\quad Humor and Philosophy}
\addcontentsline{toc}{chapter}{IX. Humor and Philosophy}
\markboth{Humor and Philosophy}{Humor and Philosophy}
\opage{145}43.~The great humor, as I have tried to describe it, is a view of life, or at any
rate a mood of life. It can draw nourishment from everything a person experiences, suffers and
enjoys, does and feels. Therefore a person's philosophy, if he has one, will also exert its
influence here. And there are, after all, as we have seen, intellectual conditions that can
assert themselves in the arising of this mood of life. But they are only contributing
conditions. And philosophy in and of itself will never be the only cause here. Philosophy can
lead to settling our spiritual accounts, clarifying the presuppositions that science tacitly
builds on and that life tacitly relies on. Just as little as philosophy can take the place of
all other science, just as little can it take the place of life. Philosophy must wish that the
other sciences unfold richly and vigorously, for thereby new tasks are constantly set in
epistemological respect; there is new work to be done in throwing light, by way of analysis, on
the nature and presuppositions of human knowledge. Philosophy must also wish that life be led
energetically and fully, for then there are new tasks for the one who wants to investigate the
basis on which people's conduct of life builds. Life and its phenomena are subjects for
philosophy; only indirectly and partly can they also be effects of philosophy. Philosophy is,
as Sibbern so finely and rightly maintained, only one element in the whole spiritual life, and
it exerts its influence in conflict with other powers in the world of thought and in the world
of life.\footnote{\textit{F.~C. Sibbern: Om Filosofiens Begreb.} 1843. (p.~82--86:
``Filosofien som Led og Moment i den hele aandelige Leven''). Cf.\ for my own part \textit{Den
menneskelige Tanke} (1910). p.~380 and my paper \textit{Aandelig Kultur} (Mindre Arbejder
III).} Personal experience of life and its \opage{146}yield do not without more ado coincide
with what philosophical thinking has found by its own path. There is still work to be done to
bring harmony between the two influences. It is not only in the development of the race that
the spiritual powers clash. This clash goes on within every wakeful personality, and in
different ways in different individuals. Here is a rich domain for psychological research, as
yet very little cultivated. It is the relation between the work of knowledge and the in part
wholly involuntary work that is done within the life of feeling and of will. The question is
whether what goes on in the two fields of work rests on the same presuppositions --- whether
the presuppositions on which the one field rests can be derived from the presuppositions of the
other --- or whether there is perhaps strife between the two groups of presuppositions.

Here we shall now investigate this point by the historical path, as regards the total feeling
we have been occupied with here. We shall investigate whether we find in great philosophers
that their course of thought and content of thought have been able to furnish a background for
humor as a total feeling. It will appear that philosophers have on the whole been most taken up
with shaping the intellectual background, whereas the basis the great humor presupposes has been
able to unfold most vigorously where the intellectual background became one with practical
experience of life. The history of philosophy has one great humorist to show, and he is a man
who belongs to life quite as much as to science. In Socrates life and thought became one as in
no other philosopher, and his mood of life cannot be better named than by the word humor, in the
sense in which it has been used in this treatise. Other philosophers are predominantly absorbed
in the work of thought and acquire only indirect significance for the history of the great
humor. Here it is seen precisely that philosophy is only one element in the whole connection of
life, and that its significance within this connection is conditioned by interaction with many
other elements. There is a division of labor that reaches right into the innermost regions of
the life of the spirit, and that makes a perfect interplay of all the elements of the inner life
so difficult and so rare. This explains why philosophy has only one great figure to set beside
the great \opage{147}poet who has given the most moving pictures of the contrasts of life, and
who has had the deepest understanding of how one must feel who lives through these contrasts and
is shaken through by them. Socrates and Shakespeare are the two greatest humorists. The one
concentrated his life on a single task; the other had the rich fullness of life as his basis.

This of course is not to say that philosophers on the whole have lacked a total feeling. But
their total feeling takes on essentially the character of the joy of knowledge, of intellectual
enthusiasm at finding points of view in thought for existence. Their total feeling is thus
different from the one we are investigating in this treatise. But this difference is so
characteristic that it may serve to illuminate our subject to dwell on it more closely.

44.~In my first year as a student I was present at a doctoral disputation where Sibbern
maintained with great emphasis the incorrectness of Kierkegaard's conception of Socrates as an
ironist. In contrast to it he laid down that Socrates was a humorist. The two things, according
to Sibbern, excluded each other. And against Kierkegaard Sibbern was undoubtedly right in this.
Kierkegaard himself, after all, especially in ``Uvidenskabelig Efterskrift,'' corrected the view
he had put forward in his treatise ``Om Begrebet Ironi med særligt Hensyn til Sokrates.'' (See
above § 16 and § 18). This correction is connected with the more detailed investigation,
mentioned earlier, of the relation between irony and humor in ``Uvidenskabelig Efterskrift.''
When Socrates is determined here as an ethicist with irony as his incognito, but out toward the
limit of the religious, this means, according to Kierkegaard's own doctrine of the ``stages,''
that he was a humorist, for irony lies before the ethical in the series of ``stages,'' and the
limit between the ethical and the religious is precisely the place where humor, according to
Kierkegaard, has its place in the series. (Skrifter VII p.~73 and 437). Already in ``Stadier paa
Livets Vej'' (VI p.~342) it is said that Socrates' earnest was hidden in jest. But earnest hidden
in jest is humor. --- In his ``Papirer'' (ed.\ by Heiberg and Kuhr, VI p.~118) Kierkegaard
corrects with special emphasis the view of the dissertation, \opage{148}which he explains by the
influence of Hegel. --- He could have added: and of theology.

But it was only against Kierkegaard's view (in the dissertation) that Sibbern was right in his
assertion that irony and humor exclude each other. Irony can, as we have seen (§ 18), very well
be a form of humor. And a historical investigation leads precisely to the result that this is
the most apt designation of Socrates' standpoint.

Socratic irony was essentially a method, the well-known analytic method. This consists in going
along with an assumption and going back from it to its presuppositions, which can then either
agree or conflict with other presuppositions or assumptions held by the one who maintains that
assumption. Thus Socrates goes along with apparent earnest with the assumption that the one he
is talking with really understands what he claims to know. By drawing all the consequences of
this assumption he then often comes to lay bare the ignorance or self-contradiction of the
participant in the conversation who was at first so confident. It looks as if it were a sport
Socrates was practicing here, and undoubtedly he had satisfaction in practicing the activity of
thought his method demanded. But for him it was not a matter of jest and play. The ironic
treatment was to incite to new, earnest and independent work of thought, and this was what
Socrates really wanted. There lay sympathy, ``erotics,'' behind the irony. The jest was steeped
in the most earnest interest in people. This shows itself especially in some statements in
Plato's ``Apology'' and ``Symposium,'' which according to recent investigations\footnote{%
\textit{Heinrich Maier: Sokrates.} Sein Werk und seine geschichtliche Stellung, Tübingen 1913.
p.~106 ff; 137 ff.} are to be regarded as biographical of Socrates. In the ``Apology'' (p.~29 C)
Socrates says: As long as I am able, I will not cease to strive for truth (``to
philosophize''), and whoever it is I meet, I will say to him: are you not ashamed to strive to
get as much money, as much reputation and honor as possible, while you do not care for insight
and truth and for your soul becoming as good as possible? And Alcibiades testifies (in the
``Symposium'' p.~215--216) how captivated and enthralled he was in conversation with
\opage{149}Socrates. His heart pounded, his tears streamed, and he felt the greatest
dissatisfaction with himself. Socrates was the only person before whom he felt ashamed. And when
Socrates opened himself to him in earnest, he had discovered through the outer, satyr-like shell
a divine interior.

If these utterances have a historical basis, irony was an outer form for a great and glorious
content, which was not, to be sure, expressed in definite and positive form, but at whose ground
lay a high concept of truth, of human dignity and of the significance of truth for human dignity.
There was an ideal background, what we can now call the principle of personality, in relation to
which Socrates' whole conduct among people must be seen. His jest and his sport of thought always
had earnest as their basis. It is not right, with Lipps,\footnote{\textit{Komik und Humor.}
p.~238.} to apply to Socrates, in connection with Aristophanes' ``Clouds,'' the saying that he
laughs best who laughs last. For at the last Socrates did not laugh, and no humorist does.
Socrates regarded everything else as slight in relation to a self-examination that could bear
life. It was a matter of becoming clear about one's real need and one's real ability, so that
one should not get onto paths that did not agree with one's personality. Socrates is the prophet
of the great striving for truth and of the great care of the self, both in inseparable union.
Science and life can both with equal right appeal to him. Already in antiquity the expression
was used of him that he was the first whose thinking concerned the conduct of life
itself.\footnote{\textit{Diogenes fra Laërte}. II, 20 (πρῶτος περὶ βίου διελέχθη).} He had high
thoughts of what it was to be a human being, especially to be a Greek, most especially to be an
Athenian. And in relation to this he found the usual goods to be of slight worth. In his own clear
thought man was to find the good that conditioned all other goods, and without which these are
in reality not indubitable and secure goods.\footnote{Even in \textit{Xenophon,} who otherwise
gives so sober and superficial a picture of Socrates, this thought comes forward with
philosophical sureness in \textit{Memorabilia} IV, 6 (the conversation with Euthydemus). But
\textit{Heinrich Maier} (\textit{Sokrates.} p.~57--62) has indeed shown that the whole part of
the Memorabilia to which this conversation belongs is owed to borrowing from or imitation of
Platonic dialogues.}

\opage{150}Through the Socratic connection of irony and humor one best understands Socrates'
relation to religion. He went along with the representations of popular belief and found in them
an expression of his trust that his life's work was not in vain, without, for the rest, laying
weight on shaping or grounding this faith more closely. He kept only to this: that ``no evil can
befall a good man, either in life or in death'' (``Apology'' p.~41). This faith sprang from his
faith in the validity of ethical striving. From the reasons with which he (in Plato's
``Phaedrus'') rejects rationalistic reinterpretations of myths, one can probably form a
representation of how he stood toward the representations of popular belief. He does not reject
such reinterpretations out of reverence for the gods, but partly because, if he began on such
things, so many monsters would stream in and demand reinterpretation that it would become
unmanageable, partly (and this is the main thing) because he still has enough to do to find his
way in himself and decide whether harmony or chaotic passion is the essential thing in his nature.
Here we have jest and earnest, irony and humor at once. Before the great, earnest task of
self-knowledge he jested with all theologizing. Plato's ``Euthyphro,'' where the dependence of
ethics on belief in the gods is disputed, can be taken as testimony that it was not only toward
rationalistic theology that he, in trust in ``the images of the gods within him'' (to use
Alcibiades' expression), took an ironic attitude.

Toward the belief in immortality he took, if one can believe Plato's ``Apology,'' a particularly
humorous attitude. He leaves it undecided there whether death is a sleep or whether it is a
passage to another life, to be led in the company of gods. It is plainly enough the first
possibility that lay nearest to him.\footnote{Cf.\ \textit{Heinrich Maier: Sokrates.} p.~435 f.
--- In contrast to this \textit{Burnet (Greek Philosophy from Thales to Platon.} London 1914.
p.~182) thinks that the ``Apology'' must be interpreted in agreement with the ``Phaedo,'' whose
line of thought he regards as Socratic (p.~193 f). I cannot find agreement between the ``Apology''
and the ``Phaedo.''} But a dogmatic decision and a great weight on such a \opage{151}decision
would not agree with the great significance he ascribes to clear knowledge, nor with the central
position the ethical viewpoint holds for him in relation to the religious. The double possibility
he discusses is not the expression of doubt and unrest in the mind, but of an inner sureness as
regards the essential task of life. Especially humorous is the statement that if there were
another life, it would be glorious for him to continue his activity there and apply his method to
the Homeric heroes, to see whether it should be with them as with so many of his dear countrymen,
that they thought themselves wise without being so! ---

It seems, then, to be fully justified when Socrates has been conceived in several places in what
precedes as a representative of the great humor. His irony stood in its service.

45.~The one point on which Socrates had concentrated himself radiated out, for his great
disciple, into a great world of thought. Plato found in the thoughts that make things intelligible
to us the true reality, and over against them the phenomena of experience and the circumstances of
life stood as fleeting images. He demanded of the thinker and the statesman a strict scientific
education, which was to lead, through mathematics as an intermediate link, to insight into the
ultimate grounds of knowledge and thereby, in his view, also of existence.\footnote{On
\textit{Plato's} significance in the history of thought see \textit{Den menneskelige Tanke.}
p.~117--121; 123; 140.} And when the high stage has been reached, the mind will be purified of all
lowness and pettiness. ``Do you think,'' he says (Republic VI p.~486), ``that he who has a high mind
and has gained a survey of all time and all being can regard human life as something great?'' ---
Plato does demand that the thinker, after he has reached the summit of knowledge, shall descend
among men and perform an office in the service of the state. But it comes as an outer necessity. It
is a matter of preventing the bad from coming to rule. Therefore the good, those who are least
willing to busy themselves with affairs of state, must take the helm. There is thus no direct
relation between ideal existence and life in the practical world. And although Plato,
\opage{152}rich human being and great poet as he was, had a lively apprehension of the contrasts of
life, and life stood for him as comedy and tragedy at once (see § 40), he nevertheless forgets in
the world of ideas both the tragic and the comic of existence. Change and contrasts rule only in the
sensuous world; but in the world of ideas unchangeableness and unity rule. Neither tragic lament nor
comic laughter is needed in the world of thought. And excessive laughter or lament is evil, because
the mind is thereby set in so strong a movement in a single direction that its firmness and unity
are weakened. Plato does not apply to comedy and tragedy, either of life or of art, what he so
insistently maintains of music, that a mood can be awakened that can persist even if no reasons can
be given for the thoughts or the deeds it gives rise to (``Republic'' III p.~401). There is here a
definite contrast between his hint at the end of the ``Symposium'' that it must belong to the same
man to write comedy and tragedy, and his standpoint in the ``Republic.'' In this latter work he
expressly withdraws the possibility of uniting comedy and tragedy. (``Republic'' III p.~295
B).\footnote{\textit{Hans Ræder (Platons philosophische Entwickelung.} Leipzig 1905 p.~208) finds
no contradiction between the two dialogues on this point, since the unity of comedy and tragedy in
the ``Symposium'' is to be set up as an ideal, while in the ``Republic'' this unity is conceived as
in fact impossible. But it is in reality two ideals that are set against each other: against the
poetic ideal of the ``Symposium'' the ``Republic'' sets the ethical ideal of a personality built on
self-limitation and inner harmony, which cannot give room for the great contrast between the tragic
and the comic.} % the Danish prints „philosophiscke“

A fusion or organization like the one lying in the concept of total feeling Plato regarded as
impossible when the greatest contrasts of life came to obtain their full, passionate expression.
Intellectual and ethical interest was in the end so strong in him that his rich human experience
could not deposit its full effect in his mind, at any rate not an effect that his thought could
acknowledge and express. For him, as for ancient thought in general, the unchangeable and the
changeable stood, once and for all, as the two greatest, irreconcilable contrasts,
\opage{153}between which a choice had in the end to be made. And for Plato there was no doubt
where true value lay.

His total feeling became intellectual enthusiasm, the joy of knowledge, awakened by the work of
thought in reaching the eternal ideas and thereby being raised above the stream of change.

Nevertheless Plato acquires great significance for the development of the total feeling we have
been occupied with, through his great thought of existence as a great order of things in relation
to which the events of human life stood as vanishing. An intellectual background was thereby
given that made possible a humor of another timbre than the Socratic. But Plato was so taken up
with working out this background, and so enthusiastic about it, that there was no time or strength
left to let the psychical connection between this background and the experiences in the world of
experience come fully into its own. There was lacking, moreover, the positive relation between
unchangeableness and changeableness that only modern times were to bring.

When nearly all energy flows to a single side of the mental life, there is naturally all the less
energy left over for the other sides. From this it is understood that those who are most taken
up with working out and working through the intellectual background of life do not themselves come
to stand as representatives of the great humor. We stand here at a new limit for humor as a view
of life or mood of life, different from the limits that the tragic on the one side and action on
the other set. The limit that intellectual work can set does not come from outside. The great
humor need not presuppose a dogmatic background, concluded once and for all, that would exclude
continued intellectual work. On the contrary, the great humor will, as we have seen (§§ 34--35),
draw in the knowledge actually won and take it ``humorously,'' because there will always be the
possibility of new riddles and new solutions. Behind every background there appears the
possibility of a still more comprehensive background. The constant re-arising of problems in new
forms stands as a comic, perhaps comitragic, phenomenon, and a smile can be called forth at the
busy work of knowledge. But this \opage{154}mood will not so easily arise in those who stand in
the midst of the purgatory of problems and feel their smart. ---

That Socrates is the only humorist in the grand style among the philosophers rests on the fact
that in him intellectual work became one with practical-pedagogical work among people. Through
irony as a method he sought to lead single persons to bethink themselves of the great background a
person can find in his own interior, whether this background can be shaped in clear thoughts or
not. Jest, irony, was here a way to earnest.

The other great humorist in the history of the human spirit was led by his experiences to produce
a world of images within which sorrow and joy, dark and light could stand as elements in a total
experience, a world of images that lets both the small and the great come into its own, without
the great crushing the small and without the small dragging the great down into the low.

Both the philosophical and the poetic humorist stood, each in his way, in a living relation to
life, however rich and exalted their inner world was.

46.~In accordance with the historical view indicated earlier (ch.~VIII), according to which we can
expect the basis and background for the great humor only in modern times, we now go straight from
Plato to the Renaissance. In the Middle Ages it will rather be in popular belief and popular custom
that one will find humor, thus in the representation of the stupid devil. An expression of a full
reckoning with the experiences of life it could not become under the presuppositions of the Middle
Ages.

Bruno and Spinoza laid the foundation, in broad strokes, of the view of the world peculiar to
modern times, gave the first pictures of the whole on the basis of what the new astronomy and the
new physics had brought to light.

I have already (§ 9) discussed Bruno's description and understanding of total feelings. His
philosophical and his poetic gifts worked together on this point, as on so many other points in his
thinking. But the total feeling he knew was not humor, but courage of life and magnanimity, the
feeling that is in general \opage{155}characteristic of the Renaissance (see § 42). It can become
a basis for humor (see § 20), but need not become so.

What Bruno contributed to the intellectual background of modern life of feeling was two things: the
widening of the physical picture of the world into the unsurveyable, and the assertion that if
divine forces are to be assumed in the world, they must be sought in the interior of the elements
and of souls, not in some place in space. Bruno's whole strength of spirit was taken up by the work
on these great thoughts. And the courage with which he held out in the tension and toil this work
caused lay at the ground of his description of the heroic mind or magnanimity. But the constant
yearning of thought laid claim to all the powers of the mind, and intellectual feeling entered into
no connection with an understanding sympathy toward the finite and limited in human life and in
nature.

For Spinoza it was not the infinity of existence but its lawfulness that stood in the front rank.
He urges as the most important means to spiritual freedom that our feelings be tied to valid
thoughts in which we can wholly find rest. (Ethica V, 4 Schol). The surge of the life of feeling is
to be stilled by our feeling ourselves one with the world substance, with what persists under all
changes, as we see the connection of our mind with the whole of nature (De emendatione
intellectus). Then the small and the changing are lost from sight, and the joy of knowledge becomes
the ruling mood. Spinoza certainly did not lack interest in the single and particular, especially
in the domain of the mental life. He wrote, after all, not only the first and second books of the
``Ethics,'' which treat of the world substance and of human knowledge, but also its third and fourth
books, which investigate human feelings and passions. But the two lines of thought --- the one that
seeks the fundamental relations of existence and of thought, and the one that explores the details
of experience --- do not come to unite in one and the same series of experiences, so that they both
come to determine the total feeling. The high point of life for Spinoza was the beholding
(intuition) that with one glance surveyed and saw through the great connection into which everything
fits. In his admiration for the infinite power and fullness of nature, which is able
\opage{156}to produce not only what in our eyes is perfect but also what in our eyes is imperfect
(Ethica I, end) --- a remark he set in sharp contrast to people's superstitious fear or narrow-minded
offense --- in this admiration he was raised above the comedy and tragedy of life, and these
contrasts did not come to determine in a positive way his disposition toward life.

It would be wrong to conclude from this that Spinoza in his personal life was untouched by the
feelings whose interplay conditions the great humor. Both fellow feeling and the feeling of laughter
played a part in him.

The misfortunes of others went nearer his heart, his first biographers report (Lucas, Colerus), than
his own, and it was to him an ugly comfort that one could find for one's own sufferings in others
suffering as we do. Only the purely passive pity that leads to blind and thoughtless intervention
did he reject. According to Spinoza's ``revaluation of all values,'' only those feelings are good in
and of themselves that are the expression of an active bearing and of the use of one's own strength
and reflection. Only joy, not sorrow, can therefore be good in itself. Yet in people who are not
guided by clear reflection, sorrow and sorrowful moods too (pity, fear and remorse, for example) can
be of value.

But he did not see the possibility that a feeling that in and of itself has no unconditional value
can enter as an element into a total feeling, which thereby comes to stand as the fruit of
many-sided experience of life and makes the mind well armed to meet what life may still bring.

A healthy and fresh laugh Spinoza counted among the good things in life, with a strong stress on
his contrast to ascetic and pietistic thinking (Eth.\ IV, 45 Schol.). He made a distinction only
between involuntary laughter (risus) and the laughter of scorn (irrisio). As soon as laughter gets a
definite object, it becomes for Spinoza an expression of scorn and stands in contrast to the
understanding that is for him the highest. It is this view that lies at the ground of his famous
statement that he endeavors neither to laugh at, mourn nor be indignant at human actions, but to
understand them (Tractatus politicus. I, 4). He applied this \opage{157}principle not only in a
treatise on politics but also toward the events of his time. The confusion of the time, he writes
in a letter (Ep.~30) during a war between Holland and England, awakens neither his laughter nor his
tears, but leads him to investigate and observe human nature; man is, after all, like everything
else, a part of nature, whose connection with the other parts and with the whole of nature we do not
know. The joy of knowledge gathered all his interest; it was his total feeling. But the great humor
rests precisely on a presupposition that Spinoza did not hold, namely that understanding and the
feeling of laughter do not exclude each other. Laughter can become the expression not of distance,
scorn or contempt, but of a feeling of unity. Understanding itself can bring laughter or a smile
with it. ---

For the men of the Renaissance it was first and foremost a matter of self-assertion. It was the
contrast to the Middle Ages' faith in authority and asceticism that asserted itself in them. They
were wholly taken up with the work of building up a new world of thought as a frame for life.
Through this work Bruno came to ``the heroic affect'' and Spinoza to ``the intellectual love.''
These were their total feelings, whereas the peculiar synthesis that humor presupposes did not come
about.

47.~While the Renaissance laid weight on self-assertion and the forms of feeling and will that could
develop from it, the English psychologists and moral philosophers of the eighteenth century --- the
series Shaftesbury, Hutcheson, Hume and Adam Smith --- were led to stress the social and sympathetic
feelings. This became of significance for the total feeling we have investigated here, since it can
arise from a sympathy that brings about a psychological connection between the personality and its
object and makes it possible for the object, however paltry, to be seen in its worth and
significance. It does not follow from this, however, that the philosophers named were themselves
humorists.

Shaftesbury was an enthusiastic worshipper of beauty. His concern was to combat fanaticism and
superstition, which would set narrow limits to the free unfolding of life. In this fight he regards
the comic as a good and legitimate ally. This is \opage{158}the fundamental thought of his little
work ``Essay on the Freedom of Wit and Humour'' (1709). (The word ``Humour'' is taken here in the
older, Jonsonian sense of the word). ``I for my part,'' he says, ``do not fear sceptical wit. No
cause in which I am in the least degree interested can be harmed by a tasteful wit.'' In this line
of thought there lies a germ of the great humor, for which it is once and for all so that even the
most valuable has a side from which it is subject to limitation and finitude, so that one can smile,
if not laugh, at it. But for Shaftesbury it was here only a matter of a means, a device, that is
not inwardly connected with a whole mood of life.

Hutcheson's main interest is ethical, and Smith's ethical and economic. Their psychological interest
was derived and subordinate. In Hume, on the other hand, one might expect to find humor. He has a
clear eye for the contradictions in the world and for the disproportion between what human feeling
demands and what reality shows. But the experiences he had in this respect seem not to have led to
any total feeling. On the contrary, he gives up striving for a unified expression of the experiences
of life. Such a striving rests, in his opinion, on an overvaluing of the significance of life. On
the other hand there are --- so Hume concludes his reflections (in the essay ``The Sceptic'') ---
natures for whom thinking about life is the most entertaining way of employing life!

48.~English literature in the eighteenth century moved, especially as regards the novel, in a
direction that offers an interesting parallel to the philosophical movement. Weight was laid on the
life of feeling, on sympathy, on the good heart, and the outward events of life were followed with a
mild eye and a friendly smile. We have to do here with the series of writers that Thackeray (The
English Humourists of the Eighteenth Century. 1853) put together under the name of the English
humorists. Thackeray's definition of humor is that it is a combination of wit and sympathy (wit and
love).\footnote{Thackeray hopes that he himself too can be counted among the humorists by this
definition, although in the ``Times'' he has been called a dull misanthrope who sees no good
anywhere. --- Thackeray himself, Dickens before him and George Eliot after him form a series of
English humorists whose characterization in relation to the writers of the same direction in the
18th century would offer great interest. For my purpose the characterization of Carlyle (see § 51)
will be more to the point.} In this special \opage{159}meaning it does not seem to fit all the men
Thackeray discusses: Swift, Addison, Steele, Pope, Fielding. Of Swift I have earlier (§ 15) tried
to show that he is a satirist, not a humorist. The one of the writers named whom the definition
mentioned best fits is Sterne, although Thackeray would rather call him sentimental. Goethe, who
otherwise did not care for ``humorists,'' made an exception for Sterne (see above § 25). In his
``Sentimental Journey'' (``a journey to learn to know the human heart'') Sterne says that his journey
is owed to ``the heart's need of nature, and of the feelings that spring from nature and make us love
one another and the world more than before.'' With this the basis is given. It leads to depicting
human characters and fates, now in earnest, now in jest, now with tears, now with laughter, with a
sense for the small things and for life in the small circles. What Shakespeare in his giant works
had been able to present, the great richness and strong contrasts of life, on the basis of deep
experience of life and faithful love of real life, was now carried by lesser minds out into wide
circles, down into the world of small things, which the master had given a place in his world only
episodically. There was often a lack here of the artistic ability to let the details take shape
sharply and individually. The naive, overflowing feeling often played with its subjects in a
capricious way, and instead of being the deep but hidden basis, the good heart liked to push forward
and speak along directly, so that heartiness became sentimentality. Add to this the smooth optimism
of the time, which made the life of feeling superficial. One passed lightly over the disharmonies.
The sense for the infinite and mysterious, which had been so strongly developed in Shakespeare and
Milton, had vanished in the poets of this time.\footnote{\textit{L.~Stephen: English Thought in the
18. Century.} II. p.~371.}

Nevertheless it was of significance that the spiritual and literary current arose that posterity has
called humor, although it was \opage{160}not the great humor here. It was the sign of a psychology
more open to the manifoldness of life. And in the course of the century the sense for Shakespeare's
greatness awoke again, especially through a great actor's rendering of his characters,\footnote{On
the significance of Garrick's acting for a Shakespearean renaissance see \textit{L.~Stephens:
English Literature and Society in the 18. Century.} p.~83 and elsewhere.} and thereby the sense for
great humor could be awakened again.

49.~The Renaissance had made a great intellectual background possible, and poetry had, now in
the grand style, now in small dimensions, awakened the sense for the manifoldness and peculiarity
of experiences. At the end of the eighteenth century there came in addition an attempt to settle
the relation between the work of thought and the other sides of the mental life. From the great
picture of the world and the great view of the world that the natural science and philosophy of
the Renaissance had made possible, Immanuel Kant descended to the power in the human spirit to
which it is owed that it can form images and intuitions at all, shape thoughts and set up ideals.
He uncovered what he called ``the hidden art in man's interior,'' the involuntary gathering
activity that expresses itself in all sensing and thinking. And when it was asked whether
validity could be ascribed to the results of this activity, his answer was that this validity
rested on whether they were necessary for being able to distinguish between experience and
imagination, between reality and dream. Thereby the ground was taken away from all speculative
systems that sought beyond every possible experience. Theology and spiritualistic psychology
became impossible as sciences. Not without inner struggle had Kant reached this result. For he
had, as he himself said, the fate of being in love with metaphysics --- and now he saw that it
was impossible. Into this divided mood his first investigations had already brought him. It was
through humor that he got beyond the dividedness. In ``Träume eines Geistersehers, erläutert durch
Träume der Metaphysik'' (1766) a humorous mood reigns. Before the high and strict concept of truth
he had now caught sight of, speculative systems stood for him in a comic light, although it had
been \opage{161}his own inmost wish to be able to produce such a system. Later too, after the
working out of his famous main work, he can look at his work with humor. In one of his notes he
jests about the pedantic character that the ``Kritik der reinen Vernunft'' undeniably in part
presents; he consoles himself that pedantry must drive out pedantry. His Socratic character,
which made him prefer thorough and independent knowledge to speculation as well as to the mere
heaping up of information, and set philosophizing above having a philosophy, gave him a firm
background for regarding fruitless but often all the more confident attempts of thought.

In Kant's life of feeling the sublime played the greatest part --- the infinity of the universe
outwardly and the inner infinity appearing in conscience were, according to a famous statement of
his, two things that filled his mind with ever new reverence. But as the man of firm, rational
rules he fixed his life and his thought more and more in definite forms, which it was painful and
at last impossible for him to think of as otherwise. The humor that could stir in him in his
younger years in the intellectual domain did not come to give its stamp to his whole personality
and mood of life.

50.~It is within German Romanticism that the concept of the great humor is first expressly
formed. % the print numbers this section "49." a second time
As everywhere, the thing here too goes before the concept. We have pointed to two great figures as
representatives of what the concept contains --- the one in antiquity, the other at the entrance
to modern times.

The English humorists of the eighteenth century do not, as already remarked, represent the great
humor. Their horizon was too limited, and their optimism did not let the contrasts of life have
their full influence. One moved on an even, bourgeois level, and since feeling had no contrast and
resistance to overcome, it became sentimentality, especially in imitators. According to Jean Paul
(Vorschule § 32) Sterne got a whole host of imitators in Germany who in reality were nothing but
``Ausplauderer lustiger Selbstbehaglichkait.'' A great idea was lacking as background. Such an
idea German Romanticism commanded.

\opage{162}Not as if the German Romantics had themselves created this background. As I have
already remarked in ``Den nyere Filosofis Historie'' (in the introduction to the eighth book), they
built essentially on the great German renaissance in the last part of the eighteenth century
(Lessing and Kant, Herder and Hamann, Schiller and Goethe), on the widening of the spiritual
horizon it had brought, and on the representation of the originality and manifoldness of
spiritual life it had awakened. Of special significance besides was that one came to know
different directions of spirit and thought, the life of faith and mood of many ages, and acquired
a sense for and a need to enter into manifold ways of seeing and standpoints, and to look with
irony on each single form of spirit when it wanted to be the only valid one. In Fr.~Schlegel we
have already found this trait (see § 16). With a mixture of irony and admiration one looked back
on the faith of earlier times, on their art and their thinking. One had a deep conviction that
what thus constantly sought new forms and burst the old ones was stirring now too in all whose
senses and minds were opened. And it was not something that stirred only in human minds. What was
the highest for thought and need must also be the innermost in existence, although no single part
of our experience can express it completely. Only a higher beholding or a speculative thought
opened access to it. Each single experience, each representation taken from a special field of
observation, could acquire only symbolic significance.

The philosophy of Romanticism is an impatient idealism. One strove unrestrainedly, through a kind
of magic, to penetrate to ``what holds the world together in its inmost,'' and one had no sense
for work in detail, from special starting points and by special methods. Such things stood, like
the tasks and interests of plain bourgeois life, as a jest in relation to the ideal world that
revealed itself in so many forms. --- Romanticism worked liberatingly by pointing beyond the
narrow forms in which life, art, religion and thought so easily stiffen. And it could thereby make
great humor possible. But romantic humor suffered from a great imperfection. For the small
circumstances of life and the small steps of thought may certainly stand as \opage{163}vanishing
and insignificant in relation to the ideal and infinite we can divine; but it is nevertheless
through faithfulness in small things that great things are carried through. Romanticism conceived
every limited striving, every one-sided direction, merely as something negative, as something that
claimed to be infinite without being so, something that set a limit to the comprehensive view.
Therefore the Medusa head of irony was turned against them. There was here the possibility of an
irony like the one wrongly found in Fr.~Schlegel (see § 16), and like the one Kierkegaard
described in his treatise on the concept of irony and later in his characterization of ``the
aesthetic stage.'' But there was also the possibility of a connection of irony and enthusiasm that
leads, and led, to the great humor.

Novalis already maintained that Fr.~Schlegel was really a humorist. But the theorist of romantic
humor was Solger, who also struck the strings of the great humor. In him the humor of the
eighteenth century, which is the combination of wit and sympathy, is united with Romanticism's idea
of the single as pointing beyond itself to a great whole. In his view of life and his conception
of art he unites irony and enthusiasm. Through irony the single is consumed in its isolation;
through enthusiasm it is seen as a member of the whole. In his ``Erwin'' it is the great longing
(§ 21) that forms the basis; his posthumous writings seem to show that in his last years he took a
more mystical direction. He is a Romantic by his longing for a form of spiritual life in which the
difference between art, religion and philosophy has vanished. Philosophy is to prepare the way for
such a high point by showing, through free knowledge, the eternal and infinite as contained in the
laws that hold for finite things. Art is to show, through images, the finite as shone through by
the infinite. And religion is an immediate living in the eternal and infinite. At the high point
none of these three forms of spiritual life holds; the difference between them falls away. The
highest state is enthusiasm; but it gives rise to irony toward reality and its phenomena.

It is a personal standpoint of life that Solger has depicted here. His abstract expressions have
caused Kierkegaard to depict his \opage{164}standpoint as purely formal and negative and to regard
him merely as an imperfect forerunner of Hegel. But while for Hegel all spiritual forms of life
were imperfect forms of philosophical speculations, Solger maintained that the highest form of life
must be one to which philosophy, art and religion each made its contribution.

The deficiency of romantic humor, a deficiency that comes out more in Solger's forerunner Jean Paul
than in Solger himself, consists in there coming no positive relation between the small and the
great. Predominant weight is laid on the impossibility of a perfect presentation of the highest.
The great humor reaches its full unfolding only when the small and the limited are seen as
necessary forms for the great and infinite and thereby acquire positive significance. Romantic
humor could lead, and did lead, to disregard for empirical science and for practical endeavors of
reform. In Carlyle we shall see this plainly. The realistic element in the great humor, which
springs from its faithfulness to reality, does not come into its own in Romanticism. Neither the
tragic (ch.~V) nor understanding (ch.~VI) nor action (ch.~VII) --- the three domains over against
which humor is to show itself at once in its positive significance and in its necessary limitation
--- are acknowledged in their independence. ---

The great humor did not reach its full development in German spiritual life. German humor keeps a
more or less romantic-speculative character. English humor is more earthbound, but on the firm
ground of reality Shakespeare's world of mighty figures in dramatic interaction grew up, a world in
which the most intimate and at the same time the most plastically shaped humor appears. --- It is
characteristic that the French language uses the English form (humour) to express humor in the
sense in which we take it here. The fusion or organization the great humor presupposes seems not to
have natural conditions on French soil. It is at any rate more characteristic of the French nature
to swing between very different moods that can unfold in a certain independence of one another.
This explains the surprises \opage{165}the French often offer the world. Thus a French philosopher
maintains that no psychological miracle has taken place with his people for a thorough observer,
but that one could already earlier have traced the excellent qualities it has shown in the present
war.\footnote{Nous avions pu discerner, sous une frivolité apparente, un esprit sérieux et capable
d'application, sous des dehors frondeurs, la résolution de se soumettre librement à une discipline
raisonnable, et sous la ``blague'' superficielle, le sentiment et le respect de l'ideal.
\textit{Terraillon:} Allocution prononcée à la Distribution des prix. Collège de Saint-Claude. 1915,
p.~9.} --- It is perhaps bold to conclude that psychical chemistry, a close connection and coherence
of different psychical tendencies, should come more naturally to Anglo-Saxons and Germans than to
the French. But there is at any rate here a task for closer investigation, when racial psychology
has once outgrown its infancy.

51.~Carlyle is the most significant personality who represents romantic humor. The basis for him
was the great earnest of life, to which his heavy mind opened his sense, and which the Calvinist
faith he had grown up in had impressed on his way of thinking. Later he lived himself into the
world of thought of German poets and philosophers. Especially (as he himself declared in his later
days) Goethe's ``Wilhelm Meister'' influenced him, and here again the idea of the three kinds of
reverence urged in the ``Wanderjahre'': reverence for the great, what is above us --- for the
equal, what is beside us --- for the small, what stands beneath us. Reverence forms the basis of
Carlyle's life of feeling and of his view of life. What he appropriated of German philosophy was
especially the great contrast between the inner core of life and the outer husk. Kant's contrast
between ``the thing in itself'' and ``the phenomenon'' he used, not, like Kant himself, to stake
out the domain of science and settle down there, but to express that scientific research leads to
nothing, since it does not lead to everything. All experience became for him a symbolism that was
related to the thing as clothes to the body. His philosophy was ``the philosophy of clothes,'' with
the main weight on clothes not making the man. (Sartor Resartus. 1836). This conditions his
\opage{166}relation to the teachings of science as well as to the dogmas of religion, and to all
practical endeavors. All these things he counts among the clothes, or among the patches put on the
clothes. He stands toward such endeavors with a mixture of mockery and pity. When Carlyle so often
says ``poor'' of a person, it is partly sympathy, partly contempt that is expressed by this word. In
a drastic way he depicts what in his eyes is the senselessness of working for truth and justice in
the world. And yet his heart was warm. And he loved healthy, spontaneous laughter, could laugh
heartily himself, and thought that the person who had once laughed from a full heart could not be
altogether bad (see § 13). But his laughter went strangely alongside the often bitter jest or broke
through as an effect of contrast to it. Carlyle is a humorist by his sense for the contrast between
kernel and shell and by his understanding that the shells change and must change. But his idealism
is, like that of Romanticism in general, of the impatient kind that demanded all or nothing, and so
by this apparently courageous demand in reality shut itself up in the world of fantasy and mood.
There was lacking a positive bond between the main thought and its special theoretical and practical
applications, which he therefore would not acknowledge as applications either. His unruly
temperament made him scorn what by virtue of his own principle he must have acknowledged as a
necessary, though imperfect, expression of reality. A secret doubt expressed itself in him here.

It has rightly been stressed that it is the complicated circumstances of modern life that confuse and
depress him. His thought always moved about the simple, broad facts of life. When from his solitude
in the Scottish Highlands he was transported to Edinburgh and London, civilized life met him with all
its manifoldness of research and striving. Like Rousseau\footnote{It was, by Carlyle's own account,
Rousseau's ``Confessions'' that first awakened him to deeper reflection. (According to what he told
Emerson, who in turn told George Eliot. See \textit{George Eliot's Life,} as related in her letters
and journals. New edition. p.~104). This was not known to me when I wrote my characterization of
Carlyle in ``Den nyere Filosofis Historie'' (ninth book) and in ``Mindre Arbejder'' III p.~125.} he
stood uncomprehending and rejecting before the culture of his time. But while \opage{167}Rousseau did
learn to understand that there is a natural culture, Carlyle never got beyond the sharp contrast that
he expressed in the most violent way in single moments, under the influence of his morbid and vehement
mind. And yet he had in his humor, and in the thoughts that formed its background, the conditions for
getting beyond this contrast.

Romantic humor did not reach full fusion or organization --- Carlyle is the greatest example of this.
The contrasts that could not enter into full inner connection then asserted themselves in successive
contrast. It is reported of Carlyle that when he had in his rough way torn down everything in the
world, he could suddenly stop and, through a powerfully breaking-out laugh, return to more level
regions. His own prophetic tirades must then have come to stand in a comic light for
himself.\footnote{Henry James the elder (the father of William James and of the novelist Henry James)
describes a delightful scene from Carlyle's tea table. An American was disputing with Carlyle about
politics. Carlyle became more and more violent in his replies, and Mrs.\ Carlyle therefore trod on the
guest's foot to get him to stop the dangerous topic. But the American cried out loud: ``Why don't you
tread on your husband's toes, Mrs.\ Carlyle? He is surely worse than I am!'' The whole company burst
into hearty laughter, and Carlyle led it. (\textit{Henry James: Some personal recollections of
Carlyle}. Literary Remains. Boston 1897. p.~454).}

52.~According to the psychological characterization of the great humor given in what precedes, it
has a decidedly realistic character and has therefore been called faithfulness to reality. When it
is, as it seems, rarely found in eminent philosophers, the reason must, as said above, be sought
especially in its being the intellectual interest about which everything gathers for them, and
which therefore comes to determine the ruling total feeling. Often the reason may lie in an
impatient idealism that leads to stressing the contrast to experience more than the connection
with it. James Sully (Essay on Laughter. p.~397. 401), who has also made the observation that
philosophers are rarely humorists, finds the reason in the philosopher's being removed by his
studies too \opage{168}far from the stage of real life. I would, however, lay the main weight on
the contrast between the life of thought and the other sides of the mental life. In every man of
science a greater or lesser contrast can appear here. And even if the philosopher has special
occasion to investigate the relation between thought and personality, it easily becomes this very
investigation that lays claim to all his interest and energy. The life of thought too has its
fate, brings its experiences, which can become determining for the final mood toward life. Even if
real life intervenes with a strong hand, it can nevertheless be the work of thought, and this
work's fortune or misfortune, that becomes the decisive thing. However important the background of
thought and science for life may be, its significance for a total feeling like humor is only
indirect. And, as we have seen (§ 21), a thorough experience of life won by a purely practical path
can have the same significance.

A distinction has been made in what precedes between background and basis. The basis, which lies
especially in the life of feeling and will, is more important than the intellectual background. It
proves, accordingly, that from a scholar's general way of thinking one cannot infer the character
of his total feeling toward life. To draw such an inference requires knowledge of the whole concrete
basis of the personality. Unfortunately this is so hard to know, and one must as a rule build here
on a hypothesis proceeding from the total feeling that has actually formed, and to which the
background has also made its contribution.

Perhaps great and strong total feelings can form more readily in artists than in thinkers. There
is a closer relation between image and feeling than between thought and feeling. The feeling at
what is experienced finds expression involuntarily in images, which present or symbolize vividly
what is experienced in its whole peculiarity, while for thought the single experience easily comes
to stand as one example among several. The artist thereby stands in a closer relation to life in
its concrete, individual forms. That is presumably why Schiller writes in a letter to Goethe that
the poet is the only true human being, while the best philosopher is only a caricature in
\opage{169}comparison with him.\footnote{Cf.\ my memorial address on Schiller in \textit{Mindre
Arbejder.} Third series. p.~179 f.} Yet the artist too can make life a mere object of
contemplation (cf.\ § 37), so that the world of fantasy becomes for him the only world. Schiller
himself advised a friend to stay in ``the bright and quiet world of ideas'' and for the present let
``poor, unworthy and immature mankind'' look after itself.

Where energy is used predominantly for life in thought or fantasy, energy will be lacking for the
fusion or organization of the feelings the course of life can awaken. In the regions beneath the
world of thought and images manifold feelings may then often tumble about chaotically.
Hypochondria or levity, phlegm or irascibility, pettiness or magnanimity, bitterness or
sentimentality, sensuality or ideality --- these and many more tendencies can form a sublunary world
that one who knows only the works of thought or fantasy stemming from the ethereal world above does
not suspect. It is not always the case that the life of thought or fantasy, besides being the domain
where energy is first and foremost put in, can also form the background for the development of a
total feeling in the grand style. Only when this happens are we at the high points of human life.

It is not, for the rest, only in thinkers and poets that this contrast can be found between a world
in which they live their most energetic life and a sublunary world of a wholly different structure.
In practical natures too, men of action, there can be a corresponding contrast, namely between the
part of their life that is absorbed in work, public activity or business, and the sides of their
mental life that are not taken up by this activity. ---

One might (cf.\ § 38) also look at the matter from the opposite side and find it doubtful whether
the life of feeling could be joined into complete totality, since thereby the one-sidedness that
earnest work and energetic action demand would become impossible. % the Danish prints „derved derved“

Gustav Fröding, in a little essay from the year 1890, maintained in a fine way the significance and
right of humor over against the one-sided part that, in his opinion, \opage{170}the striving for
truth and moral yearning played in the literature of the time. He thought there was more humor in
the old days than now, although so much is preached about love nowadays. ``Man är så kärleksful i
ord och gärning, att man alldeles glömmer att vara det i själ och sinne, hvilket ändå är det, som
verker försonande människor mellan.'' ``Humor är hundarne, som slicka såren på den lidande Lazarus,
och vänligare än den kristliga kärleken aktar den icke för rof at skänka den rike mannen den
vattendroppe, efter hvilkan hans tunga trådde.'' --- But when Fröding in the year 1910 reprinted his
essay (in which he had mainly depicted humor on the basis of sympathy), he had lost faith in the
right of the great humor and thought he had overvalued it. In his opinion it ought not to rule
alone.

Nor have I in what precedes been blind to the misgivings it can awaken, as I have brought forward
the dangers of every concluding total feeling. But through the investigation of the relation of
humor to understanding, to the tragic and to action, I have tried to show that from humor's own
nature there follow certain limitations that need by no means come from outside.

It was a brilliant idea of Kierkegaard's to conceive the great humor as the high point of what can
be reached by a purely human path, in contrast to the religious standpoints that presuppose a
relation to something lying beyond the reach of human experience and knowledge. The mere fact that
there can form such a total feeling, which holds the opposed elements of life without excluding the
possibility that these elements, if necessary, can be released anew in their particularity --- the
mere fact is already of decisive significance. Therein appears a typical example of a human view of
life.

53.~Finally we shall consider the relation of humor to a realistic way of thinking, by which I
understand a way of thinking that finds the criterion of truth and reality in the firm, lawful
connection in everything that appears as an object of observation.

Humor is, as has been said several times in what precedes, a realistic feeling, and to that extent
it might seem to be wholly \opage{171}natural at a markedly realistic standpoint. But it appears
here too that the way of thinking is only a condition, not necessarily decisive for the arising of
a total feeling. The background is not enough; it depends especially on the basis. While the
idealist's danger is that he easily becomes absorbed in pure ideas without coming into close
relation to the concrete realities of life, the realist's danger will be that he is weighed down by
the ``realities'' of life and does not come to make his understanding of them a mere background for
his life of feeling. The background then comes to check the unfolding of the basis. If for the
idealist life fades into an ethereal world, for the realist it is easily stifled in the dense fogs of
``reality.'' Such, once and for all, is man's lot, according to Goethe's famous words (Grenzen der
Menschheit): if we rise high, we can find no foothold, and if we stand firm on the earth, we do not
reach high!

If realism is to become the background for a free and fresh development of the life of feeling, the
realist must look critically at his own realism. The same condition holds for the idealist. Kant, who
in his ``Träume eines Geistersehers'' had turned critically and humorously against a dogmatic
idealism, prepared the way in his ``Kritik der reinen Vernunft'' for a critical realism.

When the realist asks himself in earnest what reality (realness) really is, he will notice that when
we are in doubt whether something is real or not, we go, in daily life as well as in science, the
way of gathering as many observations as possible about the subject concerned and seeking to find as
close a connection as possible among these observations. One experiments not only in science but
also in daily life. The criterion of reality summons to ever new work --- both to find or produce new
observations and to find a closer connection among the observations found. The concept of reality
stands as an ideal we can approach, but which always makes new demands. Armed with this ideal the
great humor pushes all dogmatic realism aside. No given reality satisfies the concept of reality.
(Cf.\ § 34).

Within such a view of existence as this at once experimental and ideal concept of reality makes
possible, \opage{172}one can unite interest in and understanding of the single in its peculiarity
and its worth with a view toward a greater whole. However vanishing the single may seem, its worth
nevertheless never becomes zero. And the greater whole itself, of which it is a member, proves in
turn to be vanishing over against more comprehensive wholes. It may perhaps be foreseen that its time
too will one day be over, as Shakespeare's Prospero divines it, and as modern scientists have often
assumed as regards our world system. One need not be an optimist at all to be a humorist; on the
contrary, optimism (cf.\ § 48) will exclude humor by preventing a sufficiently thorough penetration
into all sides of life. The humorist will not belong among those who rejoice over how gloriously far
we have brought things; such joy will rather awaken his irony. He knows that the struggle between the
human and the animal is long; it goes, like a modern war, from trench to trench. But as long as there
is still fighting, hope and trust are both possible and necessary. For Charles Darwin this slowness
of development was not the greatest stumbling block. What stood for him as a difficulty was that when
development had advanced through enormous periods of time, everything should then dissolve again into
a chaos. He writes to Hooker (9 Feb.\ 1865): ``I quite agree how humiliating man's slow progress is; but
everyone has his own hobbyhorse as regards fear (his own pet horror), and the slow progress, indeed even
personal annihilation, sinks in my eyes into an insignificance compared with the representation, or
rather the certainty, that the sun will one day grow cold and we all freeze to death. To think of a
progress through millions of years, with every continent swarming with good and enlightened men, and
everything ending thus, without there probably coming any new beginning before our planetary system
has again become red-hot gas! Sic transit gloria mundi --- with a nemesis (with a vengeance).'' --- Darwin did not, however, give up his faith that development in the
world had the upper hand over dissolution. He found, like Kant, the fixed point in the feeling of
duty. ``Man can do his duty.'' From the possibility of realizing something valuable there
\opage{173}shone at last the light in which it became possible for him to live and
strive.\footnote{See my paper on \textit{Darwin og Filosofien} in the Festschrift ``Darwin and Modern
Science'' published in Cambridge, in Danish in ``Mindre Arbejder,'' Third series. (p.~225--228).}
% Darwin's letter is translated from Høffding's Danish rendering (his English parentheses kept)

Should it be that the spring of values were one day exhausted, because everything in the end came into
absolute equilibrium, then, as already said (§ 30), the humorist will not be able to take this
humorously. Then he admits that the tragic gets the last word. Indeed, he would wish that it could
get the last word. For it will unfortunately go so that the extinction of psychical life will be
preceded by a period of weakening and dullness in which there will not even be talk of tragic
suffering. The tragic presupposes strength. But the humorist works as long as he can to find new
values, and he takes a critical attitude toward all theories, speculations and perspectives that
dogmatize without sufficient ground about ``the last things.'' Physics at bottom knows no more about
them than theology. Just as little as there is reason for, perhaps not even possibility of, thinking
of a conclusion of time and space, just as little can proof ever be given that all possibilities are
exhausted, ``all sounds closed.'' By virtue of the unsurveyability of the intellectual background
existence can always come to stand as new.

54.~From realism one does not come with necessity to humor. But critical realism can perhaps furnish
a background for humor better than most idealistic directions.

A necessary component of such a critical realism is the acknowledgment of the fact that every
personality, to a greater or lesser degree and with greater or lesser consciousness, forms and must
form its own view of life, which cannot without more ado be derived from the positive or negative
dogmas to which individuals profess allegiance. Precisely from a realistic standpoint it must be
demanded that everyone build his view of life on his own experiences and with the help of the
knowledge he has won by the work of his own thought. Only so does it build on the ground of reality,
on what is really reality for \opage{174}the single person. A view of life is worth something only
when it is really the expression of how life stirs at the definite place in existence the single
person occupies. The demand for personal truth is realistic, although it was first put forward by an
idealist.

Experience shows more and more how meaningless dogmatic declarations can be in relation to the
realities of the inner life. From what is confessed, sung and proclaimed one can draw no sure
conclusion if one does not gain access by another path to the inner life that lies behind. Behind
theology or anti-theology the realist, and the humorist with him, seeks psychology. He knows that in
the inner world of the mind great things can go on in all quietness, and the noisy disputes about the
index and the headings of the book of life only draw a smile from him. The principle of personality
works not only dissolvingly against every kind of dogmatism and every kind of superficial community;
it also expresses the inexhaustibility of the inner life.\footnote{See
\textit{Religionsfilosofi}\textsuperscript{2}. p.~284--286.} There opens here an inner world that
belongs to the intellectual background of the life of feeling. In each single personality this inner
world can find expression only in a more or less one-sided way. It goes with the art of life as with
the visual arts, which, as Julius Lange so energetically maintained, must be one-sided because ``no
single work of art can communicate the many-sided conception of the human
figure.''\footnote{See \textit{Religionsfilosofi}\textsuperscript{2}. p.~256.}

By acknowledging this view the great humor acknowledges at the same time its own limitation. Even if
it should have ever so great a significance as a characteristic total feeling in which everything
human can come into its own, it will nevertheless never be able to proclaim itself as the only right
way of taking life. It takes a humorous attitude here toward itself. On the ground of reality it sees
human life unfold from the most different starting points and under the most different conditions of
conflict, and it understands that all development proceeds sporadically. Therefore it does not wonder
that total feelings come to be of different character. It is not strange that not all are humorists.
And even if they were, they would be so each in \opage{175}his own way. Harmony and community are
great goods, but must not be bought with the sacrifice of originality and genuineness in single
persons. The sporadic character of development makes the world of spirit at once great and small,
tragic and comic --- and this contains (cf.\ § 35) an occasion for the arising of the total feeling we
have been occupied with here.

% ===== INDEX (REGISTER, pp. 177--178) =====
\chapter*{Index}
\addcontentsline{toc}{chapter}{Index}
\markboth{Index}{Index}
\noindent{\small [In the order of the Danish index; the page numbers are those of the Danish
edition, given in the margin of this translation. The Danish headword follows in italics where it
differs. Italic page numbers are the print's main places.]}
\medskip

\opage{177}\entry \textbf{A}cedia 96.
\entry Earnest \textit{(Alvor)} 68 ff. 125.
\entry Amiel 131.
\entry Antigone 100. 129.
\entry Aristotle \textit{(Aristoteles)} 101 f. 129.
\par\smallskip
\entry \textbf{B}ackground (for humor) \textit{(Baggrund)} 86. 113 f. 145 ff. 166. 171.
\entry Benn 102.
\entry Decision \textit{(Beslutning)} 10. 19. cf.\ 123 ff. % the Danish prints „fi.“
\entry Bradley (A.~C.) 98 f. 117.
\entry Brandes (G.) 100.
\entry Bruno 39. 143. 154 f.
\entry Burton 41.
\par\smallskip
\entry \textbf{C}arlyle 34. 55. \textit{165--167}
\entry Cervantes 43.
\entry Comte 114.
\entry Cusanus 121.
\par\smallskip
\entry \textbf{D}ante 96--98.
\entry Darwin 172.
\entry Descartes 9.
\entry Differentiation 125.
\entry Dowden 43. 67.
\par\smallskip
\entry \textbf{O}nce-born \textit{(Enfødte)} 13. 140.
\entry Single feeling \textit{(Enkeltfølelse)} \textit{7--12}.
\entry Ethics (and total feeling) \textit{(Etik)} 36 f. 125 f.
\par\smallskip
\entry \textbf{F}alstaff \textit{(Fallstaff)} 116 f.
\entry Feilberg (Ludvig) 10. 12.
\entry Felix culpa 80. 83. 124.
\entry Philosophy and view of life \textit{(Filosofi og Livsanskuelse)} 45. \textit{145--175}.
\entry Freud 102. cf.\ 13.
\entry Fröding \textit{(Frøding)} 169 f.
\par\smallskip
\entry \textbf{G}alileo \textit{(Galilei)} 120.
\entry Genetic psychology \textit{(Genetisk Psykologi)} 16. 126.
\entry George Eliot 54. 82. 140.
\entry Goethe 55. 57. 92 f. 100. 127. 171.
\entry Goldschmidt 131.
\entry Basis (for humor) \textit{(Grundlag)} 86. 113 f. 166. 171.
\entry Gunnarson (G.) 104.
\entry Mrs.\ Gyllembourg 83.
\par\smallskip
\entry \textbf{S}corn \textit{(Haan)} 52. 59.
\entry Action \textit{(Handling)} 125 ff.
\entry Hegel 100 f. 137.
\entry Heroic passion \textit{(Heroisk Lidenskab)} 39.
\entry Hirn (Yrjo) 68. 73.
\entry Hobbes 48. 55.
\entry Holberg 42. 120.
\entry Home (Henry) 46.
\entry Hume 158.
\entry Humor. Meaning of the word 40--45. Kinds 52 ff. \textit{75--87}. The great humor 43 f. 47.
54. \textit{56--58}. Humor and irony 68 f. Humor and reality 70 ff. Humor as an art of living
92 ff. Humor and philosophy 145 ff.
\entry Hypochondria \textit{(Hypokondri)} 95 ff. 169.
\par\smallskip
\entry \textbf{I}phigenia \textit{(Ifigencia)} 120. % so printed; the text is on p. 129
\entry Irony \textit{(Ironi)} 61 ff. 147 ff.
\par\smallskip
\entry \textbf{J}ames (W.) 28.
\entry Jean Paul 47. 141. 161. 164.
\entry Self (real) \textit{(Jeg, realt)} 11--15. 23. 27. 35. 54.
\entry Jonson (Ben) 41.
\par\smallskip
\entry \textbf{K}aalund 57.
\entry Kamm (Adèle) 90.
\entry Kant 34. 76. \textit{169 f.}
\entry Kierkegaard 21. 34. 47. 64. 74. 89. 93. \textit{104--109}. \textit{132--134}. 147.
\opage{178}\entry Christianity (primitive Christianity) \textit{(Kristendom, Urkristendom)} 81 f. 105.
\entry Cult \textit{(Kultus)} 141.
\par\smallskip
\entry \textbf{L}ange (Julius) 27. 174.
\entry Laughter \textit{(Latter)} 48 ff.
\entry Lehmann (Alfr.) 22. 29.
\entry Leopardi 119.
\entry Lessing 89.
\entry Lichtenberg 55.
\entry Lotze 58. 65.
\entry Longing \textit{(Længsel)} 53. \textit{80--83}.
\par\smallskip
\entry \textbf{M}aier (Heinrich) 148--150.
\entry Medea \textit{(Medeia)} 129.
\entry Mephistopheles \textit{(Mefistofeles)} 117.
\entry Mozart 9.
\par\smallskip
\entry \textbf{N}ietzsche 59.
\entry Novalis 163.
\par\smallskip
\entry \textbf{E}vil (the origin of evil) \textit{(Ondt, det Ondes Udspring)} 80.
\entry Optimism \textit{(Optimisme)} 140.
\entry Organization (of feelings) \textit{(Organisation)} 18. 26. \textit{30--33}.
\par\smallskip
\entry \textbf{P}ersonality \textit{(Personlighed)} 13--102.
\entry Principle of personality \textit{(Personlighedsprincipet)} 64.
\entry Pessimism \textit{(Pessimisme)} 140.
\entry Plato \textit{(Platon)} 15. 26. 38. 55. 136 f. \textit{151--153}.
\entry Poincaré (Henri) 37.
\entry Prospero 99. 107. 162.
\entry Psychoanalysis \textit{(Psykoanalyse)} 13. 102.
\entry Psychology \textit{(Psykologi)} 34. 76.
\par\smallskip
\entry \textbf{R}eal self, see Self.
\entry Religion 108 ff. cf.\ 104. 150.
\entry Ribot 26. 49.
\entry Rousseau 10. 166.
\entry Ræder (Hans) 152.
\par\smallskip
\entry \textbf{F}usion (of feelings) \textit{(Sammensmeltning)} 18. 23. \textit{26--30}. % the Danish leaves the parenthesis open
\entry Truth \textit{(Sandhed)} 115. 117.
\entry Satire 60 f.
\entry Schlegel (Fr.) 65 ff. 69. 163.
\entry Schlegel (A.~W.) 67.
\entry Shaftesbury 157 f.
\entry Shakespeare 40. 42 f. 67 f. 85. \textit{97--100}. 116 f. 129. 141. 147. 159 f.
\entry Shand 26. 31. 36.
\entry Sibbern 25. 29. 38. 73. 87. 121. 145. 147.
\entry Sidney Lee 100. 116.
\entry Disposition and emotion \textit{(Sindelag og Sindsbevægelse)} 128 ff.
\entry Loftiness of mind (magnanimity) \textit{(Sindshøjhed, Højsind)} 03. \textit{77--79}. % so printed: „03.“
\entry Pain \textit{(Smerte)} 91.
\entry Sneedorff 120.
\entry Socrates \textit{(Sokrates)} 21. 61 f. 138. \textit{146--151}.
\entry Solger 47. 73--93. 163.
\entry Spinoza 19. 97. 113. \textit{155--157}.
\entry Sporadic development \textit{(Sporadisk Udvikling)} 75.
\entry Sterne 93. 189. % so printed: „189“ (the book has 178 pages)
\entry Stoicism \textit{(Stoicisme)} 79.
\entry Stuckenberg 83.
\entry Suddard (Mary) 40.
\entry Sully (James) 49. 51. 54. 167.
\entry Swift 60. 159.
\entry Sympathy \textit{(Sympati)} 58. 84 f.
\par\smallskip
\entry \textbf{T}emperament 139.
\entry Thackeray 54. 158 f.
\entry Thorndike 16. 126.
\entry Total feeling \textit{(Totalfølelse)} \textit{18--26}. 76. 76. 77. 128. 145--148.
\entry Total state \textit{(Totaltilstand)} 1--25.
\entry The tragic \textit{(Tragik)} 89 ff. 173.
\entry Faithfulness \textit{(Troskab)} 43. 115 f.
\entry Twice-born \textit{(Tvefødte)} 14. 140.
\par\smallskip
\entry \textbf{V}elázquez \textit{(Velasquez)} 27.
\entry Wistfulness \textit{(Vemod)} 14. 30. 53. \textit{80--83}.
\par\smallskip
\entry \textbf{W}ilhelm Meister 89.
\par\smallskip
\entry \textbf{C}riterion of reality \textit{(Virkelighedskriterium)} 114 f. 160. 171.
\entry Value \textit{(Værdi)} \textit{68--73}. 118. Persistence of value \textit{(Værdiens Bestaaen)} 99. 106.
109 f. Value and reality 80 f. 118 f.

\end{document}
